%% ============================================================================
%%  CAPITOLO 6 — Applicazioni delle derivate. Studio di funzioni
%%  Fonte principale per argomenti e dimostrazioni: Marcellini–Sbordone, cap. 6
%%  (§46 massimi e minimi relativi, teorema di Fermat; §47 teoremi di Rolle e di
%%  Lagrange; §48 funzioni crescenti e decrescenti; §49 funzioni convesse e
%%  concave; §50 teorema di L'Hôpital; §51 studio del grafico di una funzione;
%%  §52 formula di Taylor; appendice §53 teorema di Cauchy e teorema di
%%  L'Hôpital nel caso generale).
%% ============================================================================
\aperturacapitolo{Applicazioni delle derivate. Studio di funzioni}{cap:applicazioni-derivate}

\epigraph{Nihil omnino in mundo contingit, in quo non maximi minimive ratio quaepiam
eluceat.\\[2pt]
\upshape(Nulla accade nel mondo in cui non risplenda una qualche regola di massimo o di
minimo.)}{Leonhard Euler, \emph{Methodus inveniendi lineas curvas maximi minimive
proprietate gaudentes} (1744)}

\lettrine[lines=3, findent=2pt, nindent=0pt]{N}{el capitolo precedente} abbiamo
imparato a \emph{calcolare} le derivate; in questo impariamo a \emph{usarle}. La
derivata è un'informazione locale: dice come si comporta una funzione in un intorno
piccolissimo di un punto. Il risultato centrale del capitolo, il teorema di Lagrange,
trasforma questa informazione locale in un'informazione \emph{globale}: dal segno della
derivata in tutto un intervallo si deduce se la funzione cresce o decresce, dal segno
della derivata seconda se il grafico volge la concavità verso l'alto o verso il basso.
Su queste basi costruiremo un procedimento sistematico per disegnare il grafico di una
funzione e, soprattutto, per trovarne i massimi e i minimi.

I problemi di massimo e di minimo sono il pane quotidiano dell'ingegnere. L'assetto di
efficienza massima di un velivolo, la velocità di minima resistenza e quella di minima
potenza necessaria, la sezione di una trave che a parità di materiale resiste di più,
il percorso di minimo tempo di un raggio di luce: sono tutti problemi in cui si cerca il
punto in cui la derivata di una opportuna funzione si annulla, e poi si dimostra che
quel punto è davvero un massimo o un minimo. Completeremo il capitolo con due strumenti
di calcolo di grande efficacia: il teorema di L'Hôpital, per le forme indeterminate, e
la formula di Taylor, che approssima una funzione con un polinomio e generalizza
l'approssimazione lineare del \cref{sec:retta-tangente}. Seguiremo, come sempre, il
testo di riferimento \cite{marcellini-sbordone}; in appendice dimostreremo il teorema
di Cauchy e, con esso, il teorema di L'Hôpital in tutta la sua generalità.

\medskip
\noindent\textbf{In questo capitolo.}
\begin{itemize}
  \item Massimi e minimi relativi; il teorema di Fermat (\S\ref{sec:fermat}).
  \item I teoremi di Rolle e di Lagrange (\S\ref{sec:rolle-lagrange}).
  \item Il criterio di monotonia e le sue conseguenze (\S\ref{sec:monotonia}).
  \item Funzioni convesse e concave; i punti di flesso (\S\ref{sec:convessita}).
  \item Approfondimento: la convessità senza derivate e la disuguaglianza di Jensen
  (\S\ref{subsec:convessita-corde}).
  \item Il teorema di L'Hôpital e le forme indeterminate (\S\ref{sec:hopital}).
  \item Lo studio del grafico di una funzione; gli asintoti (\S\ref{sec:studio-funzione}).
  \item La formula di Taylor e il criterio delle derivate successive (\S\ref{sec:taylor}).
  \item Appendice: il teorema di Cauchy e il teorema di L'Hôpital nel caso generale
  (\S\ref{sec:cauchy-hopital}).
  \item Esercizi di riepilogo, tra cui alcuni problemi di ottimo strutturale e ottico
  (\S\ref{sec:esercizi-cap6}).
  \item Punti fondamentali del capitolo (\S\ref{sec:sintesi-cap06}).
\end{itemize}

%% ----------------------------------------------------------------------------
%%  §6.1  Massimi e minimi relativi. Il teorema di Fermat
%% ----------------------------------------------------------------------------
\section{Massimi e minimi relativi. Il teorema di Fermat}
\label{sec:fermat}

Il teorema di Weierstrass (\cref{teo:weierstrass}) garantisce che una funzione continua
in un intervallo chiuso e limitato assume massimo e minimo, ma non dice \emph{dove}. In
questo paragrafo cominciamo a rispondere alla domanda. L'idea è semplice: in un punto
di massimo «interno» il grafico di una funzione derivabile non può né salire né
scendere, e quindi la retta tangente deve essere orizzontale. Per formulare l'idea in
modo preciso conviene introdurre una nozione di massimo più debole di quella usata fino
a ora, che guarda soltanto a ciò che accade \emph{vicino} al punto.

\begin{definizione}[Massimi e minimi relativi]\label{def:max-min-relativo}
\index{massimo!relativo}\index{minimo!relativo}\index{punto!di massimo relativo}
Sia $f:A\to\R$, con $A\subseteq\R$, e sia $x_{0}\in A$. Si dice che $x_{0}$ è un
\emph{punto di massimo relativo} (o \emph{locale}) per $f$ se esiste $\delta>0$ tale che
\begin{equation}\label{eq:max-relativo}
  f(x)\le f(x_{0}) \qquad\perogni x\in A\cap(x_{0}-\delta,x_{0}+\delta);
\end{equation}
si dice che $x_{0}$ è un \emph{punto di minimo relativo} (o \emph{locale}) per $f$ se
esiste $\delta>0$ tale che
\begin{equation}\label{eq:min-relativo}
  f(x)\ge f(x_{0}) \qquad\perogni x\in A\cap(x_{0}-\delta,x_{0}+\delta).
\end{equation}
Il valore $f(x_{0})$ si chiama, rispettivamente, \emph{massimo relativo} o \emph{minimo
relativo} di $f$. Se le disuguaglianze valgono in senso stretto per $x\neq x_{0}$, si
parla di massimo o minimo relativo \emph{stretto}. Se le disuguaglianze valgono per
\emph{ogni} $x\in A$ (non solo vicino a $x_{0}$), si parla di punto di massimo o di
minimo \emph{assoluto}. I punti di massimo e di minimo, relativo o assoluto, si dicono
anche \emph{punti di estremo}.
\end{definizione}

Ogni punto di massimo assoluto è anche di massimo relativo (basta prendere un $\delta$
qualunque), ma non viceversa. Nella \cref{fig:estremi-relativi} la funzione,
definita in $[a,b]$, ha un minimo relativo in $x_{1}$, che è anche il minimo assoluto,
un massimo relativo in $x_{2}$ e un minimo relativo in $x_{3}$, che non sono assoluti;
anche gli estremi dell'intervallo sono punti di estremo relativo: $a$ è di massimo
relativo (vicino ad $a$ i valori di $f$ sono minori di $f(a)$) e $b$ è il punto di
massimo assoluto. Si osservi che la definizione confronta $f(x_{0})$ soltanto con i
valori assunti \emph{nel dominio}: per questo ha senso anche negli estremi
dell'intervallo, dove l'intorno $(x_{0}-\delta,x_{0}+\delta)$ esce da $[a,b]$.

\begin{figure}[htb]
\centering
\begin{tikzpicture}
\begin{axis}[width=0.95\textwidth, height=5.6cm, xmin=0, xmax=6.8, ymin=0, ymax=3.8,
  axis lines=left, xtick={0.3,1.2,3.4,5.3,6.2},
  xticklabels={$a$,$x_{1}$,$x_{2}$,$x_{3}$,$b$}, ytick=\empty, clip=false]
  \fill[blu!8] (axis cs:2.6,0) rectangle (axis cs:4.2,3.55);
  \addplot[curva, domain=0.3:6.2, samples=200]
    {0.0875*x^4-1.155*x^3+4.9805*x^2-7.5684*x+4.72456};
  \addplot[tratteggio] coordinates {(2.6,2.8634) (4.2,2.8634)};
  \addplot[tratteggio] coordinates {(0.3,0) (0.3,2.8718)};
  \addplot[tratteggio] coordinates {(1.2,0) (1.2,1.0)};
  \addplot[tratteggio] coordinates {(3.4,0) (3.4,2.8634)};
  \addplot[tratteggio] coordinates {(5.3,0) (5.3,1.6031)};
  \addplot[tratteggio] coordinates {(6.2,0) (6.2,3.275)};
  \addplot[only marks, mark=*, mark size=1.6pt, rame] coordinates
    {(0.3,2.8718) (1.2,1.0) (3.4,2.8634) (5.3,1.6031) (6.2,3.275)};
  \node[font=\scriptsize, grigio, anchor=south] at (axis cs:3.4,3.56)
    {$(x_{2}-\delta,x_{2}+\delta)$};
  \node[font=\scriptsize, rame, above] at (axis cs:3.4,2.95) {max.\ rel.};
  \node[font=\scriptsize, rame, below right] at (axis cs:1.2,0.95) {min.\ ass.};
  \node[font=\scriptsize, rame, anchor=north west] at (axis cs:5.35,1.55) {min.\ rel.};
  \node[font=\scriptsize, rame, right] at (axis cs:6.25,3.275) {max.\ ass.};
  \node[font=\scriptsize, rame, anchor=south west] at (axis cs:0.3,2.95) {max.\ rel.};
\end{axis}
\end{tikzpicture}
\caption{Punti di estremo relativo e assoluto di una funzione definita in $[a,b]$. Nella
striscia evidenziata, intorno di $x_{2}$, il grafico sta sotto la retta orizzontale
$y=f(x_{2})$: $x_{2}$ è di massimo relativo, anche se la funzione assume altrove valori
maggiori di $f(x_{2})$. In $x_{1}$, $x_{2}$, $x_{3}$, interni all'intervallo, la tangente
è orizzontale; negli estremi $a$ e $b$ non lo è.}
\label{fig:estremi-relativi}
\end{figure}

La figura suggerisce la proprietà che stiamo cercando: nei punti di estremo
relativo \emph{interni} all'intervallo la tangente al grafico è orizzontale.

\begin{teorema}[di Fermat]\label{teo:fermat}\index{teorema!di Fermat}
Sia $f$ una funzione definita in un intervallo $[a,b]$ e sia $x_{0}\in(a,b)$ un punto
\emph{interno} di massimo o di minimo relativo per $f$. Se $f$ è derivabile in $x_{0}$,
allora
\begin{equation}\label{eq:fermat}
  f'(x_{0})=0 .
\end{equation}
\end{teorema}

\begin{schema}
\Ipotesi{\begin{itemize}[nosep, leftmargin=1.2em]
  \item $f:[a,b]\to\R$;
  \item $x_{0}\in(a,b)$ punto di massimo (o di minimo) relativo;
  \item $f$ derivabile in $x_{0}$.
\end{itemize}}
\Tesi{$f'(x_{0})=0$.}
\end{schema}

\begin{strategia}
Studiamo il \emph{segno} del rapporto incrementale $\frac{f(x_{0}+h)-f(x_{0})}{h}$ in un
punto di massimo. Il numeratore è $\le0$ per $h$ piccolo, qualunque sia il segno di $h$;
quindi il rapporto è $\le0$ se $h>0$ ed è $\ge0$ se $h<0$. Passando al limite, la
derivata destra è $\le0$ e la derivata sinistra è $\ge0$; ma la funzione è derivabile,
quindi le due derivate unilaterali coincidono con $f'(x_{0})$, che deve essere
contemporaneamente $\le0$ e $\ge0$. L'ipotesi che $x_{0}$ sia \emph{interno} serve a
poter prendere incrementi $h$ di entrambi i segni.
\end{strategia}

\begin{dimostrazione}
Supponiamo che $x_{0}$ sia un punto di massimo relativo (il caso del minimo si
riconduce a questo considerando la funzione $-f$, che ha un massimo relativo in $x_{0}$ e
derivata $-f'(x_{0})$). Per la \eqref{eq:max-relativo} esiste $\delta>0$ tale che
$f(x)\le f(x_{0})$ per ogni $x\in[a,b]\cap(x_{0}-\delta,x_{0}+\delta)$. Poiché $x_{0}$ è
interno, possiamo ridurre $\delta$ in modo che $\delta\le\min\{x_{0}-a,\,b-x_{0}\}$:
allora l'intero intervallo $(x_{0}-\delta,x_{0}+\delta)$ è contenuto in $[a,b]$, e
\begin{equation}\label{eq:fermat-numeratore}
  f(x_{0}+h)-f(x_{0})\le0 \qquad\perogni h\in(-\delta,\delta).
\end{equation}
Dividendo per $h$ si ottiene, per $0<h<\delta$ e per $-\delta<h<0$ rispettivamente,
\begin{equation}\label{eq:fermat-rapporti}
\begin{aligned}
  &\frac{f(x_{0}+h)-f(x_{0})}{h}\le0 \qquad\text{per } 0<h<\delta,\\[2pt]
  &\frac{f(x_{0}+h)-f(x_{0})}{h}\ge0 \qquad\text{per } -\delta<h<0.
\end{aligned}
\end{equation}
Poiché $f$ è derivabile in $x_{0}$, il rapporto incrementale tende a $f'(x_{0})$ sia per
$h\to0^{+}$ sia per $h\to0^{-}$ (\cref{prop:dx-sx}). Per il teorema del confronto dei
limiti (\cref{teo:confronto-funzioni}\,(ii)), il limite di una funzione $\le0$ è $\le0$,
e quello di una funzione $\ge0$ è $\ge0$; dalle \eqref{eq:fermat-rapporti} segue quindi
\begin{align*}
  f'(x_{0})&=\lim_{h\to0^{+}}\frac{f(x_{0}+h)-f(x_{0})}{h}\le0,\\
  f'(x_{0})&=\lim_{h\to0^{-}}\frac{f(x_{0}+h)-f(x_{0})}{h}\ge0 ,
\end{align*}
cioè $f'(x_{0})=0$.
\end{dimostrazione}

\begin{figure}[htb]
\centering
\begin{tikzpicture}
\begin{axis}[width=0.8\textwidth, height=5cm, xmin=-0.2, xmax=4.4, ymin=0.6, ymax=3.6,
  axis lines=left, xtick={1,2,3.2}, xticklabels={$x_{0}+h$,$x_{0}$,$x_{0}+h$},
  ytick=\empty, clip=false]
  \addplot[curva, domain=0:4, samples=100] {3-0.5*(x-2)^2};
  \addplot[curva3, domain=0.2:2.8] {3+0.5*(x-2)};
  \addplot[curva2, domain=1.2:4.2] {3-0.6*(x-2)};
  \addplot[tratteggio] coordinates {(0.4,3) (3.6,3)};
  \addplot[tratteggio] coordinates {(1,0.6) (1,2.5)};
  \addplot[tratteggio] coordinates {(3.2,0.6) (3.2,2.28)};
  \addplot[only marks, mark=*, mark size=1.6pt, blu] coordinates {(1,2.5) (2,3) (3.2,2.28)};
  \node[font=\scriptsize, verde!50!black, anchor=south east] at (axis cs:0.55,2.3) {$h<0$};
  \node[font=\scriptsize, rame, anchor=west] at (axis cs:4.25,1.68) {$h>0$};
\end{axis}
\end{tikzpicture}
\caption{La dimostrazione del teorema di Fermat. Vicino a un punto di massimo relativo
interno $x_{0}$, le secanti con $h<0$ (verde) hanno pendenza $\ge0$, quelle con $h>0$
(rame) pendenza $\le0$; al limite, la pendenza della tangente è insieme $\ge0$ e $\le0$.}
\label{fig:fermat}
\end{figure}

\begin{attenzione}[Il teorema di Fermat dà una condizione necessaria, non sufficiente]
Il viceversa del teorema di Fermat è falso: la derivata può annullarsi in un punto
che non è né di massimo né di minimo. L'esempio tipico è $f(x)=x^{3}$: risulta
$f'(0)=0$, ma $f(x)<0=f(0)$ per $x<0$ e $f(x)>0=f(0)$ per $x>0$, quindi in ogni intorno
di $0$ la funzione assume valori sia maggiori sia minori di $f(0)$. Inoltre il teorema
non dice nulla nei punti in cui $f$ non è derivabile: $f(x)=\abs x$ ha in $0$ il minimo
assoluto, ma non ha derivata in $0$ (\cref{es:modulo-non-derivabile}). Infine, l'ipotesi
che $x_{0}$ sia interno è essenziale: $f(x)=x$ in $[0,1]$ ha il massimo in $x_{0}=1$,
dove la derivata vale $1$.
\end{attenzione}

I punti $x_{0}$ in cui $f'(x_{0})=0$ si chiamano \emph{punti stazionari} (o
\emph{critici}) di $f$\index{punto!stazionario}. Che cosa accade negli estremi
dell'intervallo? Il ragionamento della dimostrazione si può ancora fare, ma con
incrementi di un solo segno, e fornisce una disuguaglianza invece di un'uguaglianza.

\begin{proposizione}[Estremi relativi e segno della derivata]\label{prop:fermat-estremi}
Sia $f:[a,b]\to\R$ e sia $x_{0}\in[a,b]$ un punto di massimo relativo per $f$, nel quale
$f$ è derivabile (con derivata unilaterale se $x_{0}=a$ oppure $x_{0}=b$). Allora
\begin{equation}\label{eq:fermat-max-generale}
  f'(x_{0})\,(x-x_{0})\le0 \qquad\perogni x\in[a,b].
\end{equation}
Se invece $x_{0}$ è un punto di minimo relativo, risulta
\begin{equation}\label{eq:fermat-min-generale}
  f'(x_{0})\,(x-x_{0})\ge0 \qquad\perogni x\in[a,b].
\end{equation}
In particolare: se $a$ è di massimo relativo allora $f'(a)\le0$, se $b$ è di massimo
relativo allora $f'(b)\ge0$; se $a$ è di minimo relativo allora $f'(a)\ge0$, se $b$ è di
minimo relativo allora $f'(b)\le0$; se $x_{0}\in(a,b)$ allora $f'(x_{0})=0$.
\end{proposizione}

\begin{schema}
\Ipotesi{$f:[a,b]\to\R$; \ $x_{0}\in[a,b]$ di massimo (minimo) relativo; \ $f$
derivabile in $x_{0}$ (unilateralmente se $x_{0}$ è un estremo).}
\Tesi{$f'(x_{0})(x-x_{0})\le0$ \ (rispettivamente $\ge0$) \ per ogni $x\in[a,b]$.}
\end{schema}

\begin{strategia}
Ripetiamo la dimostrazione del teorema di Fermat tenendo conto che, se $x_{0}=a$, sono
ammessi soltanto incrementi $h>0$, e se $x_{0}=b$ soltanto incrementi $h<0$. La forma
\eqref{eq:fermat-max-generale} riassume in una sola disuguaglianza tutti i casi: il
prodotto $f'(x_{0})(x-x_{0})$ è la variazione «prevista dalla tangente» quando ci si
sposta da $x_{0}$ a $x$, e in un punto di massimo non può essere positiva.
\end{strategia}

\begin{dimostrazione}
Sia $x_{0}$ di massimo relativo e sia $\delta>0$ come nella \eqref{eq:max-relativo}.
Se $x_{0}<b$, per $0<h<\min\{\delta,b-x_{0}\}$ il rapporto incrementale è $\le0$, e
passando al limite per $h\to0^{+}$ si ottiene $f'(x_{0})\le0$. Se $x_{0}>a$, per
$-\min\{\delta,x_{0}-a\}<h<0$ il rapporto incrementale è $\ge0$, e al limite per
$h\to0^{-}$ si ottiene $f'(x_{0})\ge0$. Quindi:
\begin{itemize}
  \item se $x_{0}=a$, vale solo il primo caso: $f'(a)\le0$, e poiché $x-a\ge0$ per ogni
  $x\in[a,b]$ risulta $f'(a)(x-a)\le0$;
  \item se $x_{0}=b$, vale solo il secondo caso: $f'(b)\ge0$, e poiché $x-b\le0$ si ha
  $f'(b)(x-b)\le0$;
  \item se $a<x_{0}<b$, valgono entrambi: $f'(x_{0})=0$, e la
  \eqref{eq:fermat-max-generale} è banalmente vera.
\end{itemize}
Il caso del minimo si ottiene applicando quanto provato alla funzione $-f$.
\end{dimostrazione}

\subsection{Un metodo per trovare massimi e minimi assoluti}

Combinando il teorema di Weierstrass con quello di Fermat otteniamo un procedimento
concreto. Sia $f$ continua in $[a,b]$ e derivabile in $(a,b)$, eccetto al più un numero
finito di punti. Per il teorema di Weierstrass $f$ ammette massimo e minimo assoluti;
un punto $x_{0}$ di massimo o di minimo assoluto è anche di estremo relativo, e
quindi, per il teorema di Fermat, deve appartenere a una delle tre categorie seguenti:
\begin{enumerate}[label=\textup{(\roman*)}]
  \item $x_{0}$ è uno degli estremi $a$, $b$ dell'intervallo;
  \item $x_{0}$ è un punto interno in cui $f$ non è derivabile;
  \item $x_{0}$ è un punto interno stazionario, cioè $f'(x_{0})=0$.
\end{enumerate}
Nei casi più comuni i punti delle tre categorie sono in numero finito: basta allora
calcolare i valori di $f$ in tali punti e confrontarli. Il più grande è il massimo
assoluto, il più piccolo è il minimo assoluto. Si noti che il ragionamento è
\emph{per esclusione}: non serve sapere a priori quali punti stazionari siano di
massimo o di minimo, perché l'esistenza degli estremi assoluti è garantita dal teorema
di Weierstrass. Il programma seguente esegue il procedimento con \texttt{sympy} per la
funzione $f(x)=x^{3}-3x$ nell'intervallo $[-2,3]$.

\begin{python}[Massimo e minimo assoluti in un intervallo chiuso]
import sympy as sp

x = sp.symbols("x", real=True)

def estremi_assoluti(f, a, b):
    """Massimo e minimo assoluti di f (derivabile) in [a, b]."""
    critici = sp.solveset(sp.diff(f, x), x, sp.Interval.open(a, b))
    candidati = sorted(set(critici) | {a, b})
    valori = {c: f.subs(x, c) for c in candidati}
    for c in candidati:
        print(f"  f({c}) = {valori[c]}")
    M, m = max(valori.values()), min(valori.values())
    dove = lambda y: [c for c in candidati if valori[c] == y]
    print("  max =", M, "in", dove(M))
    print("  min =", m, "in", dove(m))

print("f(x) = x^3 - 3x in [-2, 3]:")
estremi_assoluti(x**3 - 3*x, -2, 3)
\end{python}
\begin{risultato}
f(x) = x^3 - 3x in [-2, 3]:
  f(-2) = -2
  f(-1) = 2
  f(1) = -2
  f(3) = 18
  max = 18 in [3]
  min = -2 in [-2, 1]
\end{risultato}

Il massimo assoluto $18$ è assunto nell'estremo $b=3$, dove la derivata
$f'(3)=24$ non è nulla (ed è positiva, in accordo con la \cref{prop:fermat-estremi}); il
minimo assoluto $-2$ è assunto in due punti: nell'estremo $a=-2$, dove $f'(-2)=9\ge0$,
e nel punto stazionario interno $x=1$. Il punto stazionario $x=-1$ è di massimo
relativo, ma non assoluto. Ritroveremo il grafico di questa funzione nella
\cref{fig:cubica-monotonia}.

\begin{esempio}[La gittata massima]\label{es:gittata}
Un carico viene lanciato con velocità iniziale di modulo $v$, inclinata dell'angolo
$\theta$ sull'orizzontale. Trascurando la resistenza dell'aria, la distanza orizzontale
percorsa prima di tornare alla quota di partenza (la \emph{gittata}) è
\begin{equation}\label{eq:gittata}
  R(\theta)=\frac{v^{2}\sen2\theta}{g}, \qquad \theta\in\Bigl[0,\frac\pi2\Bigr],
\end{equation}
dove $g$ è l'accelerazione di gravità. La funzione $R$ è continua nell'intervallo
chiuso $[0,\frac\pi2]$, quindi ha massimo assoluto; è derivabile, con
$R'(\theta)=\frac{2v^{2}}{g}\cos2\theta$, che nell'intervallo aperto si annulla soltanto
per $2\theta=\frac\pi2$, cioè $\theta=\frac\pi4$. I candidati sono dunque $0$,
$\frac\pi4$, $\frac\pi2$, con $R(0)=R(\frac\pi2)=0$ e $R(\frac\pi4)=\frac{v^{2}}g$: la
gittata massima si ottiene con l'inclinazione di $45^{\circ}$, mentre i due estremi
sono punti di minimo assoluto. Coerentemente con la \cref{prop:fermat-estremi}, risulta
$R'(0)=\frac{2v^{2}}g\ge0$ (minimo nell'estremo sinistro) e
$R'(\frac\pi2)=-\frac{2v^{2}}g\le0$ (minimo nell'estremo destro).
\end{esempio}

\begin{python}[La gittata: punti critici e derivate negli estremi]
import sympy as sp

th, v, g = sp.symbols("theta v g", positive=True)
R = v**2 * sp.sin(2 * th) / g             # gittata nel vuoto
dR = sp.diff(R, th)
print("R'(theta) =", dR)
print("punti critici in (0, pi/2):",
      sp.solveset(dR, th, sp.Interval.open(0, sp.pi / 2)))
print("R'(0) =", dR.subs(th, 0),
      "   R'(pi/2) =", dR.subs(th, sp.pi / 2))
print("R(pi/4) =", R.subs(th, sp.pi / 4))
\end{python}
\begin{risultato}
R'(theta) = 2*v**2*cos(2*theta)/g
punti critici in (0, pi/2): {pi/4}
R'(0) = 2*v**2/g    R'(pi/2) = -2*v**2/g
R(pi/4) = v**2/g
\end{risultato}

\begin{idea}[L'efficienza massima, finalmente con tutti i crismi]
Nell'idea «Ritorno all'efficienza massima» del \cref{sec:operazioni-derivate} abbiamo
visto che la derivata dell'efficienza $E(C_{L})=\frac{C_{L}}{C_{D0}+kC_{L}^{2}}$ si annulla
in $C_{L}^{*}=\sqrt{C_{D0}/k}$. Ora possiamo concludere con il metodo appena descritto:
$E$ è continua in $[0,C_{L,\max}]$ (con $C_{L,\max}>C_{L}^{*}$), quindi ha massimo
assoluto; l'unico punto stazionario interno è $C_{L}^{*}$; negli estremi
$E(0)=0<E(C_{L}^{*})$ e $E(C_{L,\max})<E(C_{L}^{*})$, perché, come vedremo nel
\cref{sec:monotonia}, $E$ è strettamente decrescente per $C_{L}\ge C_{L}^{*}$. Il massimo assoluto è
quindi $E_{\max}=E(C_{L}^{*})=\frac1{2\sqrt{kC_{D0}}}$, in accordo con il
\cref{cap:funzioni-continue}.
\end{idea}

\begin{notastorica}[Fermat e il metodo dell'«adequalitas»]
Pierre de Fermat (1601?--1665), magistrato al parlamento di Tolosa e matematico per
passione, descrisse intorno al 1636--1638, in un breve scritto intitolato
\emph{Methodus ad disquirendam maximam et minimam} fatto circolare per lettera, un
procedimento per trovare i massimi e i minimi. In linguaggio moderno: si confronta
$f(x)$ con $f(x+E)$, si «adeguagliano» le due espressioni (\emph{adaequalitas}), si
divide per $E$ e infine si pone $E=0$. Nel primo esempio si divide un segmento di
lunghezza $b$ in due parti in modo che il rettangolo $x(b-x)$ sia massimo: da
$x(b-x)\sim(x+E)(b-x-E)$ segue $E(b-2x)-E^{2}\sim0$, poi $b-2x-E\sim0$, e per $E=0$ si
trova $x=\frac b2$: il rettangolo massimo è il quadrato. È, in sostanza, la condizione
$f'(x)=0$, ottenuta decenni prima del calcolo differenziale; ma dividere per $E$ e poi
porre $E=0$ non è passare al limite. Per i polinomi le due operazioni si equivalgono; se
invece $f$ contiene radicali, come $x+\sqrt{1-x}$, prima di dividere per $E$ bisogna
eliminarli elevando al quadrato, e nei casi meno semplici i calcoli diventano
proibitivi. Il nuovo calcolo superò l'ostacolo separando due momenti, il calcolo di $f'$
con regole generali e la soluzione di $f'(x)=0$: non a caso Leibniz vantava un metodo
«che non si arresta davanti alle quantità fratte o irrazionali» (epigrafe del
\cref{cap:derivate}). Già Keplero, nella \emph{Nova stereometria doliorum vinariorum}
(1615), aveva notato che vicino a un massimo le variazioni di una grandezza diventano
impercettibili. Fermat applicò lo stesso tipo di ragionamento all'ottica, enunciando nel
1662 il \emph{principio del tempo minimo}: la
luce, per andare da un punto a un altro, segue il percorso che richiede il tempo
minore (si veda l'\cref{es:snell}). Le sue opere furono pubblicate postume dal figlio
Samuel, nei \emph{Varia opera mathematica} (1679).
\end{notastorica}

%% ----------------------------------------------------------------------------
%%  §6.2  I teoremi di Rolle e di Lagrange
%% ----------------------------------------------------------------------------
\section{I teoremi di Rolle e di Lagrange}
\label{sec:rolle-lagrange}

Il teorema di Fermat riguarda un singolo punto. I due teoremi di questo paragrafo
riguardano invece un intero intervallo, e sono la chiave di volta di tutto il capitolo:
mettono in relazione la variazione \emph{complessiva} $f(b)-f(a)$ di una funzione con
il valore della derivata in un punto opportuno. Il primo è un caso particolare del
secondo, ma è anche lo strumento con cui il secondo si dimostra.

\begin{teorema}[di Rolle]\label{teo:rolle}\index{teorema!di Rolle}
Sia $f$ una funzione continua nell'intervallo chiuso $[a,b]$ e derivabile
nell'intervallo aperto $(a,b)$. Se $f(a)=f(b)$, esiste un punto $x_{0}\in(a,b)$ tale
che
\begin{equation}\label{eq:rolle}
  f'(x_{0})=0 .
\end{equation}
\end{teorema}

\begin{schema}
\Ipotesi{\begin{itemize}[nosep, leftmargin=1.2em]
  \item $f$ continua in $[a,b]$;
  \item $f$ derivabile in $(a,b)$;
  \item $f(a)=f(b)$.
\end{itemize}}
\Tesi{esiste $x_{0}\in(a,b)$ con $f'(x_{0})=0$.}
\end{schema}

\begin{strategia}
Cerchiamo $x_{0}$ tra i punti di massimo e di minimo di $f$. La continuità
nell'intervallo chiuso e limitato ci permette di usare il teorema di Weierstrass: esistono
un punto di massimo e un punto di minimo assoluti. Se almeno uno dei due è
\emph{interno}, il teorema di Fermat (applicabile, perché $f$ è derivabile in $(a,b)$)
dà subito $f'(x_{0})=0$. Resta il caso in cui massimo e minimo cadono entrambi negli
estremi: qui interviene l'ipotesi $f(a)=f(b)$, che costringe massimo e minimo a
coincidere, cioè $f$ a essere costante; e una funzione costante ha derivata nulla
ovunque.
\end{strategia}

\begin{dimostrazione}
Poiché $f$ è continua in $[a,b]$, per il teorema di Weierstrass (\cref{teo:weierstrass})
esistono $x_{1},x_{2}\in[a,b]$ tali che
\begin{equation}\label{eq:rolle-weierstrass}
  m=f(x_{1})\le f(x)\le f(x_{2})=M \qquad\perogni x\in[a,b].
\end{equation}
Distinguiamo due casi.

\emph{Primo caso: almeno uno dei due punti $x_{1}$, $x_{2}$ è interno ad $[a,b]$.}
Supponiamo, per fissare le idee, che sia $x_{2}\in(a,b)$. Allora $x_{2}$ è un punto di
massimo (assoluto, e quindi relativo) interno, nel quale $f$ è derivabile; per il
teorema di Fermat (\cref{teo:fermat}) risulta $f'(x_{2})=0$, e la tesi vale con
$x_{0}=x_{2}$. Se invece è interno $x_{1}$, si ragiona allo stesso modo con il minimo.

\emph{Secondo caso: entrambi i punti $x_{1}$, $x_{2}$ sono estremi dell'intervallo.}
Allora $m$ e $M$ sono due dei valori $f(a)$, $f(b)$, che per ipotesi sono uguali;
quindi $m=M$, e dalla \eqref{eq:rolle-weierstrass} segue $f(x)=M$ per ogni $x\in[a,b]$.
La funzione è costante, e la sua derivata è nulla in ogni punto
\eqref{eq:derivata-costante}: la tesi vale con un qualunque $x_{0}\in(a,b)$.
\end{dimostrazione}

\begin{figure}[htb]
\centering
\begin{tikzpicture}
\begin{axis}[width=0.85\textwidth, height=5cm, xmin=0, xmax=6.4, ymin=0, ymax=3.3,
  axis lines=left, xtick={0.5,1.809,4.427,5.736},
  xticklabels={$a$,$x_{0}$,$x_{0}'$,$b$}, ytick=\empty, clip=false]
  \addplot[curva, domain=0.5:5.736, samples=150] {2+0.8*sin(deg(1.2*(x-0.5)))};
  \addplot[curva2, domain=1.0:2.6] {2.8};
  \addplot[curva2, domain=3.6:5.2] {1.2};
  \addplot[tratteggio] coordinates {(0.5,0) (0.5,2) (5.736,2) (5.736,0)};
  \addplot[tratteggio] coordinates {(1.809,0) (1.809,2.8)};
  \addplot[tratteggio] coordinates {(4.427,0) (4.427,1.2)};
  \addplot[only marks, mark=*, mark size=1.6pt, blu] coordinates
    {(0.5,2) (5.736,2) (1.809,2.8) (4.427,1.2)};
  \node[font=\scriptsize, grigio, anchor=west] at (axis cs:5.8,2.0) {$f(a)=f(b)$};
\end{axis}
\end{tikzpicture}
\caption{Il teorema di Rolle: se $f(a)=f(b)$, il grafico ha almeno un punto con
tangente orizzontale. Qui i punti sono due: il massimo $x_{0}$ e il minimo $x_{0}'$; il
teorema garantisce l'esistenza di almeno uno, non l'unicità.}
\label{fig:rolle}
\end{figure}

\begin{osservazione}[Le ipotesi del teorema di Rolle sono tutte necessarie]
\label{oss:rolle-ipotesi}
Se si toglie una qualunque delle tre ipotesi la tesi può essere falsa
(\cref{fig:rolle-ipotesi}).
\begin{enumerate}[label=\textup{(\alph*)}]
  \item \emph{Continuità in $[a,b]$.} La funzione $f(x)=x$ per $x\in[0,1)$, con
  $f(1)=0$, è derivabile in $(0,1)$, con derivata $1$, e verifica $f(0)=f(1)$; ma è
  discontinua in $x=1$, e la sua derivata non si annulla mai. La continuità serviva per
  il teorema di Weierstrass: infatti $f$ non ha massimo in $[0,1]$.
  \item \emph{Derivabilità in $(a,b)$.} La funzione $f(x)=1-\abs x$ in $[-1,1]$ è
  continua, verifica $f(-1)=f(1)=0$ e ha massimo nel punto interno $0$; ma lì non è
  derivabile, e altrove la derivata vale $\pm1$. Il teorema di Fermat non era
  applicabile.
  \item \emph{Uguaglianza $f(a)=f(b)$.} La funzione $f(x)=x$ in $[0,1]$ è continua e
  derivabile, ma $f'(x)=1$ per ogni $x$. Massimo e minimo cadono negli estremi senza che
  la funzione sia costante.
\end{enumerate}
Si noti invece che \emph{non} si richiede la derivabilità negli estremi: la funzione
$f(x)=\sqrt{1-x^{2}}$ in $[-1,1]$ (il cui grafico è una semicirconferenza) ha
derivata infinita in $\pm1$, ma verifica le ipotesi, e infatti $f'(0)=0$.
\end{osservazione}

\begin{figure}[htb]
\centering
\begin{tikzpicture}
\begin{axis}[name=r1, width=0.36\textwidth, height=0.32\textwidth, axis lines=middle,
  xmin=-0.3, xmax=1.35, ymin=-0.25, ymax=1.3, xtick={1}, ytick=\empty,
  title={\scriptsize (a) discontinua in $b$}]
  \addplot[curva, domain=0:1, samples=2] {x};
  \addplot[only marks, mark=o, mark size=1.6pt, blu, fill=white] coordinates {(1,1)};
  \addplot[only marks, mark=*, mark size=1.6pt, blu] coordinates {(1,0) (0,0)};
\end{axis}
\begin{axis}[at={(r1.east)}, anchor=west, xshift=0.6cm, name=r2, width=0.36\textwidth,
  height=0.32\textwidth, axis lines=middle, xmin=-1.35, xmax=1.35, ymin=-0.25,
  ymax=1.3, xtick={-1,1}, ytick=\empty, title={\scriptsize (b) non derivabile in $0$}]
  \addplot[curva, domain=-1:1, samples=3] {1-abs(x)};
\end{axis}
\begin{axis}[at={(r2.east)}, anchor=west, xshift=0.6cm, width=0.36\textwidth,
  height=0.32\textwidth, axis lines=middle, xmin=-0.3, xmax=1.35, ymin=-0.25,
  ymax=1.3, xtick={1}, ytick=\empty, title={\scriptsize (c) $f(a)\neq f(b)$}]
  \addplot[curva, domain=0:1, samples=2] {x};
  \addplot[only marks, mark=*, mark size=1.6pt, blu] coordinates {(0,0) (1,1)};
\end{axis}
\end{tikzpicture}
\caption{Tre funzioni che violano una sola delle ipotesi del teorema di Rolle e per le
quali la tesi è falsa: nessuna ha un punto a tangente orizzontale.}
\label{fig:rolle-ipotesi}
\end{figure}

Il teorema di Rolle ha un'interpretazione geometrica evidente: se la corda che
congiunge gli estremi del grafico è orizzontale, in qualche punto anche la tangente lo
è. Se ora «inclinassimo» la figura, ci aspetteremmo di trovare un punto in cui la
tangente è \emph{parallela alla corda}. È il contenuto del teorema seguente.

\begin{teorema}[di Lagrange, o del valor medio]\label{teo:lagrange}
\index{teorema!di Lagrange}\index{teorema!del valor medio}
Sia $f$ una funzione continua nell'intervallo chiuso $[a,b]$ e derivabile
nell'intervallo aperto $(a,b)$. Esiste un punto $x_{0}\in(a,b)$ tale che
\begin{equation}\label{eq:lagrange}
  \frac{f(b)-f(a)}{b-a}=f'(x_{0}) .
\end{equation}
\end{teorema}

\begin{schema}
\Ipotesi{$f$ continua in $[a,b]$ e derivabile in $(a,b)$.}
\Tesi{esiste $x_{0}\in(a,b)$ con $f(b)-f(a)=f'(x_{0})\,(b-a)$.}
\end{schema}

\begin{strategia}
Il primo membro della \eqref{eq:lagrange} è la pendenza della corda che congiunge
$(a,f(a))$ e $(b,f(b))$; il secondo è la pendenza della tangente in $x_{0}$. Ci
riconduciamo al teorema di Rolle \emph{sottraendo la corda}: la funzione ausiliaria
$g(x)=f(x)-r(x)$, dove $y=r(x)$ è l'equazione della retta secante, misura la distanza
verticale tra il grafico e la corda. Questa distanza è nulla negli estremi, quindi $g$
verifica le ipotesi del teorema di Rolle; nel punto in cui $g'$ si annulla, la
pendenza di $f$ eguaglia quella della retta. È un procedimento tipico: per dimostrare
un risultato più generale, si costruisce una funzione alla quale si può applicare il
caso particolare già noto.
\end{strategia}

\begin{dimostrazione}
La retta passante per i punti $(a,f(a))$ e $(b,f(b))$ ha equazione
$y=r(x)=f(a)+\frac{f(b)-f(a)}{b-a}(x-a)$. Consideriamo la funzione
\begin{equation}\label{eq:lagrange-g}
  g(x)=f(x)-\Bigl[f(a)+\frac{f(b)-f(a)}{b-a}(x-a)\Bigr], \qquad x\in[a,b].
\end{equation}
La funzione $g$ è continua in $[a,b]$ e derivabile in $(a,b)$, perché differenza di $f$
e di un polinomio di primo grado (\cref{teo:operazioni-continue,teo:operazioni-derivate});
inoltre
\[
  g(a)=f(a)-f(a)=0, \qquad g(b)=f(b)-\bigl[f(a)+f(b)-f(a)\bigr]=0 .
\]
Per il teorema di Rolle (\cref{teo:rolle}) esiste $x_{0}\in(a,b)$ tale che
$g'(x_{0})=0$. Poiché
\[
  g'(x)=f'(x)-\frac{f(b)-f(a)}{b-a} \qquad\perogni x\in(a,b),
\]
la condizione $g'(x_{0})=0$ è esattamente la \eqref{eq:lagrange}.
\end{dimostrazione}

\begin{figure}[htb]
\centering
\begin{tikzpicture}
\begin{axis}[width=0.85\textwidth, height=5.4cm, xmin=0, xmax=4.8, ymin=0, ymax=4.6,
  axis lines=left, xtick={0.25,1.5625,4}, xticklabels={$a$,$x_{0}$,$b$},
  ytick=\empty, clip=false]
  \addplot[curva, domain=0.1:4.3, samples=120] {2*sqrt(x)};
  \addplot[curva2, domain=0.25:4] {1+0.8*(x-0.25)};
  \addplot[curva3, domain=0.2:3.6] {2.5+0.8*(x-1.5625)};
  \addplot[tratteggio] coordinates {(0.25,0) (0.25,1)};
  \addplot[tratteggio] coordinates {(4,0) (4,4)};
  \addplot[tratteggio] coordinates {(1.5625,0) (1.5625,2.5)};
  \addplot[only marks, mark=*, mark size=1.6pt, blu] coordinates
    {(0.25,1) (4,4) (1.5625,2.5)};
  \node[font=\scriptsize, rame, anchor=north west] at (axis cs:2.5,2.75) {corda};
  \node[font=\scriptsize, verde!50!black, anchor=south east] at (axis cs:3.2,3.95)
    {tangente parallela};
\end{axis}
\end{tikzpicture}
\caption{Il teorema di Lagrange per $f(x)=2\sqrt x$ in $[a,b]=[\frac14,4]$: la corda ha
pendenza $\frac{4-1}{4-1/4}=\frac45$, e la tangente nel punto $x_{0}=\frac{25}{16}$, dove
$f'(x_{0})=\frac1{\sqrt{x_{0}}}=\frac45$, le è parallela.}
\label{fig:lagrange}
\end{figure}

\begin{idea}[Il «Tutor» e il velivolo: velocità media e velocità istantanea]
Se $s(t)$ è la legge oraria di un moto rettilineo, il primo membro della
\eqref{eq:lagrange} è la velocità media nell'intervallo di tempo $[a,b]$, il secondo la
velocità istantanea all'istante $x_{0}$. Il teorema di Lagrange afferma dunque che
\emph{esiste un istante in cui la velocità istantanea è uguale alla velocità media}.
È il principio su cui si fonda il sistema «Tutor» che controlla la velocità sulle
autostrade italiane: il sistema misura soltanto i tempi di passaggio sotto due portali
distanti, ad esempio, $20$~km; se l'automobile li ha percorsi in $8$~min, la
velocità media è di $150$~km/h e, per il teorema di Lagrange, il tachimetro ha
segnato $150$~km/h in almeno un istante, anche se nessuno l'ha visto. Allo stesso
modo, un velivolo che percorre i $300$~km tra due punti di riporto in un'ora ha
avuto, almeno in un istante, una velocità al suolo (lungo la rotta) di esattamente
$300$~km/h.
\end{idea}

Il punto $x_{0}$ di cui il teorema garantisce l'esistenza in generale non si può
calcolare esplicitamente, ma spesso si può trovare numericamente risolvendo
l'equazione $f'(x)=\frac{f(b)-f(a)}{b-a}$, ad esempio con il metodo di bisezione del
\cref{sec:bisezione} o con la funzione \texttt{brentq} di \texttt{scipy}, purché $f'$
sia continua e la differenza $f'(x)-\frac{f(b)-f(a)}{b-a}$ abbia segni opposti negli
estremi (per $x^{3}$ in $[-1,1]$, ad esempio, ciò non accade, benché i punti
$\pm\frac1{\sqrt3}$ esistano). Il programma
seguente lo fa per $f(x)=x^{3}$ in $[0,2]$, dove l'equazione $3x^{2}=4$ dà
$x_{0}=\frac2{\sqrt3}$, e per $f(x)=\sen x$ in $[0,\frac\pi2]$, dove $\cos x_{0}=\frac2\pi$.

\begin{python}[Il punto di Lagrange calcolato numericamente]
import math
from scipy.optimize import brentq

def xi_lagrange(f, df, a, b):
    """Un punto xi in (a, b) con f'(xi) = [f(b) - f(a)] / (b - a)."""
    m = (f(b) - f(a)) / (b - a)          # pendenza della corda
    return brentq(lambda t: df(t) - m, a + 1e-12, b - 1e-12), m

xi, m = xi_lagrange(lambda t: t**3, lambda t: 3 * t**2, 0, 2)
print(f"x^3:   corda = {m:.4f}   xi = {xi:.6f}"
      f"   2/sqrt(3)    = {2 / math.sqrt(3):.6f}")
xi, m = xi_lagrange(math.sin, math.cos, 0, math.pi / 2)
print(f"sen x: corda = {m:.4f}   xi = {xi:.6f}"
      f"   arccos(2/pi) = {math.acos(2 / math.pi):.6f}")

# stima di sqrt(101) con il teorema di Lagrange (vedi il testo)
sup = 10 + 1 / 20
inf = 10 + 1 / (2 * sup)
print(f"{inf:.6f} < sqrt(101) = {math.sqrt(101):.6f} < {sup:.6f}")
\end{python}
\begin{risultato}
x^3:   corda = 4.0000   xi = 1.154701   2/sqrt(3)    = 1.154701
sen x: corda = 0.6366   xi = 0.880689   arccos(2/pi) = 0.880689
10.049751 < sqrt(101) = 10.049876 < 10.050000
\end{risultato}

Posto $a=x_{0}$, $b=x_{0}+h$ (con $h>0$) e scritto il punto intermedio nella forma
$x_{0}+\vartheta h$, con $0<\vartheta<1$, la tesi del teorema di Lagrange diventa la
\emph{formula degli accrescimenti finiti}: esiste $\vartheta\in(0,1)$ tale che
\begin{equation}\label{eq:accrescimenti-finiti}
  f(x_{0}+h)=f(x_{0})+f'(x_{0}+\vartheta h)\,h ,
\end{equation}
valida, con lo stesso ragionamento applicato all'intervallo $[x_{0}+h,x_{0}]$, anche
per $h<0$: in entrambi i casi basta che $f$ sia continua nell'intervallo chiuso di
estremi $x_{0}$ e $x_{0}+h$ e derivabile nei suoi punti interni. Il confronto con
l'approssimazione lineare del \cref{sec:retta-tangente},
$f(x_{0}+h)=f(x_{0})+f'(x_{0})h+o(h)$, è istruttivo: lì la derivata è calcolata nel
punto noto $x_{0}$ e si commette un errore $o(h)$; qui l'uguaglianza è esatta, ma la
derivata è calcolata in un punto sconosciuto. Proprio questo la rende adatta a
ottenere \emph{stime}, come nel corollario seguente.

\needspace{9\baselineskip}
\begin{corollario}[Stima degli incrementi]\label{cor:lipschitz}
Sia $f$ derivabile in un intervallo $I$ e supponiamo che esista $M\ge0$ tale che
$\abs{f'(x)}\le M$ per ogni $x\in I$. Allora
\begin{equation}\label{eq:lipschitz}
  \abs{f(x)-f(y)}\le M\abs{x-y} \qquad\perogni x,y\in I .
\end{equation}
\end{corollario}

\begin{schema}
\Ipotesi{$f$ derivabile nell'intervallo $I$; \ $\abs{f'}\le M$ in $I$.}
\Tesi{$\abs{f(x)-f(y)}\le M\abs{x-y}$ per ogni $x,y\in I$.}
\end{schema}

\begin{strategia}
Basta applicare il teorema di Lagrange all'intervallo di estremi $x$ e $y$ (contenuto in
$I$, perché $I$ è un intervallo) e maggiorare la derivata nel punto intermedio.
\end{strategia}

\begin{dimostrazione}
Se $x=y$ non c'è nulla da dimostrare. Se, ad esempio, $x<y$, l'intervallo $[x,y]$ è
contenuto in $I$, e in esso $f$ è derivabile, quindi anche continua
(\cref{teo:derivabile-continua}). Per il teorema di Lagrange esiste $x_{0}\in(x,y)$ con
$f(y)-f(x)=f'(x_{0})(y-x)$, e quindi
$\abs{f(y)-f(x)}=\abs{f'(x_{0})}\,\abs{y-x}\le M\abs{y-x}$.
\end{dimostrazione}

Una funzione che verifica la \eqref{eq:lipschitz} per qualche costante $M$ si dice
\emph{lipschitziana}\index{funzione!lipschitziana} in $I$, dal nome del matematico tedesco
Rudolf Lipschitz. Poiché $\abs{D\sen x}=\abs{\cos x}\le1$ e
$\abs{D\operatorname{arctg}x}=\frac1{1+x^{2}}\le1$, si ottengono le disuguaglianze, valide
per ogni $x,y\in\R$,
\begin{equation}\label{eq:sen-arctg-lipschitz}
  \abs{\sen x-\sen y}\le\abs{x-y}, \qquad
  \abs{\operatorname{arctg}x-\operatorname{arctg}y}\le\abs{x-y} ,
\end{equation}
la prima delle quali era stata ricavata per via trigonometrica nell'\cref{es:seno-lipschitz}.
In termini applicativi il corollario è una regola di \emph{propagazione degli errori}:
se una grandezza $x$ è misurata con un errore non superiore a $\varepsilon$, il valore
calcolato $f(x)$ è affetto da un errore non superiore a $M\varepsilon$, dove $M$ è il
massimo di $\abs{f'}$ nell'intervallo di incertezza.

\begin{esempio}[Una stima di $\sqrt{101}$]\label{es:stima-radice}
Applichiamo il teorema di Lagrange a $f(x)=\sqrt x$ nell'intervallo $[100,101]$:
esiste $x_{0}\in(100,101)$ tale che
\[
  \sqrt{101}-10=\frac{1}{2\sqrt{x_{0}}}\,(101-100)=\frac{1}{2\sqrt{x_{0}}} .
\]
Poiché $\sqrt{x_{0}}>10$, risulta $\sqrt{101}-10<\frac1{20}$, cioè $\sqrt{101}<10{,}05$.
Ma allora anche $\sqrt{x_{0}}<\sqrt{101}<10{,}05$, e quindi
$\sqrt{101}-10>\frac{1}{20{,}1}$. In conclusione
$10{,}04975<\sqrt{101}<10{,}05$, come conferma l'ultima riga dell'output precedente.
Si noti come la stima superiore sia stata usata per migliorare quella inferiore.
\end{esempio}

\needspace{8\baselineskip}
\begin{notastorica}[Rolle, Lagrange, Bonnet]
Michel Rolle (1652--1719), membro dell'Académie des Sciences, enunciò nel 1691 il
teorema che porta il suo nome soltanto per i \emph{polinomi}, nel contesto della
ricerca delle radici delle equazioni algebriche: tra due radici consecutive di un
polinomio cade una radice di un polinomio da esso derivato secondo regole puramente
algebriche (la «cascata» di Rolle). Paradossalmente, Rolle fu per anni uno dei più
accesi critici del nuovo calcolo infinitesimale, che giudicava privo di fondamenti
solidi. Il nome «teorema di Rolle» fu attribuito soltanto nell'Ottocento. Il teorema
del valor medio compare nella \emph{Théorie des fonctions analytiques} (1797) di
Joseph-Louis Lagrange, in una forma legata allo sviluppo in serie delle funzioni;
Cauchy, nel 1823, ne diede una dimostrazione che presupponeva la continuità della
derivata. La dimostrazione che abbiamo presentato, che deduce il teorema da quello di
Rolle mediante la funzione ausiliaria \eqref{eq:lagrange-g}, è dovuta a Ossian Bonnet e
fu pubblicata nel \emph{Cours de calcul différentiel et intégral} di Joseph Serret
(1868); da allora è quella di tutti i manuali.
\end{notastorica}

%% ----------------------------------------------------------------------------
%%  §6.3  Funzioni crescenti e decrescenti
%% ----------------------------------------------------------------------------
\section{Funzioni crescenti e decrescenti}
\label{sec:monotonia}

Ricordiamo (\cref{sec:invertibili}) che una funzione $f$ si dice
\emph{crescente} in un insieme se $x_{1}<x_{2}$ implica $f(x_{1})\le f(x_{2})$, e
\emph{strettamente crescente} se $x_{1}<x_{2}$ implica $f(x_{1})<f(x_{2})$; analogamente
per le funzioni decrescenti. L'intuizione dice che una funzione con derivata positiva
«sale». Il teorema di Lagrange permette di trasformare questa intuizione in un
teorema, e anzi in una caratterizzazione completa.

\begin{teorema}[Criterio di monotonia]\label{teo:criterio-monotonia}
\index{criterio!di monotonia}\index{funzione!crescente!criterio}
Sia $f$ una funzione continua nell'intervallo $[a,b]$ e derivabile in $(a,b)$. Allora
\begin{equation}\label{eq:criterio-monotonia}
  f \text{ è crescente in } [a,b] \quad\sse\quad f'(x)\ge0 \quad\perogni x\in(a,b);
\end{equation}
analogamente, $f$ è decrescente in $[a,b]$ se e solo se $f'(x)\le0$ per ogni
$x\in(a,b)$.
\end{teorema}

\begin{schema}
\Ipotesi{$f$ continua in $[a,b]$, derivabile in $(a,b)$.}
\Tesi{$f$ crescente in $[a,b]$ \ $\sse$ \ $f'\ge0$ in $(a,b)$.}
\end{schema}

\begin{strategia}
Le due implicazioni si dimostrano con strumenti diversi. L'implicazione «$\Leftarrow$»
va dal locale al globale: dobbiamo confrontare $f(x_{1})$ e $f(x_{2})$ conoscendo solo il
segno della derivata. Il teorema di Lagrange scrive esattamente la differenza
$f(x_{2})-f(x_{1})$ come prodotto della derivata in un punto intermedio per
$x_{2}-x_{1}>0$. L'implicazione «$\Rightarrow$» va dal globale al locale ed è più
semplice: per una funzione crescente il rapporto incrementale è sempre $\ge0$, qualunque
sia il segno dell'incremento, e la disuguaglianza passa al limite.
\end{strategia}

\begin{dimostrazione}
\emph{($\Leftarrow$)} Supponiamo $f'(x)\ge0$ per ogni $x\in(a,b)$ e siano
$x_{1},x_{2}\in[a,b]$ con $x_{1}<x_{2}$. La funzione $f$ è continua in $[x_{1},x_{2}]$ e
derivabile in $(x_{1},x_{2})\subseteq(a,b)$; per il teorema di Lagrange
(\cref{teo:lagrange}) esiste $x_{0}\in(x_{1},x_{2})$ tale che
\begin{equation}\label{eq:monotonia-lagrange}
  f(x_{2})-f(x_{1})=f'(x_{0})\,(x_{2}-x_{1}) .
\end{equation}
Il secondo membro è il prodotto di $f'(x_{0})\ge0$ per $x_{2}-x_{1}>0$, quindi è $\ge0$:
risulta $f(x_{1})\le f(x_{2})$, e $f$ è crescente.

\emph{($\Rightarrow$)} Supponiamo $f$ crescente in $[a,b]$ e sia $x\in(a,b)$. Per ogni
$h\neq0$ tale che $x+h\in[a,b]$, il numeratore e il denominatore del rapporto
incrementale $\frac{f(x+h)-f(x)}{h}$ hanno lo stesso segno (se $h>0$ è $f(x+h)\ge f(x)$,
se $h<0$ è $f(x+h)\le f(x)$); quindi il rapporto è $\ge0$. Per il teorema del confronto
(\cref{teo:confronto-funzioni}\,(ii)) anche il suo limite per $h\to0$, cioè $f'(x)$, è
$\ge0$.

Il caso delle funzioni decrescenti si ottiene applicando quanto provato a $-f$.
\end{dimostrazione}

\begin{osservazione}[Intervalli qualunque]\label{oss:monotonia-intervalli}
Il criterio vale in un intervallo $I$ qualunque (aperto, chiuso, semiaperto, limitato o
no), purché $f$ sia continua in $I$ e derivabile nei punti interni di $I$. Infatti la
monotonia in $I$ è una proprietà che riguarda le coppie di punti $x_{1}<x_{2}$ di $I$,
e per ciascuna coppia basta applicare il teorema all'intervallo $[x_{1},x_{2}]\subseteq I$.
\end{osservazione}

\begin{attenzione}[L'ipotesi che il dominio sia un intervallo è essenziale]
La funzione $f(x)=\frac1x$ ha derivata $f'(x)=-\frac1{x^{2}}<0$ in ogni punto del suo
dominio $\R\setminus\{0\}$, ma \emph{non} è decrescente in tale insieme: ad esempio
$-1<1$ e $f(-1)=-1<1=f(1)$. È invece decrescente in ciascuno dei due intervalli
$(-\infty,0)$ e $(0,+\infty)$ (\cref{fig:x3-iperbole}, a destra). Nella dimostrazione,
il teorema di Lagrange non si può applicare a $[-1,1]$, che contiene il punto $0$ in cui
$f$ non è definita. Dire «$f$ è decrescente perché ha derivata negativa» senza
specificare l'intervallo è un errore frequente.
\end{attenzione}

\begin{esempio}[Monotonia delle funzioni elementari]\label{es:monotonia-elementari}
Le funzioni $e^{x}$, $\log x$, $\operatorname{arctg}x$ hanno derivate $e^{x}$,
$\frac1x$, $\frac1{1+x^{2}}$, positive nei rispettivi domini, che sono intervalli; per il
criterio sono crescenti (anzi strettamente crescenti, come vedremo tra poco). La
funzione $x^{2}$ ha derivata $2x$, negativa per $x<0$ e positiva per $x>0$: è
decrescente in $(-\infty,0]$ e crescente in $[0,+\infty)$. Si noti che questi fatti erano
già noti dai capitoli precedenti, ma ora discendono tutti da un unico calcolo.

La funzione $f(x)=x^{3}-3x$ ha derivata $f'(x)=3x^{2}-3=3(x-1)(x+1)$, positiva per
$\abs x>1$ e negativa per $\abs x<1$. Quindi $f$ è crescente in $(-\infty,-1]$,
decrescente in $[-1,1]$ e di nuovo crescente in $[1,+\infty)$. Ne segue che $x=-1$ è un
punto di massimo relativo, con $f(-1)=2$, e $x=1$ è un punto di minimo relativo, con
$f(1)=-2$ (\cref{fig:cubica-monotonia}). Nessuno dei due è assoluto, perché
$f(x)\to\pm\infty$ per $x\to\pm\infty$.
\end{esempio}

\begin{figure}[htb]
\centering
\begin{tikzpicture}
\begin{axis}[width=0.8\textwidth, height=6.2cm, xmin=-2.5, xmax=2.5, ymin=-3.3, ymax=3.3,
  xtick={-1.732,-1,1,1.732}, xticklabels={}, ytick={-2,2},
  yticklabels={$-2$,{}}, xlabel={$x$}, ylabel={$y$}, clip=false]
  \addplot[curva, domain=-2.1:2.1, samples=120] {x^3-3*x};
  \addplot[tratteggio] coordinates {(-1,0) (-1,2) (0,2)};
  \addplot[tratteggio] coordinates {(1,0) (1,-2) (0,-2)};
  \node[font=\scriptsize, anchor=west] at (axis cs:0.05,2) {$2$};
  \addplot[only marks, mark=*, mark size=1.6pt, rame] coordinates {(-1,2) (1,-2)};
  \node[font=\scriptsize, rame, above] at (axis cs:-1,2.05) {max.\ rel.};
  \node[font=\scriptsize, rame, below] at (axis cs:1,-2.05) {min.\ rel.};
  \node[font=\scriptsize, blu, anchor=east] at (axis cs:1.9,2.75) {$y=x^{3}-3x$};
  \node[font=\scriptsize, anchor=north] at (axis cs:-1,-0.05) {$-1$};
  \node[font=\scriptsize, anchor=south] at (axis cs:1,0.05) {$1$};
  \node[font=\scriptsize, anchor=south east] at (axis cs:-1.78,0.05) {$-\sqrt3$};
  \node[font=\scriptsize, anchor=north west] at (axis cs:1.78,-0.05) {$\sqrt3$};
  % schema dei segni di f'
  \draw[grigio, thin] (axis cs:-2.4,-3.9) -- (axis cs:2.4,-3.9);
  \node[font=\scriptsize, grigio, anchor=east] at (axis cs:-2.45,-3.9) {$f'$};
  \draw[grigio] (axis cs:-1,-3.75) -- (axis cs:-1,-4.05);
  \draw[grigio] (axis cs:1,-3.75) -- (axis cs:1,-4.05);
  \node[font=\scriptsize] at (axis cs:-1,-4.35) {$0$};
  \node[font=\scriptsize] at (axis cs:1,-4.35) {$0$};
  \node[font=\scriptsize, verde!50!black] at (axis cs:-1.75,-3.68) {$+\ \ \nearrow$};
  \node[font=\scriptsize, rosso] at (axis cs:0,-3.68) {$-\ \ \searrow$};
  \node[font=\scriptsize, verde!50!black] at (axis cs:1.75,-3.68) {$+\ \ \nearrow$};
\end{axis}
\end{tikzpicture}
\caption{La funzione $f(x)=x^{3}-3x$ e lo schema del segno della derivata
$f'(x)=3(x-1)(x+1)$: dove $f'>0$ la funzione cresce, dove $f'<0$ decresce. Il passaggio
da crescente a decrescente in $x=-1$ individua un massimo relativo, quello da
decrescente a crescente in $x=1$ un minimo relativo.}
\label{fig:cubica-monotonia}
\end{figure}

Lo studio del segno della derivata si presta bene a essere automatizzato. Il programma
seguente trova i punti in cui $f'$ si annulla e determina il segno di $f'$ in ciascuno
degli intervalli che essi individuano, valutando la derivata in un punto di prova
interno. Il procedimento è corretto perché $f'$ è continua e, per il teorema degli zeri
(\cref{teo:zeri}), non può cambiare segno all'interno di un intervallo in cui non si
annulla.

\begin{python}[Gli intervalli di monotonia dal segno della derivata]
import sympy as sp

x = sp.symbols("x", real=True)

def punto_interno(p, q):
    """Un punto di prova nell'intervallo aperto (p, q)."""
    if p.is_finite and q.is_finite:
        return (p + q) / 2
    if p.is_finite:
        return p + 1
    return q - 1 if q.is_finite else sp.Integer(0)

def tabella_monotonia(f):
    """Intervalli di monotonia di f su R, dal segno di f'."""
    df = sp.factor(sp.diff(f, x))
    critici = sorted(sp.solveset(df, x, sp.S.Reals))
    print(f"f = {f},   f' = {df},   f' = 0 in {critici}")
    nodi = [-sp.oo] + critici + [sp.oo]
    for p, q in zip(nodi[:-1], nodi[1:]):
        s = sp.sign(df.subs(x, punto_interno(p, q)))
        verso = "crescente" if s > 0 else "decrescente"
        print(f"  ({p}, {q}):  f' {'+' if s > 0 else '-'}  {verso}")

tabella_monotonia(x**3 - 3*x)
tabella_monotonia(x**4 - 2*x**2)
\end{python}
\begin{risultato}
f = x**3 - 3*x,   f' = 3*(x - 1)*(x + 1),   f' = 0 in [-1, 1]
  (-oo, -1):  f' +  crescente
  (-1, 1):  f' -  decrescente
  (1, oo):  f' +  crescente
f = x**4 - 2*x**2,   f' = 4*x*(x - 1)*(x + 1),   f' = 0 in [-1, 0, 1]
  (-oo, -1):  f' -  decrescente
  (-1, 0):  f' +  crescente
  (0, 1):  f' -  decrescente
  (1, oo):  f' +  crescente
\end{risultato}

\subsection{Funzioni con derivata nulla}

Un caso particolare del criterio di monotonia è di importanza fondamentale.

\begin{corollario}[Caratterizzazione delle funzioni costanti]\label{cor:funzioni-costanti}
\index{funzione!costante!caratterizzazione}
Sia $f$ continua in $[a,b]$ e derivabile in $(a,b)$. Allora $f$ è costante in $[a,b]$
se e solo se $f'(x)=0$ per ogni $x\in(a,b)$.
\end{corollario}

\begin{schema}
\Ipotesi{$f$ continua in $[a,b]$, derivabile in $(a,b)$.}
\Tesi{$f$ costante in $[a,b]$ \ $\sse$ \ $f'=0$ in $(a,b)$.}
\end{schema}

\begin{strategia}
Una funzione è costante se e solo se è contemporaneamente crescente e decrescente; e la
derivata è nulla se e solo se è contemporaneamente $\ge0$ e $\le0$. Basta applicare due
volte il criterio di monotonia.
\end{strategia}

\begin{dimostrazione}
Se $f$ è costante, la sua derivata è nulla \eqref{eq:derivata-costante}. Viceversa, se
$f'(x)=0$ in $(a,b)$, risulta sia $f'\ge0$ sia $f'\le0$; per il
\cref{teo:criterio-monotonia} $f$ è sia crescente sia decrescente in $[a,b]$. Allora,
per ogni $x\in[a,b]$, da $a\le x$ segue $f(a)\le f(x)$ e $f(a)\ge f(x)$, cioè
$f(x)=f(a)$.
\end{dimostrazione}

Anche qui il dominio deve essere un intervallo: la funzione che vale $1$ per $x>0$ e
$-1$ per $x<0$ ha derivata nulla in tutto il suo dominio $\R\setminus\{0\}$, ma non è
costante. Un esempio meno artificioso è offerto dalla funzione
$\operatorname{arctg}x+\operatorname{arctg}\frac1x$, la cui derivata è nulla per ogni
$x\neq0$ (\cref{es:arctg-1-su-x}): la funzione vale $\frac\pi2$ in $(0,+\infty)$ e
$-\frac\pi2$ in $(-\infty,0)$. Analogamente, l'identità
$\arcsen x+\arccos x=\frac\pi2$, dimostrata per via trigonometrica
nell'\cref{oss:arcsen-arccos}, si ritrova ora senza calcoli: la funzione
$\arcsen x+\arccos x$ è continua in $[-1,1]$ e ha derivata
$\frac1{\sqrt{1-x^{2}}}-\frac1{\sqrt{1-x^{2}}}=0$ in $(-1,1)$; quindi è costante, e il
suo valore è quello assunto in $x=0$, cioè $0+\frac\pi2$.

\begin{corollario}[Funzioni con la stessa derivata]\label{cor:stessa-derivata}
Se $f$ e $g$ sono continue in un intervallo $I$, derivabili nei punti interni e
$f'(x)=g'(x)$ in tali punti, allora esiste una costante $c$ tale che $f(x)=g(x)+c$ per
ogni $x\in I$.
\end{corollario}

\begin{schema}
\Ipotesi{$f,g$ continue nell'intervallo $I$, derivabili nei punti interni, con $f'=g'$.}
\Tesi{$f-g$ è costante in $I$.}
\end{schema}

\begin{strategia}
Si applica la caratterizzazione delle funzioni costanti alla differenza $f-g$, la cui
derivata è $f'-g'=0$.
\end{strategia}

\begin{dimostrazione}
La funzione $f-g$ ha derivata nulla nei punti interni di $I$, ed è quindi costante per
il \cref{cor:funzioni-costanti} (applicato a ogni intervallo chiuso contenuto in $I$).
\end{dimostrazione}

Questo corollario, apparentemente modesto, è alla base del calcolo integrale: le
funzioni che hanno una derivata assegnata (le \emph{primitive}) differiscono tra loro
per una costante. Ne vediamo subito un'applicazione a un problema fisico.

\begin{esempio}[La pressione nella stratosfera]\label{es:atmosfera-isoterma}
Nell'atmosfera standard internazionale (ISA), tra le quote $z_{1}=11\,000$~m e
$20\,000$~m la temperatura è costante, $T=216{,}65$~K. L'equilibrio idrostatico di
uno strato d'aria impone $p'(z)=-\rho(z)g$, e l'equazione di stato dei gas perfetti dà
$\rho=\frac{p}{RT}$; quindi la pressione verifica
\begin{equation}\label{eq:isoterma}
  p'(z)=-\frac{p(z)}{H}, \qquad H=\frac{RT}{g},
\end{equation}
con la condizione $p(z_{1})=p_{1}=22\,632$~Pa. Una soluzione è
$p(z)=p_{1}e^{-(z-z_{1})/H}$, come si verifica derivando. \emph{È l'unica?} Sia $p$ una
qualunque funzione derivabile che verifica la \eqref{eq:isoterma}, e consideriamo
$\varphi(z)=p(z)\,e^{(z-z_{1})/H}$. Per la regola del prodotto e per la
\eqref{eq:isoterma},
\begin{align*}
  \varphi'(z)&=p'(z)\,e^{(z-z_{1})/H}+p(z)\,\frac{e^{(z-z_{1})/H}}{H}\\
  &=\Bigl(p'(z)+\frac{p(z)}H\Bigr)e^{(z-z_{1})/H}=0 .
\end{align*}
Per il \cref{cor:funzioni-costanti} $\varphi$ è costante, e quindi
$\varphi(z)=\varphi(z_{1})=p_{1}$, cioè $p(z)=p_{1}e^{-(z-z_{1})/H}$: la soluzione è
unica. Lo stesso ragionamento mostra che \emph{l'unica funzione derivabile su $\R$ tale
che $f'=f$ e $f(0)=1$ è $e^{x}$}. La costante $H$, detta \emph{altezza di scala}, è la
quota di cui bisogna salire perché la pressione si riduca di un fattore $e$.
\end{esempio}

\begin{python}[La pressione nello strato isotermo dell'atmosfera ISA]
import math

# ISA, strato isotermo tra 11 e 20 km
R, g, T = 287.05, 9.80665, 216.65        # J/(kg K), m/s^2, K
z1, p1 = 11000.0, 22632.0                # m, Pa (tropopausa)
H = R * T / g                            # altezza di scala
p = lambda z: p1 * math.exp(-(z - z1) / H)
print(f"H = {H:.1f} m")
for z in (11000, 14000, 17000, 20000):
    rho = p(z) / (R * T)
    print(f"z = {z:5d} m   p = {p(z):8.1f} Pa"
          f"   rho = {rho:.4f} kg/m^3")
\end{python}
\begin{risultato}
H = 6341.6 m
z = 11000 m   p =  22632.0 Pa   rho = 0.3639 kg/m^3
z = 14000 m   p =  14101.7 Pa   rho = 0.2268 kg/m^3
z = 17000 m   p =   8786.6 Pa   rho = 0.1413 kg/m^3
z = 20000 m   p =   5474.8 Pa   rho = 0.0880 kg/m^3
\end{risultato}

I valori coincidono con quelli tabulati dall'ISA (a $20$~km, $p\approx5475$~Pa e
$\rho\approx0{,}0880$~kg/m$^{3}$): le tabelle dell'atmosfera standard, in questo strato,
non sono altro che la soluzione unica della \eqref{eq:isoterma}.

\subsection{Monotonia stretta}

Una funzione con derivata positiva è strettamente crescente; ma vale anche il
viceversa? La funzione $x^{3}$ è strettamente crescente su $\R$, perché
$x_{1}<x_{2}\sse x_{1}^{3}<x_{2}^{3}$, eppure la sua derivata $3x^{2}$ si annulla in
$x=0$ (\cref{fig:x3-iperbole}, a sinistra). La derivata di una funzione strettamente
crescente può dunque annullarsi in qualche punto; ciò che non può fare è annullarsi in
tutto un intervallo.

\begin{teorema}[Criterio di stretta monotonia]\label{teo:stretta-monotonia}
\index{criterio!di stretta monotonia}
Sia $f$ continua in $[a,b]$ e derivabile in $(a,b)$. Allora $f$ è strettamente
crescente in $[a,b]$ se e solo se
\begin{enumerate}[label=\textup{(\roman*)}]
  \item $f'(x)\ge0$ per ogni $x\in(a,b)$, e
  \item $f'$ non è identicamente nulla in alcun intervallo $[x_{1},x_{2}]\subseteq(a,b)$
  con $x_{1}<x_{2}$.
\end{enumerate}
\end{teorema}

\begin{schema}
\Ipotesi{$f$ continua in $[a,b]$, derivabile in $(a,b)$.}
\Tesi{$f$ strettamente crescente \ $\sse$ \ $f'\ge0$ in $(a,b)$ e $f'$ non si annulla
identicamente in alcun sottointervallo.}
\end{schema}

\begin{strategia}
Si confronta la monotonia stretta con quella debole: una funzione crescente che non è
strettamente crescente assume lo stesso valore in due punti $x_{1}<x_{2}$, e quindi,
essendo crescente, è costante in tutto $[x_{1},x_{2}]$. Le funzioni costanti su un
intervallo sono esattamente quelle con derivata identicamente nulla
(\cref{cor:funzioni-costanti}). Tutto si riduce quindi al criterio di monotonia e alla
caratterizzazione delle funzioni costanti.
\end{strategia}

\begin{dimostrazione}
\emph{($\Leftarrow$)} Per la (i) e il \cref{teo:criterio-monotonia}, $f$ è crescente in
$[a,b]$. Se non fosse strettamente crescente, esisterebbero $x_{1}<x_{2}$ in $[a,b]$
con $f(x_{1})=f(x_{2})$; per ogni $x\in[x_{1},x_{2}]$ sarebbe allora
$f(x_{1})\le f(x)\le f(x_{2})=f(x_{1})$, cioè $f$ sarebbe costante in $[x_{1},x_{2}]$ e
$f'$ sarebbe nulla in $(x_{1},x_{2})$, quindi in ogni intervallo chiuso contenuto in
$(x_{1},x_{2})$, contro la (ii).

\emph{($\Rightarrow$)} Se $f$ è strettamente crescente, è in particolare crescente, e
quindi vale la (i) per il criterio di monotonia. Se $f'$ fosse identicamente nulla in
un intervallo $[x_{1},x_{2}]\subseteq(a,b)$ con $x_{1}<x_{2}$, per il
\cref{cor:funzioni-costanti} $f$ sarebbe costante in $[x_{1},x_{2}]$, contro l'ipotesi
di stretta monotonia.
\end{dimostrazione}

\begin{corollario}\label{cor:derivata-positiva}
Se $f$ è continua in un intervallo $I$, derivabile nei punti interni e $f'(x)>0$ in tali
punti, allora $f$ è strettamente crescente in $I$. Lo stesso vale se $f'\ge0$ e, in ogni
intervallo limitato contenuto in $I$, $f'$ si annulla soltanto in un numero finito di
punti.
\end{corollario}

\begin{schema}
\Ipotesi{$f$ continua nell'intervallo $I$, derivabile nei punti interni; $f'>0$, oppure
$f'\ge0$ con un numero finito di zeri in ogni intervallo limitato.}
\Tesi{$f$ strettamente crescente in $I$.}
\end{schema}

\begin{strategia}
Basta verificare, in ogni intervallo $[x_{1},x_{2}]\subseteq I$, le due condizioni del
criterio di stretta monotonia; la seconda vale perché un intervallo non degenere
contiene infiniti punti, e quindi non può essere formato da zeri di $f'$.
\end{strategia}

\begin{dimostrazione}
Siano $x_{1}<x_{2}$ in $I$. In $[x_{1},x_{2}]$ la funzione $f$ è continua e derivabile
nei punti interni, con $f'\ge0$; inoltre $f'$ non è identicamente nulla in alcun
intervallo $[s,t]\subseteq(x_{1},x_{2})$ con $s<t$, perché $[s,t]$ contiene infiniti punti
mentre gli zeri di $f'$ in $[x_{1},x_{2}]$ sono in numero finito (o non esistono). Per il
\cref{teo:stretta-monotonia}, $f$ è strettamente crescente in $[x_{1},x_{2}]$, e in
particolare $f(x_{1})<f(x_{2})$.
\end{dimostrazione}

Ad esempio, la funzione $x-\sen x$ ha derivata $1-\cos x\ge0$, che si annulla soltanto
nei punti $2k\pi$, in numero finito in ogni intervallo limitato: quindi $x-\sen x$ è
strettamente crescente su $\R$.

\begin{figure}[htb]
\centering
\begin{tikzpicture}
\begin{axis}[name=m1, width=0.48\textwidth, height=0.42\textwidth, xmin=-1.5, xmax=1.5,
  ymin=-1.6, ymax=1.6, xtick={-1,1}, ytick={-1,1}, xlabel={$x$}, ylabel={$y$},
  title style={align=center, font=\scriptsize, yshift=8pt},
  title={$x^{3}$: strettamente crescente,\\ con $f'(0)=0$}]
  \addplot[curva, domain=-1.17:1.17, samples=100] {x^3};
  \addplot[curva2, domain=-0.8:0.8, samples=2] {0};
  \addplot[only marks, mark=*, mark size=1.4pt, rame] coordinates {(0,0)};
\end{axis}
\begin{axis}[at={(m1.east)}, anchor=west, xshift=1cm, width=0.48\textwidth,
  height=0.42\textwidth, xmin=-2.6, xmax=2.6, ymin=-2.6, ymax=2.6, xtick={-1,1},
  xticklabels={{},$1$}, ytick={-1,1}, yticklabels={{},$1$}, xlabel={$x$}, ylabel={$y$},
  title style={align=center, font=\scriptsize, yshift=8pt},
  title={$\frac1x$: $f'<0$, ma non decrescente\\ in $\R\setminus\{0\}$}]
  \addplot[curva, domain=0.39:2.5, samples=100] {1/x};
  \addplot[curva, domain=-2.5:-0.39, samples=100] {1/x};
  \addplot[tratteggio] coordinates {(-1,0) (-1,-1) (0,-1)};
  \addplot[tratteggio] coordinates {(1,0) (1,1) (0,1)};
  \addplot[only marks, mark=*, mark size=1.4pt, rame] coordinates {(-1,-1) (1,1)};
  \node[font=\scriptsize, anchor=south] at (axis cs:-1,0.08) {$-1$};
  \node[font=\scriptsize, anchor=west] at (axis cs:0.08,-1) {$-1$};
\end{axis}
\end{tikzpicture}
\caption{A sinistra: la derivata di una funzione strettamente crescente può annullarsi
in un punto isolato (la tangente in $0$ è orizzontale). A destra: una derivata negativa
in tutto il dominio non garantisce la monotonia se il dominio non è un intervallo, come
mostrano i punti $(-1,-1)$ e $(1,1)$.}
\label{fig:x3-iperbole}
\end{figure}

\subsection{Disuguaglianze dimostrate con le derivate}

Il criterio di monotonia è uno strumento molto efficace per dimostrare disuguaglianze.
Lo schema è sempre lo stesso: per provare che $u(x)\ge v(x)$ in un intervallo, si studia
il segno della derivata della differenza $g=u-v$ e se ne deduce che il minimo di $g$ è
$\ge0$.

\begin{esempio}[Due disuguaglianze fondamentali]\label{es:disuguaglianze-exp-log}
Dimostriamo che
\begin{equation}\label{eq:exp-log-disuguaglianze}
  e^{x}\ge1+x \quad\perogni x\in\R, \qquad\qquad \log(1+x)\le x \quad\perogni x>-1,
\end{equation}
con l'uguaglianza soltanto per $x=0$. Sia $g(x)=e^{x}-1-x$; allora $g'(x)=e^{x}-1$, che
è negativa per $x<0$ e positiva per $x>0$ (perché l'esponenziale è strettamente
crescente ed $e^{0}=1$). Per il \cref{cor:derivata-positiva}, $g$ è strettamente
decrescente in $(-\infty,0]$ e strettamente crescente in $[0,+\infty)$; quindi $g(x)>g(0)=0$
per ogni $x\neq0$. Per la seconda, sia $h(x)=x-\log(1+x)$ per $x>-1$: si ha
$h'(x)=1-\frac1{1+x}=\frac{x}{1+x}$, negativa in $(-1,0)$ e positiva in $(0,+\infty)$, e
come prima $h(x)>h(0)=0$ per $x\neq0$. Geometricamente: il grafico di $e^{x}$ sta sopra
la sua tangente $y=1+x$ nel punto $0$, e quello di $\log(1+x)$ sta sotto la sua tangente
$y=x$. Nel prossimo paragrafo vedremo che questo è un fatto generale per le funzioni
convesse e concave.
\end{esempio}

\begin{attenzione}[Attenzione ai circoli viziosi]
Con lo stesso metodo si può «dimostrare» che $\sen x<x$ per $x>0$: la funzione
$x-\sen x$ è strettamente crescente e vale $0$ in $0$. Questo ragionamento è corretto
come \emph{verifica}, ma non può sostituire la dimostrazione geometrica delle
disuguaglianze $\sen x<x<\tg x$ data nel \cref{sec:trigonometriche} \eqref{eq:sen-x-tg}:
la derivata del seno, infatti, è stata calcolata nel \cref{sec:derivate-elementari}
proprio a partire dal limite notevole $\sen x/x\to1$, che a sua volta poggia su quelle
disuguaglianze. Usare la derivata del seno per dimostrarle sarebbe un circolo vizioso:
prima di usare un risultato, bisogna sempre chiedersi su che cosa poggia.
\end{attenzione}

\subsection{Il test della derivata prima}

Lo schema dei segni della \cref{fig:cubica-monotonia} suggerisce un criterio per
riconoscere i punti di estremo, che useremo sistematicamente nello studio di funzioni.

\begin{corollario}[Test della derivata prima]\label{cor:test-derivata-prima}
\index{test della derivata prima}
Sia $f$ continua in un intervallo $(x_{0}-\delta,x_{0}+\delta)$ e derivabile in tale
intervallo, eccetto al più nel punto $x_{0}$. Se
\begin{equation}\label{eq:test-derivata-prima}
  f'(x)\ge0 \ \text{ per } x_{0}-\delta<x<x_{0}, \qquad
  f'(x)\le0 \ \text{ per } x_{0}<x<x_{0}+\delta,
\end{equation}
allora $x_{0}$ è un punto di massimo relativo per $f$; se le disuguaglianze sono strette,
il massimo relativo è stretto. Con le disuguaglianze invertite, $x_{0}$ è un punto di
minimo relativo.
\end{corollario}

\begin{schema}
\Ipotesi{$f$ continua in $(x_{0}-\delta,x_{0}+\delta)$, derivabile per $x\neq x_{0}$;
$f'\ge0$ a sinistra di $x_{0}$, $f'\le0$ a destra.}
\Tesi{$x_{0}$ è di massimo relativo (stretto, se le disuguaglianze sono strette).}
\end{schema}

\begin{strategia}
Il segno della derivata dà la monotonia separatamente a sinistra e a destra di $x_{0}$:
la funzione sale fino a $x_{0}$ e poi scende, quindi $f(x_{0})$ è il valore più alto
nell'intorno. La derivabilità in $x_{0}$ non serve, perché il criterio di monotonia
richiede la derivabilità solo nei punti interni degli intervalli.
\end{strategia}

\begin{dimostrazione}
Per il criterio di monotonia, applicato agli intervalli $[x_{0}-\delta',x_{0}]$ e
$[x_{0},x_{0}+\delta']$ per ogni $0<\delta'<\delta$, $f$ è crescente in
$(x_{0}-\delta,x_{0}]$ e decrescente in $[x_{0},x_{0}+\delta)$; quindi $f(x)\le f(x_{0})$
sia per $x<x_{0}$ sia per $x>x_{0}$. Se le disuguaglianze sono strette, la monotonia è
stretta (\cref{cor:derivata-positiva}) e $f(x)<f(x_{0})$ per $x\neq x_{0}$.
\end{dimostrazione}

Il test non richiede la derivabilità in $x_{0}$: per $f(x)=\abs x$ si ha $f'(x)=-1$ per
$x<0$ e $f'(x)=1$ per $x>0$, e il test riconosce correttamente il minimo in $0$, che il
teorema di Fermat non poteva vedere. Applicato all'efficienza
$E(C_{L})=\frac{C_{L}}{C_{D0}+kC_{L}^{2}}$, la cui derivata è positiva per
$C_{L}<C_{L}^{*}$ e negativa per $C_{L}>C_{L}^{*}$ (\cref{sec:operazioni-derivate}), il
\cref{cor:derivata-positiva} mostra che $E$ è strettamente crescente in $[0,C_{L}^{*}]$ e
strettamente decrescente in $[C_{L}^{*},+\infty)$: il massimo in $C_{L}^{*}$ è assoluto.

%% ----------------------------------------------------------------------------
%%  §6.4  Funzioni convesse e concave
%% ----------------------------------------------------------------------------
\section{Funzioni convesse e concave}
\label{sec:convessita}

Il segno della derivata prima dice se il grafico sale o scende; non dice però
\emph{come} sale o scende. Le funzioni $x^{2}$ e $\sqrt x$ sono entrambe crescenti in
$[0,+\infty)$, ma la prima cresce sempre più rapidamente e il suo grafico volge la
concavità verso l'alto, la seconda cresce sempre più lentamente e il suo grafico volge
la concavità verso il basso. Per descrivere questa proprietà introduciamo una nuova
definizione.

Si dice che una funzione è \emph{convessa} in un intervallo se, per ogni punto
$x_{0}$ dell'intervallo, il grafico della funzione sta \emph{al di sopra} della retta
tangente nel punto $(x_{0},f(x_{0}))$; si dice \emph{concava} se il grafico sta \emph{al di
sotto} di ogni retta tangente. Usando l'equazione della retta tangente
\eqref{eq:retta-tangente}, possiamo dare la definizione in forma precisa.

\begin{definizione}[Funzioni convesse e concave]\label{def:convessa}
\index{funzione!convessa}\index{funzione!concava}
Sia $f$ una funzione derivabile nell'intervallo $[a,b]$ (con derivate unilaterali negli
estremi). Si dice che $f$ è \emph{convessa} in $[a,b]$ se
\begin{equation}\label{eq:convessa}
  f(x)\ge f(x_{0})+f'(x_{0})(x-x_{0}) \qquad\perogni x,x_{0}\in[a,b];
\end{equation}
si dice che $f$ è \emph{concava} in $[a,b]$ se
\begin{equation}\label{eq:concava}
  f(x)\le f(x_{0})+f'(x_{0})(x-x_{0}) \qquad\perogni x,x_{0}\in[a,b].
\end{equation}
Se le disuguaglianze sono strette per $x\neq x_{0}$, la funzione si dice
\emph{strettamente} convessa o concava. Una funzione è convessa (concava) in un
intervallo qualunque $I$ se lo è in ogni intervallo chiuso e limitato contenuto in $I$.
\end{definizione}

Evidentemente $f$ è concava se e solo se $-f$ è convessa: per questo basta
studiare le funzioni convesse. Nella \cref{fig:convessa-concava} la funzione è convessa
nell'intervallo $[a,c]$ e concava nell'intervallo $[c,b]$: in $c$ la concavità cambia
verso.

\begin{definizione}[Punto di flesso]\label{def:flesso}\index{punto!di flesso}\index{flesso}
Sia $f$ derivabile in un intorno di $x_{0}$. Si dice che $x_{0}$ è un \emph{punto di
flesso} per $f$ se esiste $\delta>0$ tale che $f$ è convessa in $[x_{0}-\delta,x_{0}]$ e
concava in $[x_{0},x_{0}+\delta]$, oppure concava in $[x_{0}-\delta,x_{0}]$ e convessa in
$[x_{0},x_{0}+\delta]$.
\end{definizione}

In un punto di flesso la retta tangente \emph{attraversa} il grafico: da una parte il
grafico sta sopra la tangente, dall'altra sotto. Se la tangente è orizzontale si parla
di \emph{flesso a tangente orizzontale} (come per $x^{3}$ in $0$), altrimenti di flesso a
tangente obliqua.

\begin{figure}[htb]
\centering
\begin{tikzpicture}
\begin{axis}[width=0.9\textwidth, height=5.6cm, xmin=0.4, xmax=5.6, ymin=1.15, ymax=2.85,
  axis lines=left, xtick={0.8,3,5.2}, xticklabels={$a$,$c$,$b$}, ytick=\empty,
  clip=false]
  \addplot[curva, domain=0.8:5.2, samples=120] {2-0.15*(x-3)^3+0.5*(x-3)};
  \addplot[curva2, domain=0.55:2.0] {1.9748-0.958*(x-1.2)};
  \addplot[curva2, domain=2.2:3.8] {2+0.5*(x-3)};
  \addplot[curva2, domain=4.0:5.45] {2.0252-0.958*(x-4.8)};
  \addplot[tratteggio] coordinates {(3,1.15) (3,2)};
  \addplot[only marks, mark=*, mark size=1.6pt, blu] coordinates
    {(1.2,1.9748) (3,2) (4.8,2.0252)};
  \node[font=\scriptsize, grigio] at (axis cs:1.9,2.65) {convessa};
  \node[font=\scriptsize, grigio] at (axis cs:3.6,2.7) {concava};
  \node[font=\scriptsize, rame, anchor=north west] at (axis cs:3.05,1.95) {flesso};
\end{axis}
\end{tikzpicture}
\caption{Una funzione convessa in $[a,c]$, dove il grafico sta sopra le tangenti, e
concava in $[c,b]$, dove sta sotto. Nel punto di flesso $c$ la tangente attraversa il
grafico.}
\label{fig:convessa-concava}
\end{figure}

Come per la monotonia, esiste un criterio che riconduce la convessità al segno di una
derivata, questa volta della derivata seconda.

\needspace{12\baselineskip}
\begin{teorema}[Criterio di convessità]\label{teo:criterio-convessita}
\index{criterio!di convessità}
Sia $f$ una funzione derivabile in $[a,b]$ che ammette derivata seconda in $(a,b)$. Le
seguenti condizioni sono fra loro equivalenti:
\begin{enumerate}[label=\textup{(\alph*)}]
  \item $f$ è convessa in $[a,b]$;
  \item $f'$ è crescente in $[a,b]$;
  \item $f''(x)\ge0$ per ogni $x\in(a,b)$.
\end{enumerate}
Analogamente, $f$ è concava in $[a,b]$ se e solo se $f'$ è decrescente in $[a,b]$, se e
solo se $f''\le0$ in $(a,b)$.
\end{teorema}

\begin{schema}
\Ipotesi{$f$ derivabile in $[a,b]$, derivabile due volte in $(a,b)$.}
\Tesi{(a) $f$ convessa \ $\sse$ \ (b) $f'$ crescente \ $\sse$ \ (c) $f''\ge0$ in $(a,b)$.}
\end{schema}

\begin{strategia}
Dimostriamo le equivalenze (a)$\sse$(b) e (b)$\sse$(c).
\begin{itemize}[nosep]
  \item (a)$\Rightarrow$(b): scriviamo la disuguaglianza di convessità con il punto di
  tangenza in $x_{1}$ e poi in $x_{2}$, e sommiamo: i valori di $f$ si cancellano e resta
  una disuguaglianza tra le derivate.
  \item (b)$\Rightarrow$(a): per il teorema di Lagrange la differenza $f(x)-f(x_{0})$ è
  uguale a $f'(x_{1})(x-x_{0})$, con $x_{1}$ tra $x_{0}$ e $x$; la monotonia di $f'$
  permette di confrontare $f'(x_{1})$ con $f'(x_{0})$.
  \item (b)$\sse$(c): è il criterio di monotonia applicato alla funzione $f'$, con una
  piccola cautela negli estremi $a$ e $b$, dove non sappiamo se $f'$ sia continua.
\end{itemize}
\end{strategia}

\begin{dimostrazione}
\emph{(a)$\Rightarrow$(b).} Siano $x_{1},x_{2}\in[a,b]$ con $x_{1}<x_{2}$. Scrivendo la
\eqref{eq:convessa} con $x_{0}=x_{1}$, $x=x_{2}$, e poi con $x_{0}=x_{2}$, $x=x_{1}$,
otteniamo
\[
  f(x_{2})\ge f(x_{1})+f'(x_{1})(x_{2}-x_{1}), \qquad
  f(x_{1})\ge f(x_{2})+f'(x_{2})(x_{1}-x_{2}) .
\]
Sommando membro a membro e semplificando,
\begin{equation}\label{eq:convessa-somma}
  0\ge f'(x_{1})(x_{2}-x_{1})+f'(x_{2})(x_{1}-x_{2})=-\bigl[f'(x_{2})-f'(x_{1})\bigr](x_{2}-x_{1}),
\end{equation}
cioè $\bigl[f'(x_{2})-f'(x_{1})\bigr](x_{2}-x_{1})\ge0$. Essendo $x_{2}-x_{1}>0$, ne segue
$f'(x_{1})\le f'(x_{2})$: $f'$ è crescente.

\emph{(b)$\Rightarrow$(a).} Siano $x,x_{0}\in[a,b]$ con $x\neq x_{0}$. Per il teorema di
Lagrange (\cref{teo:lagrange}) esiste $x_{1}$ strettamente compreso tra $x_{0}$ e $x$
tale che
\begin{equation}\label{eq:convessa-lagrange}
  f(x)-f(x_{0})=f'(x_{1})(x-x_{0}) .
\end{equation}
Se $x>x_{0}$, allora $x_{1}>x_{0}$ e, per la monotonia di $f'$, $f'(x_{1})\ge f'(x_{0})$;
moltiplicando per $x-x_{0}>0$ si ottiene $f'(x_{1})(x-x_{0})\ge f'(x_{0})(x-x_{0})$. Se
$x<x_{0}$, allora $x_{1}<x_{0}$ e $f'(x_{1})\le f'(x_{0})$; moltiplicando per
$x-x_{0}<0$ il verso si inverte, e si ottiene di nuovo
$f'(x_{1})(x-x_{0})\ge f'(x_{0})(x-x_{0})$. In entrambi i casi, per la
\eqref{eq:convessa-lagrange}, $f(x)-f(x_{0})\ge f'(x_{0})(x-x_{0})$, che è la
\eqref{eq:convessa}. Per $x=x_{0}$ la disuguaglianza è un'uguaglianza.

\emph{(b)$\Rightarrow$(c).} Se $f'$ è crescente, per ogni $x\in(a,b)$ i rapporti
incrementali $\frac{f'(x+h)-f'(x)}{h}$ sono $\ge0$ (come nella dimostrazione del
\cref{teo:criterio-monotonia}), e il loro limite $f''(x)$ è $\ge0$.

\emph{(c)$\Rightarrow$(b).} Siano $a<c<d<b$. Nell'intervallo $[c,d]$ la funzione $f'$ è
derivabile, quindi continua, e ha derivata $f''\ge0$; per il criterio di monotonia
(\cref{teo:criterio-monotonia}) $f'$ è crescente in $[c,d]$. Poiché due punti qualsiasi
di $(a,b)$ sono contenuti in un tale intervallo, $f'$ è crescente in $(a,b)$. Resta da
includere gli estremi, nei quali $f'$ potrebbe non essere continua. Sia $x\in(a,b)$; per
ogni $h$ con $0<h<x-a$ il teorema di Lagrange in $[a,a+h]$ fornisce un punto
$\xi_{h}\in(a,a+h)$ tale che
\[
  \frac{f(a+h)-f(a)}{h}=f'(\xi_{h})\le f'(x),
\]
dove la disuguaglianza vale perché $\xi_{h}<a+h<x$ e $f'$ è crescente in $(a,b)$. Per
$h\to0^{+}$ il primo membro tende a $f'(a)$, e per il teorema del confronto
$f'(a)\le f'(x)$. Allo stesso modo si prova che $f'(x)\le f'(b)$; infine
$f'(a)\le f'(b)$ si ottiene confrontando entrambi con un qualunque $x\in(a,b)$. Dunque
$f'$ è crescente in $[a,b]$.

Il caso delle funzioni concave si ottiene applicando quanto provato a $-f$.
\end{dimostrazione}

\begin{osservazione}[Convessità stretta]\label{oss:convessita-stretta}
Se $f''(x)>0$ per ogni $x\in(a,b)$, allora $f$ è \emph{strettamente} convessa in $[a,b]$.
Infatti, per il \cref{cor:derivata-positiva} applicato a $f'$, la derivata prima è
strettamente crescente in ogni $[c,d]\subset(a,b)$, e quindi in $(a,b)$; negli estremi,
per $a<t<x$ si ha $f'(a)\le f'(t)<f'(x)$. Nella dimostrazione di (b)$\Rightarrow$(a) le
disuguaglianze tra $f'(x_{1})$ e $f'(x_{0})$ diventano allora strette, e così la
\eqref{eq:convessa} per $x\neq x_{0}$.
\end{osservazione}

Negli enunciati che seguono useremo le derivate successive calcolate in un singolo
punto; precisiamo il significato dell'espressione.

\begin{notazione}[Derivate successive in un punto]\label{not:derivabile-n-volte}
\index{derivata!successiva in un punto}
Si dice che $f$ è \emph{derivabile due volte nel punto $x_{0}$} se $f'$ esiste in tutto un
intorno di $x_{0}$ ed è derivabile in $x_{0}$; la derivata seconda $f''(x_{0})$ è allora la
derivata di $f'$ nel punto $x_{0}$. Più in generale, $f$ è \emph{derivabile $n$ volte in
$x_{0}$} se la derivata $f^{(n-1)}$ esiste in tutto un intorno di $x_{0}$ ed è derivabile in
$x_{0}$. In tal caso $f,f',\dots,f^{(n-2)}$ sono derivabili, e quindi continue, in tale
intorno; nulla si richiede invece a $f^{(n-1)}$ fuori del punto $x_{0}$.
\end{notazione}

\begin{corollario}[Condizione necessaria per i flessi]\label{cor:flesso-necessaria}
Se $x_{0}$ è un punto di flesso per $f$ e $f$ è derivabile due volte in $x_{0}$, allora
$f''(x_{0})=0$.
\end{corollario}

\begin{schema}
\Ipotesi{$x_{0}$ punto di flesso; $f$ derivabile due volte in $x_{0}$.}
\Tesi{$f''(x_{0})=0$.}
\end{schema}

\begin{strategia}
Il cambiamento di concavità si traduce, per il criterio di convessità, in un cambiamento
di monotonia della derivata prima: $f'$ cresce da una parte di $x_{0}$ e decresce
dall'altra. Dunque $x_{0}$ è un punto di estremo per la funzione $f'$, e il teorema di
Fermat, applicato a $f'$, dà $(f')'(x_{0})=0$.
\end{strategia}

\begin{dimostrazione}
Supponiamo $f$ convessa in $[x_{0}-\delta,x_{0}]$ e concava in $[x_{0},x_{0}+\delta]$.
Per il criterio di convessità, nella forma (a)$\Rightarrow$(b) (che non usa la derivata
seconda), $f'$ è crescente in $[x_{0}-\delta,x_{0}]$ e decrescente in
$[x_{0},x_{0}+\delta]$. Dunque $x_{0}$ è un punto di massimo relativo \emph{interno} per
la funzione $f'$, che è derivabile in $x_{0}$: per il teorema di Fermat, applicato a
$f'$, risulta $f''(x_{0})=0$. Il caso opposto è analogo (con un minimo di $f'$).
\end{dimostrazione}

La condizione non è sufficiente: la funzione $f(x)=x^{4}$ ha $f''(0)=0$, ma è convessa su
tutto $\R$, perché $f''(x)=12x^{2}\ge0$; il punto $0$ non è di flesso.

\begin{esempio}[Convessità delle funzioni elementari]\label{es:convessita-elementari}
La funzione $e^{x}$ è convessa su tutto $\R$, perché la sua derivata seconda $e^{x}$ è
positiva; la disuguaglianza $e^{x}\ge1+x$ dell'\cref{es:disuguaglianze-exp-log} è la
\eqref{eq:convessa} con $x_{0}=0$. La funzione $\log x$ è concava in $(0,+\infty)$,
perché $D^{2}\log x=-\frac1{x^{2}}<0$. Per $\operatorname{arctg}x$ si ha
\[
  D\operatorname{arctg}x=\frac1{1+x^{2}}, \qquad
  D^{2}\operatorname{arctg}x=-\frac{2x}{(1+x^{2})^{2}},
\]
quindi l'arcotangente è convessa per $x\le0$ e concava per $x\ge0$, e il punto $0$ è di
flesso (\cref{fig:arctg}). La funzione $x^{2}$ è convessa su $\R$. Infine
$f(x)=x^{3}-3x$ ha $f''(x)=6x$: è concava per $x\le0$ e convessa per $x\ge0$, con un
flesso nell'origine, come si vede nella \cref{fig:cubica-monotonia}.
\end{esempio}

\begin{esempio}[Il grafico di $\sen x$]\label{es:grafico-seno}
Le proprietà stabilite negli ultimi due paragrafi permettono di \emph{motivare} il
grafico, ben noto, della funzione $\sen x$ nell'intervallo $[0,2\pi]$. Le derivate sono
$f'(x)=\cos x$ e $f''(x)=-\sen x$, e i loro segni sono riassunti nella
\cref{tab:seno-segni}. Ne segue che $\sen x$ cresce in $[0,\frac\pi2]$, decresce in
$[\frac\pi2,\frac32\pi]$ e cresce di nuovo in $[\frac32\pi,2\pi]$; è concava in
$[0,\pi]$ e convessa in $[\pi,2\pi]$. Il punto $\frac\pi2$ è di massimo, $\frac32\pi$ di
minimo, $\pi$ di flesso (a tangente obliqua, con pendenza $\cos\pi=-1$). Queste
informazioni, insieme a pochi valori facilmente calcolabili, determinano il grafico
della \cref{fig:seno-studio}. In modo analogo il lettore può studiare $\cos x$.
\end{esempio}

\begin{table}[htb]
\centering
\small
\renewcommand{\arraystretch}{1.35}
\setlength{\tabcolsep}{4.5pt}
\begin{tabular}{l|cccc}
  $f(x)=\sen x$ & $0<x<\frac\pi2$ & $\frac\pi2<x<\pi$ & $\pi<x<\frac32\pi$ &
  $\frac32\pi<x<2\pi$\\\hline
  segno di $f$   & $+$ & $+$ & $-$ & $-$\\
  segno di $f'$  & $+$ & $-$ & $-$ & $+$\\
  segno di $f''$ & $-$ & $-$ & $+$ & $+$\\\hline
  monotonia   & crescente & \multicolumn{2}{c}{decrescente} & crescente\\
  concavità   & \multicolumn{2}{c}{concava} & \multicolumn{2}{c}{convessa}
\end{tabular}
\caption{Segni di $f(x)=\sen x$ e delle sue derivate $f'(x)=\cos x$ e
$f''(x)=-\sen x$ in $(0,2\pi)$, e le conseguenti informazioni di monotonia e di
concavità.}
\label{tab:seno-segni}
\end{table}

\begin{figure}[htb]
\centering
\begin{tikzpicture}
\begin{axis}[width=0.9\textwidth, height=5cm, xmin=-0.3, xmax=7, ymin=-1.35, ymax=1.35,
  xtick={1.5708,3.1416,4.7124,6.2832},
  xticklabels={$\frac\pi2$,$\pi$,{},$2\pi$}, ytick={-1,1},
  xlabel={$x$}, ylabel={$y$}, clip=false]
  \addplot[curva, domain=0:6.2832, samples=150] {sin(deg(x))};
  \addplot[curva2, domain=2.3:4.0, samples=2] {-(x-3.1416)};
  \addplot[tratteggio] coordinates {(1.5708,0) (1.5708,1) (0,1)};
  \addplot[tratteggio] coordinates {(4.7124,0) (4.7124,-1) (0,-1)};
  \addplot[only marks, mark=*, mark size=1.6pt, rame] coordinates
    {(1.5708,1) (3.1416,0) (4.7124,-1)};
  \node[font=\scriptsize, rame, above] at (axis cs:1.5708,1.03) {massimo};
  \node[font=\scriptsize, rame, below] at (axis cs:4.7124,-1.03) {minimo};
  \node[font=\scriptsize, rame, anchor=south west] at (axis cs:3.2,0.08) {flesso};
  \node[font=\scriptsize, anchor=south] at (axis cs:4.7124,0.06) {$\frac32\pi$};
\end{axis}
\end{tikzpicture}
\caption{Il grafico di $\sen x$ in $[0,2\pi]$, ricostruito dalla \cref{tab:seno-segni}.
Nel flesso $x=\pi$ la tangente $y=-(x-\pi)$ attraversa il grafico.}
\label{fig:seno-studio}
\end{figure}

Il programma seguente studia la convessità della funzione $e^{-x^{2}}$, il cui grafico
è la celebre «curva a campana» di Gauss che descrive la distribuzione degli errori
casuali di misura.

\begin{python}[Convessità e flessi della curva di Gauss]
import sympy as sp

x = sp.symbols("x", real=True)
f = sp.exp(-x**2)                       # curva di Gauss
d1 = sp.factor(sp.diff(f, x))
d2 = sp.factor(sp.diff(f, x, 2))
print("f'  =", d1)
print("f'' =", d2)
flessi = sorted(sp.solveset(d2, x, sp.S.Reals))
approx = [round(float(c), 4) for c in flessi]
print("f'' = 0 in", flessi, "~", approx)
for t in (-1, 0, 1):              # segno di f'' negli intervalli
    print(f"  f''({t}) = {d2.subs(x, t)}  ->",
          "convessa" if d2.subs(x, t) > 0 else "concava")
print("f(1/sqrt(2)) =", f.subs(x, flessi[1]),
      "~", round(float(f.subs(x, flessi[1])), 4))
print("punto critico x = 0:  f''(0) =",
      d2.subs(x, 0), "-> massimo relativo")
\end{python}
\begin{risultato}
f'  = -2*x*exp(-x**2)
f'' = 2*(2*x**2 - 1)*exp(-x**2)
f'' = 0 in [-sqrt(2)/2, sqrt(2)/2] ~ [-0.7071, 0.7071]
  f''(-1) = 2*exp(-1)  -> convessa
  f''(0) = -2  -> concava
  f''(1) = 2*exp(-1)  -> convessa
f(1/sqrt(2)) = exp(-1/2) ~ 0.6065
punto critico x = 0:  f''(0) = -2 -> massimo relativo
\end{risultato}

Poiché il fattore $e^{-x^{2}}$ è positivo, il segno di $f''$ è quello di $2x^{2}-1$: la
curva di Gauss è convessa per $\abs x\ge\frac1{\sqrt2}$ e concava per
$\abs x\le\frac1{\sqrt2}$, e i due punti $x=\pm\frac1{\sqrt2}\approx\pm0{,}7071$ sono di
flesso, con ordinata $e^{-1/2}\approx0{,}6065$. L'ultima riga dell'output anticipa il
criterio che dimostriamo alla fine del paragrafo.

\begin{idea}[L'efficienza massima vista con la convessità]
Ritorniamo ancora una volta alla polare parabolica, pensando ora il coefficiente di
resistenza come funzione del coefficiente di portanza:
$C_{D}(C_{L})=C_{D0}+kC_{L}^{2}$. Poiché $C_{D}''=2k>0$, la polare è strettamente
convessa. L'efficienza $E=\frac{C_{L}}{C_{D}}$ è l'inverso della pendenza della retta che
congiunge l'origine con il punto $(C_{L},C_{D}(C_{L}))$; massimizzare $E$ significa
quindi cercare la retta per l'origine di pendenza \emph{minima} che tocca la polare
(\cref{fig:polare-tangente}). La tangente in $C_{L}^{*}$,
$y=C_{D}(C_{L}^{*})+2kC_{L}^{*}(C_{L}-C_{L}^{*})$, passa per l'origine se e solo se
$C_{D0}+kC_{L}^{*2}-2kC_{L}^{*2}=0$, cioè $C_{L}^{*}=\sqrt{C_{D0}/k}$, e in tal caso ha
equazione $y=2kC_{L}^{*}C_{L}=2\sqrt{kC_{D0}}\,C_{L}$. Per la stretta convessità la polare
sta sopra questa tangente:
\[
  C_{D}(C_{L})>2\sqrt{kC_{D0}}\;C_{L} \quad\text{per } C_{L}\neq C_{L}^{*},
\]
e quindi, dividendo per $C_{L}\,C_{D}(C_{L})>0$,
\[
  E(C_{L})=\frac{C_{L}}{C_{D}(C_{L})}<\frac1{2\sqrt{kC_{D0}}}=E_{\max}
  \qquad\text{per } C_{L}>0,\ C_{L}\neq C_{L}^{*}.
\]
Ritroviamo il risultato del \cref{cap:funzioni-continue} con una dimostrazione puramente
geometrica, che è anche la costruzione grafica usata dagli aerodinamici: \emph{l'assetto
di massima efficienza si trova tracciando dall'origine la tangente alla polare}.
\end{idea}

\begin{figure}[htb]
\centering
\begin{tikzpicture}
\begin{axis}[width=0.85\textwidth, height=6cm, xmin=0, xmax=1.62, ymin=0, ymax=0.15,
  axis lines=left, xlabel={$C_{L}$}, ylabel={$C_{D}$}, xtick={0.709,1.2},
  xticklabels={$C_{L}^{*}$,$C_{L}$}, ytick={0.025,0.05},
  yticklabels={$C_{D0}$,$0{,}05$}, clip=false,
  scaled y ticks=false]
  \addplot[curva, domain=0:1.55, samples=100] {0.025+0.0497359*x^2};
  \addplot[curva2, domain=0:1.55, samples=2] {x/14.17963};
  \addplot[curva3, domain=0:1.4, samples=2] {0.0805164*x};
  \addplot[tratteggio] coordinates {(0.709,0) (0.709,0.05)};
  \addplot[tratteggio] coordinates {(1.2,0) (1.2,0.09662)};
  \addplot[only marks, mark=*, mark size=1.6pt, rame] coordinates {(0.709,0.05)};
  \addplot[only marks, mark=*, mark size=1.6pt, verde] coordinates {(1.2,0.09662)};
  \node[font=\scriptsize, blu, anchor=east] at (axis cs:1.38,0.13) {polare};
  \node[font=\scriptsize, rame, anchor=west] at (axis cs:1.56,0.108)
    {pendenza $\frac1{E_{\max}}$};
  \node[font=\scriptsize, verde!50!black, anchor=south east] at (axis cs:1.15,0.098)
    {pendenza $\frac1{E(C_{L})}$};
\end{axis}
\end{tikzpicture}
\caption{La polare parabolica $C_{D}=C_{D0}+kC_{L}^{2}$ per $C_{D0}=0{,}025$, $AR=8$,
$e=0{,}8$. Ogni retta per l'origine che incontra la polare ha pendenza $\frac1E$; la
tangente dall'origine (rame) ha la pendenza minima e tocca la polare in
$C_{L}^{*}\approx0{,}709$, dove $E=E_{\max}\approx14{,}18$. La secante verde, per
$C_{L}=1{,}2$, ha pendenza maggiore: $E\approx12{,}4$.}
\label{fig:polare-tangente}
\end{figure}

\subsection{Il criterio della derivata seconda}

Chiudiamo il paragrafo con un criterio, basato sul segno della derivata seconda, per
stabilire se un punto stazionario è di massimo o di minimo relativo. Lo generalizzeremo
nel \cref{sec:taylor}.

\begin{teorema}[Criterio della derivata seconda]\label{teo:test-derivata-seconda}
\index{criterio!della derivata seconda}
Sia $f$ derivabile due volte in un intorno di $x_{0}$, con derivata seconda continua in
$x_{0}$. Allora
\begin{equation}\label{eq:test-derivata-seconda}
  \begin{aligned}
  &f'(x_{0})=0,\ f''(x_{0})>0 \ \implica\ x_{0} \text{ di minimo relativo stretto};\\
  &f'(x_{0})=0,\ f''(x_{0})<0 \ \implica\ x_{0} \text{ di massimo relativo stretto}.
  \end{aligned}
\end{equation}
\end{teorema}

\begin{schema}
\Ipotesi{$f$ derivabile due volte vicino a $x_{0}$, $f''$ continua in $x_{0}$;
$f'(x_{0})=0$ e $f''(x_{0})>0$ (risp.\ $<0$).}
\Tesi{$x_{0}$ è di minimo (risp.\ massimo) relativo stretto.}
\end{schema}

\begin{strategia}
Se $f''(x_{0})>0$ e $f''$ è continua, per la permanenza del segno $f''$ è positiva in
tutto un intorno di $x_{0}$; lì la funzione è strettamente convessa e il grafico sta
sopra la tangente in $x_{0}$. Ma la tangente in un punto stazionario è la retta
orizzontale $y=f(x_{0})$: dunque il grafico sta sopra il livello $f(x_{0})$, che è la
definizione di minimo relativo.
\end{strategia}

\begin{dimostrazione}
Sia $f''(x_{0})>0$. Per il teorema della permanenza del segno per le funzioni continue
(\cref{teo:permanenza-segno-continue}) esiste $\delta>0$ tale che $f''(x)>0$ per ogni
$x\in(x_{0}-\delta,x_{0}+\delta)$, intervallo che possiamo supporre contenuto
nell'intorno in cui $f$ è derivabile due volte. Fissato $0<\delta'<\delta$, per il
criterio di convessità e l'\cref{oss:convessita-stretta}, $f$ è strettamente convessa in
$[x_{0}-\delta',x_{0}+\delta']$, e quindi, ricordando che $f'(x_{0})=0$,
\begin{equation}\label{eq:test-derivata-seconda-dim}
  f(x)>f(x_{0})+f'(x_{0})(x-x_{0})=f(x_{0})
\end{equation}
per ogni $x\in[x_{0}-\delta',x_{0}+\delta']$ con $x\neq x_{0}$.
Dunque $x_{0}$ è un punto di minimo relativo stretto. Il caso $f''(x_{0})<0$ si ottiene
applicando quanto provato a $-f$.
\end{dimostrazione}

\begin{osservazione}[Basta l'esistenza di $f''(x_{0})$]\label{oss:test-derivata-seconda}
La continuità di $f''$ in $x_{0}$ non è necessaria. Supponiamo soltanto che $f$ sia
derivabile in un intorno di $x_{0}$, con $f'(x_{0})=0$, e che esista $f''(x_{0})>0$.
Allora
\[
  f''(x_{0})=\lim_{x\to x_{0}}\frac{f'(x)-f'(x_{0})}{x-x_{0}}=\lim_{x\to x_{0}}\frac{f'(x)}{x-x_{0}}>0,
\]
e per la permanenza del segno (\cref{teo:confronto-funzioni}\,(i)) il rapporto
$\frac{f'(x)}{x-x_{0}}$ è positivo per $0<\abs{x-x_{0}}<\delta$. Quindi $f'(x)<0$ per
$x_{0}-\delta<x<x_{0}$ e $f'(x)>0$ per $x_{0}<x<x_{0}+\delta$, e il test della derivata
prima (\cref{cor:test-derivata-prima}) mostra che $x_{0}$ è di minimo relativo stretto.
Per la curva di Gauss, $f'(0)=0$ e $f''(0)=-2<0$: l'origine è di massimo relativo
(anzi assoluto, perché $e^{-x^{2}}\le1$). Se invece $f'(x_{0})=f''(x_{0})=0$, il criterio
non decide: $x^{4}$, $-x^{4}$ e $x^{3}$ hanno in $0$, rispettivamente, un minimo, un
massimo e un flesso. Il \cref{sec:taylor} risolverà questi casi.
\end{osservazione}

\subsection{Approfondimento: la convessità senza derivate}
\label{subsec:convessita-corde}

La \cref{def:convessa} confronta il grafico con le tangenti, e richiede quindi la
derivabilità; eppure anche $\abs x$, che ha uno spigolo nell'origine, «volge la concavità
verso l'alto». Seguendo \cite{giusti}, usiamo le \emph{corde} al posto delle tangenti:
per le funzioni derivabili si ritrova la nostra definizione, e dalla sola convessità
discendono la continuità e l'esistenza delle derivate unilaterali nei punti interni.

\begin{definizione}[Convessità nel senso delle corde]\label{def:convessa-corde}
\index{funzione!convessa!nel senso delle corde}
Sia $I$ un intervallo e sia $f:I\to\R$. Si dice che $f$ è \emph{convessa nel senso delle
corde} in $I$ se, per ogni $x_{1},x_{2}\in I$ e per ogni $t\in[0,1]$, risulta
\begin{equation}\label{eq:convessa-corde}
  f\bigl((1-t)x_{1}+tx_{2}\bigr)\le(1-t)f(x_{1})+tf(x_{2}) ;
\end{equation}
si dice che $f$ è \emph{concava nel senso delle corde} se $-f$ è convessa nel senso
delle corde.
\end{definizione}

Se $x_{1}<x_{2}$, al variare di $t$ in $[0,1]$ il punto $x=(1-t)x_{1}+tx_{2}$ descrive
$[x_{1},x_{2}]$, con $t=\frac{x-x_{1}}{x_{2}-x_{1}}$, e la condizione diventa
\begin{equation}\label{eq:corda-sopra}
  f(x)\le f(x_{1})+\frac{f(x_{2})-f(x_{1})}{x_{2}-x_{1}}\,(x-x_{1})
  \qquad (x_{1}\le x\le x_{2}) :
\end{equation}
sopra $[x_{1},x_{2}]$ il grafico sta sotto la corda di estremi $(x_{1},f(x_{1}))$ e
$(x_{2},f(x_{2}))$. Il caso $x_{1}=x_{2}$ è banale, e scambiare $x_{1}$ con $x_{2}$
equivale a sostituire $t$ con $1-t$. Poiché la condizione coinvolge solo i punti di
$[x_{1},x_{2}]$, $f$ è convessa nel senso delle corde in $I$ se e solo se lo è in ogni
intervallo chiuso e limitato contenuto in $I$. Ad esempio $\abs x$ lo è su $\R$, perché
$\abs{(1-t)x_{1}+tx_{2}}\le(1-t)\abs{x_{1}}+t\abs{x_{2}}$ per la disuguaglianza
triangolare (\cref{teo:triangolare}).

\begin{lemma}[delle tre corde]\label{lem:tre-corde}\index{lemma delle tre corde}
Sia $f$ convessa nel senso delle corde nell'intervallo $I$.
\begin{enumerate}[label=\textup{(\roman*)}]
  \item Se $x_{1}<x<x_{2}$ sono punti di $I$, allora (\cref{fig:tre-corde})
  \begin{equation}\label{eq:tre-corde}
    \frac{f(x)-f(x_{1})}{x-x_{1}}\le\frac{f(x_{2})-f(x_{1})}{x_{2}-x_{1}}
    \le\frac{f(x_{2})-f(x)}{x_{2}-x} .
  \end{equation}
  \item Per ogni $x_{0}\in I$ il rapporto incrementale
  $R(x)=\frac{f(x)-f(x_{0})}{x-x_{0}}$ è una funzione crescente in $I\setminus\{x_{0}\}$.
\end{enumerate}
\end{lemma}

\begin{schema}
\Ipotesi{$f$ convessa nel senso delle corde nell'intervallo $I$.}
\Tesi{(i) per $x_{1}<x<x_{2}$ in $I$ vale la \eqref{eq:tre-corde}; \ (ii) per ogni
$x_{0}\in I$, $R(x)=\frac{f(x)-f(x_{0})}{x-x_{0}}$ è crescente in $I\setminus\{x_{0}\}$.}
\end{schema}

\begin{strategia}
Il punto $P=(x,f(x))$ sta sotto la retta $P_{1}P_{2}$: abbassando l'estremo destro di un
segmento che parte da $P_{1}$ la pendenza diminuisce, abbassando l'estremo sinistro di un
segmento che arriva in $P_{2}$ la pendenza aumenta. Per la (ii) si applica la (i) alla
terna formata da $x_{0}$ e dai due punti da confrontare.
\end{strategia}

\begin{dimostrazione}
(i) Sia $m=\frac{f(x_{2})-f(x_{1})}{x_{2}-x_{1}}$ e sia $r(s)=f(x_{1})+m(s-x_{1})$, cosicché
$r(x_{1})=f(x_{1})$, $r(x_{2})=f(x_{2})$ e, per la \eqref{eq:corda-sopra}, $f(x)\le r(x)$.
Poiché $x-x_{1}>0$ e $x_{2}-x>0$,
\[
\begin{gathered}
  \frac{f(x)-f(x_{1})}{x-x_{1}}\le\frac{r(x)-r(x_{1})}{x-x_{1}}=m ,\\[2pt]
  \frac{f(x_{2})-f(x)}{x_{2}-x}\ge\frac{r(x_{2})-r(x)}{x_{2}-x}=m ,
\end{gathered}
\]
perché $r$ è un polinomio di primo grado di pendenza $m$: è la \eqref{eq:tre-corde}.

(ii) Il rapporto incrementale non cambia scambiando i due punti:
$R(x)=\frac{f(x_{0})-f(x)}{x_{0}-x}$. Siano $x<y$ due punti di $I\setminus\{x_{0}\}$. Se
$x_{0}<x<y$, la prima disuguaglianza della \eqref{eq:tre-corde}, scritta per la terna
$x_{0}<x<y$, dice che $R(x)\le R(y)$; se $x<x_{0}<y$, lo dice il confronto tra il primo e
il terzo membro per la terna $x<x_{0}<y$; se $x<y<x_{0}$, lo dice la seconda
disuguaglianza per la terna $x<y<x_{0}$. Dunque $R$ è crescente.
\end{dimostrazione}

\begin{figure}[htb]
\centering
\begin{tikzpicture}
\begin{axis}[width=0.8\textwidth, height=5cm, xmin=0, xmax=4.1, ymin=0, ymax=4.9,
  axis lines=left, xtick={0.5,2,3.5}, xticklabels={$x_{1}$,$x$,$x_{2}$},
  ytick=\empty, clip=false]
  \addplot[curva, domain=0.2:3.7, samples=120] {0.35*x^2+0.3};
  \addplot[tratteggio] coordinates {(0.5,0) (0.5,0.3875)};
  \addplot[tratteggio] coordinates {(2,0) (2,1.7)};
  \addplot[tratteggio] coordinates {(3.5,0) (3.5,4.5875)};
  \draw[curva2] (axis cs:0.5,0.3875) -- (axis cs:2,1.7);
  \draw[curva3] (axis cs:0.5,0.3875) -- (axis cs:3.5,4.5875);
  \draw[grigio, line width=1pt] (axis cs:2,1.7) -- (axis cs:3.5,4.5875);
  \addplot[only marks, mark=*, mark size=1.6pt, blu] coordinates
    {(0.5,0.3875) (2,1.7) (3.5,4.5875)};
  \node[font=\footnotesize, rame, anchor=south] at (axis cs:1.75,1.56) {$s_{1}$};
  \node[font=\footnotesize, verde!50!black, anchor=south east] at (axis cs:1.45,1.85)
    {$s$};
  \node[font=\footnotesize, grigio, anchor=south] at (axis cs:2.2,2.2) {$s_{2}$};
  \node[font=\scriptsize, anchor=south east] at (axis cs:0.47,0.42) {$P_{1}$};
  \node[font=\scriptsize, anchor=north west] at (axis cs:2.04,1.66) {$P$};
  \node[font=\scriptsize, anchor=north west] at (axis cs:3.54,4.55) {$P_{2}$};
\end{axis}
\end{tikzpicture}
\caption{Il lemma delle tre corde. Se $f$ è convessa nel senso delle corde e
$x_{1}<x<x_{2}$, il punto $P$ sta sotto la corda $P_{1}P_{2}$, e le pendenze delle corde
$P_{1}P$, $P_{1}P_{2}$, $PP_{2}$ crescono: $s_{1}\le s\le s_{2}$. Qui
$f(x)=0{,}35\,x^{2}+0{,}3$, con $s_{1}=0{,}875$, $s=1{,}4$ e $s_{2}=1{,}925$.}
\label{fig:tre-corde}
\end{figure}

\begin{proposizione}[Le due definizioni coincidono per le funzioni derivabili]
\label{prop:convessita-equivalenti}
Sia $f$ derivabile in $[a,b]$ (con derivate unilaterali negli estremi). Allora $f$ è
convessa in $[a,b]$ secondo la \cref{def:convessa} se e solo se è convessa nel senso
delle corde in $[a,b]$.
\end{proposizione}

\begin{schema}
\Ipotesi{$f$ derivabile in $[a,b]$.}
\Tesi{grafico sopra le tangenti, \eqref{eq:convessa} \ $\sse$ \ grafico sotto le corde,
\eqref{eq:convessa-corde}.}
\end{schema}

\begin{strategia}
Dalle tangenti alle corde: si scrive la \eqref{eq:convessa} con il punto di tangenza nel
punto $z$ in cui si stima $f$, e si combinano le disuguaglianze relative agli estremi
della corda. Dalle corde alle tangenti: il rapporto incrementale $R$, crescente per il
lemma delle tre corde, a destra di $x_{0}$ non è minore del suo limite $f'(x_{0})$, a
sinistra non è maggiore; si moltiplica poi per $x-x_{0}$.
\end{strategia}

\begin{dimostrazione}
($\Rightarrow$) È l'argomento dell'\cref{es:corde-convesse}. Siano $x_{1},x_{2}\in[a,b]$,
$t\in[0,1]$ e $z=(1-t)x_{1}+tx_{2}$. Per la \eqref{eq:convessa} con punto di tangenza $z$,
\[
  f(x_{1})\ge f(z)+f'(z)(x_{1}-z), \qquad f(x_{2})\ge f(z)+f'(z)(x_{2}-z) .
\]
Moltiplicando la prima per $1-t\ge0$, la seconda per $t\ge0$ e sommando, si ottiene
$(1-t)f(x_{1})+tf(x_{2})\ge f(z)+f'(z)\bigl[(1-t)x_{1}+tx_{2}-z\bigr]=f(z)$.

($\Leftarrow$) Siano $x_{0},x\in[a,b]$, $x\neq x_{0}$, e sia $R$ il rapporto incrementale
del \cref{lem:tre-corde}\,(ii). Se $x>x_{0}$, per ogni $s\in(x_{0},x)$ si ha
$R(s)\le R(x)$; per $s\to x_{0}^{+}$ il primo membro tende a $f'(x_{0})$
(\eqref{eq:derivata-xi} e \cref{prop:dx-sx}; se $x_{0}=a$ si tratta della derivata
destra), e per il teorema del confronto (\cref{teo:confronto-funzioni}\,(ii)) risulta
$f'(x_{0})\le R(x)$. Moltiplicando per $x-x_{0}>0$ si ottiene
$f(x)-f(x_{0})\ge f'(x_{0})(x-x_{0})$. Se $x<x_{0}$, allo stesso modo $R(x)\le f'(x_{0})$,
e moltiplicando per $x-x_{0}<0$ si ritrova la stessa disuguaglianza, cioè la
\eqref{eq:convessa}.
\end{dimostrazione}

Dunque $x^{2}$, $e^{x}$, $-\log x$ e tutte le funzioni convesse incontrate finora lo sono
anche nel senso delle corde. Una funzione convessa nel senso delle corde può non essere
derivabile, come $\abs x$; ma nei punti interni non può essere troppo irregolare.

\begin{teorema}[Regolarità delle funzioni convesse]\label{teo:convesse-continue}
Sia $f$ convessa nel senso delle corde nell'intervallo $I$, e sia $x_{0}$ un punto
interno di $I$. Allora:
\begin{enumerate}[label=\textup{(\roman*)}]
  \item esistono finite la derivata sinistra $f'_{-}(x_{0})$ e la derivata destra
  $f'_{+}(x_{0})$, e risulta $f'_{-}(x_{0})\le f'_{+}(x_{0})$;
  \item $f$ è continua in $x_{0}$;
  \item per ogni $m\in\bigl[f'_{-}(x_{0}),f'_{+}(x_{0})\bigr]$ si ha
  \begin{equation}\label{eq:retta-appoggio}
    f(x)\ge f(x_{0})+m(x-x_{0}) \qquad\perogni x\in I .
  \end{equation}
\end{enumerate}
\end{teorema}

\begin{schema}
\Ipotesi{$f$ convessa nel senso delle corde nell'intervallo $I$; \ $x_{0}$ interno a $I$.}
\Tesi{(i) $f'_{-}(x_{0})\le f'_{+}(x_{0})$, entrambe finite; \ (ii) $f$ continua in
$x_{0}$; \ (iii) vale la \eqref{eq:retta-appoggio}.}
\end{schema}

\begin{strategia}
Tutto discende dalla monotonia di $R$. Una funzione monotòna ha limiti unilaterali, finiti
se è limitata dalla parte giusta: qui interviene l'ipotesi che $x_{0}$ sia interno, perché
un punto $a'<x_{0}$ fornisce il minorante $R(a')$ a destra di $x_{0}$, e un punto
$b'>x_{0}$ un maggiorante a sinistra. La continuità si legge nell'identità
$f(x)-f(x_{0})=R(x)(x-x_{0})$; la (iii) si ottiene confrontando $R$ con i suoi limiti.
\end{strategia}

\begin{dimostrazione}
Siano $a',b'\in I$ con $a'<x_{0}<b'$; per il \cref{lem:tre-corde}, $R$ è crescente in
$I\setminus\{x_{0}\}$.

(i) La funzione che vale $R(a')$ in $x_{0}$ e coincide con $R$ in $(x_{0},b']$ è crescente
in $[x_{0},b']$, perché $R(a')\le R(x)$ per $x>x_{0}$. Per il \cref{teo:limite-monotone-funzioni}
esiste finito il limite di $R(x)$ per $x\to x_{0}^{+}$, che per la \eqref{eq:derivata-xi}
è $f'_{+}(x_{0})$. Allo stesso modo, la funzione che vale $R(b')$ in $x_{0}$ e coincide
con $R$ in $[a',x_{0})$ è crescente in $[a',x_{0}]$, e quindi esiste finita
$f'_{-}(x_{0})=\lim_{x\to x_{0}^{-}}R(x)$. Infine, se $x<x_{0}<y$, da $R(x)\le R(y)$ si
ottiene, per $x\to x_{0}^{-}$, $f'_{-}(x_{0})\le R(y)$ e poi, per $y\to x_{0}^{+}$,
$f'_{-}(x_{0})\le f'_{+}(x_{0})$ (\cref{teo:confronto-funzioni}\,(ii)).

(ii) Per $x\to x_{0}^{\pm}$ il fattore $R(x)$ ha limite finito e $x-x_{0}\to0$, quindi
$f(x)-f(x_{0})=R(x)(x-x_{0})\to0$ da entrambe le parti; per la \cref{prop:dx-sx}, $f$ è
continua in $x_{0}$.

(iii) Se $x>x_{0}$, per ogni $s\in(x_{0},x)$ è $R(s)\le R(x)$, e al limite per
$s\to x_{0}^{+}$ si ha $m\le f'_{+}(x_{0})\le R(x)$; moltiplicando per $x-x_{0}>0$ si
ottiene la \eqref{eq:retta-appoggio}. Se $x<x_{0}$, allo stesso modo
$R(x)\le f'_{-}(x_{0})\le m$, e si moltiplica per $x-x_{0}<0$. Per $x=x_{0}$ vale
l'uguaglianza.
\end{dimostrazione}

\begin{osservazione}[Rette d'appoggio, spigoli ed estremi dell'intervallo]
\label{oss:convesse-esempi}\index{retta d'appoggio}
La retta $y=f(x_{0})+m(x-x_{0})$ della \eqref{eq:retta-appoggio} si dice \emph{retta
d'appoggio}: passa per $(x_{0},f(x_{0}))$ e lascia tutto il grafico da una parte. Se $f$ è
derivabile in $x_{0}$ l'unica retta d'appoggio è la tangente; in uno spigolo le rette
d'appoggio sono infinite: per $\abs x$ si ha $f'_{-}(0)=-1$, $f'_{+}(0)=1$, e
$\abs x\ge mx$ per ogni $m\in[-1,1]$. Se $f$ è convessa e derivabile in un punto interno
$x_{0}$ con $f'(x_{0})=0$, la \eqref{eq:retta-appoggio} con $m=0$ dà $f(x)\ge f(x_{0})$
per ogni $x\in I$: \emph{per una funzione convessa ogni punto stazionario interno è di
minimo assoluto}. Infine, l'ipotesi che $x_{0}$ sia interno non si può togliere: la
funzione che vale $0$ in $[0,1)$ e $1$ in $1$ è convessa nel senso delle corde in $[0,1]$
(le corde hanno ordinate $\ge0$), ma è discontinua nell'estremo $1$.
\end{osservazione}

La \eqref{eq:convessa-corde} riguarda medie pesate di \emph{due} punti; per induzione la
si estende a un numero qualunque di punti.

\begin{teorema}[Disuguaglianza di Jensen]\label{teo:jensen}\index{disuguaglianza di Jensen}
Sia $f$ convessa nel senso delle corde nell'intervallo $I$, siano $x_{1},\dots,x_{n}\in I$
e siano $t_{1},\dots,t_{n}\ge0$ numeri con $t_{1}+\dots+t_{n}=1$. Allora il punto
$\sum_{k=1}^{n}t_{k}x_{k}$ appartiene a $I$ e
\begin{equation}\label{eq:jensen}
  f\Bigl(\sum_{k=1}^{n}t_{k}x_{k}\Bigr)\le\sum_{k=1}^{n}t_{k}f(x_{k}) .
\end{equation}
\end{teorema}

\begin{schema}
\Ipotesi{$f$ convessa nel senso delle corde in $I$; \ $x_{k}\in I$, $t_{k}\ge0$,
$\sum_{k}t_{k}=1$.}
\Tesi{$\sum_{k}t_{k}x_{k}\in I$ \ e \
$f\bigl(\sum_{k}t_{k}x_{k}\bigr)\le\sum_{k}t_{k}f(x_{k})$.}
\end{schema}

\begin{strategia}
Per induzione su $n$: per $n=2$ la tesi è la definizione. Per passare da $n$ a $n+1$ punti
si raggruppano i primi $n$ nella loro media pesata $y$, con i pesi divisi per la loro
somma: la combinazione diventa una combinazione dei \emph{due} punti $y$ e $x_{n+1}$, a cui
si applica la definizione, mentre $f(y)$ si stima con l'ipotesi induttiva. In termini
meccanici: il baricentro di masse poste sul grafico sta sopra il grafico.
\end{strategia}

\begin{dimostrazione}
Se $p$ e $q$ sono il minimo e il massimo dei punti $x_{k}$, risulta
$p=\sum_{k}t_{k}p\le\sum_{k}t_{k}x_{k}\le\sum_{k}t_{k}q=q$; poiché $I$ è un intervallo che
contiene $p$ e $q$, contiene anche $\sum_{k}t_{k}x_{k}$.

Proviamo la \eqref{eq:jensen} per induzione (\cref{teo:induzione}). Per $n=1$ è $t_{1}=1$
e i due membri coincidono. Supponiamola vera per $n$ punti e siano dati $n+1$ punti con
pesi $t_{1},\dots,t_{n+1}$; sia $s=t_{1}+\dots+t_{n}=1-t_{n+1}$. Se $s=0$ i due membri
valgono $f(x_{n+1})$. Se $s>0$, il punto $y=\sum_{k=1}^{n}\frac{t_{k}}{s}x_{k}$ è una media
pesata con pesi $\ge0$ di somma $1$: quindi $y\in I$ e, per l'ipotesi induttiva,
$f(y)\le\sum_{k=1}^{n}\frac{t_{k}}{s}f(x_{k})$. Poiché
$\sum_{k=1}^{n+1}t_{k}x_{k}=s\,y+t_{n+1}x_{n+1}$, la \eqref{eq:convessa-corde} (con
$t=t_{n+1}$) dà
\[
\begin{aligned}
  f\Bigl(\sum_{k=1}^{n+1}t_{k}x_{k}\Bigr)&\le s\,f(y)+t_{n+1}f(x_{n+1})\\
  &\le\sum_{k=1}^{n}t_{k}f(x_{k})+t_{n+1}f(x_{n+1})=\sum_{k=1}^{n+1}t_{k}f(x_{k}) .
\end{aligned}
\]
La tesi vale dunque per $n+1$ punti, e quindi per ogni $n$.
\end{dimostrazione}

\begin{corollario}[Media aritmetica e media geometrica]
\label{cor:media-aritmetica-geometrica}
\index{disuguaglianza tra media aritmetica e geometrica}
Per ogni $n\in\N$ e per ogni scelta di numeri positivi $a_{1},\dots,a_{n}$,
\begin{equation}\label{eq:am-gm}
  \sqrt[n]{a_{1}a_{2}\cdots a_{n}}\le\frac{a_{1}+a_{2}+\dots+a_{n}}{n} .
\end{equation}
\end{corollario}

\begin{schema}
\Ipotesi{$n\in\N$; \ $a_{1},\dots,a_{n}>0$.}
\Tesi{la media geometrica non supera la media aritmetica: vale la \eqref{eq:am-gm}.}
\end{schema}

\begin{strategia}
Il logaritmo della media geometrica è la media aritmetica dei logaritmi. La tesi diventa
così «il logaritmo della media non è minore della media dei logaritmi»: è la
disuguaglianza di Jensen per la funzione convessa $-\log x$, con i pesi $t_{k}=\frac1n$.
\end{strategia}

\begin{dimostrazione}
La funzione $\varphi(x)=-\log x$ ha $\varphi''(x)=\frac1{x^{2}}>0$ in $(0,+\infty)$
(\cref{es:convessita-elementari}); per il \cref{teo:criterio-convessita} e la
\cref{prop:convessita-equivalenti} è convessa nel senso delle corde in ogni intervallo
$[a,b]\subset(0,+\infty)$, e quindi in $(0,+\infty)$. La \eqref{eq:jensen} con
$x_{k}=a_{k}$ e $t_{k}=\frac1n$ e le proprietà dei logaritmi (\cref{prop:log-proprieta})
danno
\[
  -\log\Bigl(\frac1n\sum_{k=1}^{n}a_{k}\Bigr)\le-\frac1n\sum_{k=1}^{n}\log a_{k}
  =-\log\sqrt[n]{a_{1}a_{2}\cdots a_{n}} .
\]
Cambiando segno, il logaritmo della media geometrica non supera quello della media
aritmetica; poiché il logaritmo è strettamente crescente (\cref{prop:log-monotono}), lo
stesso vale per le due medie.
\end{dimostrazione}

Per $n=2$ si ritrova la disuguaglianza $\sqrt{ab}\le\frac{a+b}2$, dimostrata per via
algebrica negli esercizi del \cref{cap:numeri}. Se gli $a_{k}$ sono tutti uguali, i due
membri della \eqref{eq:am-gm} coincidono; con la convessità stretta di $-\log x$ si
dimostra che è l'unico caso di uguaglianza.

\begin{esempio}[La potenza minima senza derivate]\label{es:potenza-am-gm}
Con termini scelti ad arte, la \eqref{eq:am-gm} trova minimi senza derivate: si maggiora
(o si minora) la funzione con una costante, e si cerca dove la disuguaglianza diventa
un'uguaglianza. Nel \cref{sec:studio-funzione} vedremo che la potenza necessaria al volo
livellato è $P(V)=aV^{3}+\frac bV$, con $a,b>0$. Spezzando $\frac bV$ in tre parti uguali,
in modo che nel prodotto dei quattro termini $V$ si semplifichi,
\[
  P(V)=aV^{3}+3\cdot\frac{b}{3V}\ge4\sqrt[4]{aV^{3}\Bigl(\frac{b}{3V}\Bigr)^{3}}
  =4\sqrt[4]{\frac{ab^{3}}{27}} ,
\]
con l'uguaglianza quando i quattro termini sono uguali: $aV^{3}=\frac{b}{3V}$, cioè
$V^{4}=\frac{b}{3a}$, la velocità di minima potenza \eqref{eq:V-min-P}. Per il velivolo
del \cref{sec:studio-funzione} si trova $P_{\min}\approx23{,}47$~kW, come là con le
derivate. Con due soli termini lo stesso procedimento dà la resistenza minima e, nel
\cref{cap:funzioni-continue}, l'efficienza massima.
\end{esempio}

\begin{notastorica}[Le funzioni convesse]
Le funzioni convesse comparvero in modo implicito in molti lavori dell'Ottocento sulle
disuguaglianze tra medie, ma la loro teoria sistematica fu iniziata dal matematico
danese Johan Jensen, che nel 1906 pubblicò negli \emph{Acta Mathematica} la memoria
\emph{Sur les fonctions convexes et les inégalités entre les valeurs moyennes}. Jensen
adottò una definizione che non richiede la derivabilità, basata sui punti medi delle
corde anziché sulle tangenti: $f\bigl(\frac{x+y}2\bigr)\le\frac{f(x)+f(y)}2$ per ogni $x,y$.
Per le funzioni continue essa equivale alla \eqref{eq:convessa-corde}, cioè al fatto che
ogni corda stia sopra il grafico, che è la definizione oggi più comune: per le funzioni
derivabili coincide con la nostra (\cref{prop:convessita-equivalenti}), ma si applica
anche a funzioni non derivabili, come $\abs x$. Nella stessa memoria Jensen dimostrò in
generale la disuguaglianza \eqref{eq:jensen}, già ottenuta da Otto Hölder nel 1889 per
le funzioni con $f''\ge0$. Oggi la convessità è al centro dell'\emph{ottimizzazione
convessa}, la teoria che garantisce che un minimo locale sia anche globale e che sta
alla base di molti algoritmi di progetto ottimo in ingegneria.
\end{notastorica}

%% ----------------------------------------------------------------------------
%%  §6.5  Il teorema di L'Hôpital
%% ----------------------------------------------------------------------------
\section{Il teorema di L'Hôpital}
\label{sec:hopital}

Siano $f$ e $g$ due funzioni che tendono a zero per $x\to x_{0}$. Abbiamo visto nel
\cref{cap:funzioni-continue} che il rapporto $\frac{f(x)}{g(x)}$ è allora una forma
indeterminata: in generale non è possibile dedurre il risultato del limite dai limiti
di $f$ e di $g$, ma occorre trasformare preliminarmente il rapporto in modo da togliere
l'indeterminazione. Il teorema di L'Hôpital serve a questo scopo, e sostituisce al
rapporto delle funzioni il rapporto delle loro derivate.

L'idea si vede bene nella \cref{fig:hopital-tangenti}. Se $f(x_{0})=g(x_{0})=0$, vicino a
$x_{0}$ le due funzioni sono approssimate dalle loro rette tangenti,
$f(x)\approx f'(x_{0})(x-x_{0})$ e $g(x)\approx g'(x_{0})(x-x_{0})$; il fattore comune
$x-x_{0}$, responsabile dell'indeterminazione, si semplifica, e il rapporto
$\frac{f(x)}{g(x)}$ tende al rapporto delle pendenze $\frac{f'(x_{0})}{g'(x_{0})}$.

\begin{figure}[htb]
\centering
\begin{tikzpicture}
\begin{axis}[width=0.75\textwidth, height=6cm, xmin=-0.75, xmax=0.75, ymin=-1.15,
  ymax=1.15, xtick={-0.5,0.5}, ytick={-1,1}, xlabel={$x$}, ylabel={$y$}, clip=false]
  \addplot[curva, domain=-0.6:0.6, samples=100] {exp(x)-1};
  \addplot[curva2, domain=-0.6:0.6, samples=100] {sin(deg(2*x))};
  \addplot[tratteggio, domain=-0.6:0.6, samples=2] {x};
  \addplot[tratteggio, domain=-0.52:0.52, samples=2] {2*x};
  \node[font=\scriptsize, rame, anchor=west] at (axis cs:0.6,0.96) {$g(x)=\sen2x$};
  \node[font=\scriptsize, blu, anchor=west] at (axis cs:0.6,0.79) {$f(x)=e^{x}-1$};
  \node[font=\scriptsize, grigio, anchor=west] at (axis cs:0.6,0.6) {$y=x$};
  \node[font=\scriptsize, grigio, anchor=south] at (axis cs:0.52,1.04) {$y=2x$};
\end{axis}
\end{tikzpicture}
\caption{Vicino all'origine $f(x)=e^{x}-1$ e $g(x)=\sen2x$ si confondono con le loro
tangenti $y=x$ e $y=2x$ (tratteggiate): il rapporto $\frac{f(x)}{g(x)}$ si comporta come
$\frac{x}{2x}=\frac12$, il rapporto delle pendenze.}
\label{fig:hopital-tangenti}
\end{figure}

\begin{teorema}[di L'Hôpital]\label{teo:hopital}\index{teorema!di L'Hôpital}
Sia $I$ un intervallo aperto contenente il punto $x_{0}$ e siano $f$, $g$ funzioni
derivabili in $I\setminus\{x_{0}\}$, tali che
\begin{equation}\label{eq:hopital-ipotesi}
  \lim_{x\to x_{0}}f(x)=0, \qquad \lim_{x\to x_{0}}g(x)=0 .
\end{equation}
Se $g'(x)\neq0$ per ogni $x\in I\setminus\{x_{0}\}$ e se esiste, finito o infinito, il
limite del rapporto delle derivate, allora esiste anche il limite del rapporto delle
funzioni, e risulta
\begin{equation}\label{eq:hopital}
  \lim_{x\to x_{0}}\frac{f(x)}{g(x)}=\lim_{x\to x_{0}}\frac{f'(x)}{g'(x)} .
\end{equation}
Il teorema vale anche:
\begin{enumerate}[label=\textup{(\roman*)}]
  \item per i limiti destri ($x\to x_{0}^{+}$) e sinistri ($x\to x_{0}^{-}$), con $f$, $g$
  derivabili soltanto da una parte di $x_{0}$;
  \item per le forme indeterminate $\frac\infty\infty$, cioè supponendo, al posto della
  \eqref{eq:hopital-ipotesi}, che $f$ e $g$ tendano all'infinito (anzi, basta che
  $\abs{g(x)}\to+\infty$, senza alcuna ipotesi sul limite di $f$);
  \item per $x\to+\infty$ e per $x\to-\infty$, con $f$, $g$ derivabili in una semiretta.
\end{enumerate}
\end{teorema}

\begin{schema}
\Ipotesi{\begin{itemize}[nosep, leftmargin=1.2em]
  \item $f,g$ derivabili in $I\setminus\{x_{0}\}$, con $g'\neq0$;
  \item $f(x)\to0$ e $g(x)\to0$ (oppure $\abs{g(x)}\to+\infty$) per $x\to x_{0}$;
  \item esiste $\displaystyle\lim_{x\to x_{0}}\frac{f'(x)}{g'(x)}=\ell\in\overline{\R}$.
\end{itemize}}
\Tesi{$\displaystyle\lim_{x\to x_{0}}\frac{f(x)}{g(x)}=\ell$.}
\end{schema}

Le ipotesi implicano che anche $g(x)\neq0$ per $x\neq x_{0}$ vicino a $x_{0}$, cosicché
il rapporto $\frac{f}{g}$ è definito. La dimostrazione nel caso generale richiede un
nuovo strumento, il teorema di Cauchy, ed è svolta nel \cref{sec:cauchy-hopital} in
appendice. Qui ci limitiamo a provare il teorema in un caso particolare, ma
significativo, che rende rigorosa l'idea della \cref{fig:hopital-tangenti}.

\begin{proposizione}[Il caso delle funzioni derivabili in $x_{0}$]\label{prop:hopital-semplice}
Siano $f$, $g$ definite in un intorno di $x_{0}$, derivabili in $x_{0}$, con
\[
  f(x_{0})=g(x_{0})=0, \qquad g'(x_{0})\neq0 .
\]
Allora $g(x)\neq0$ per $x\neq x_{0}$ in un intorno di $x_{0}$, e
\begin{equation}\label{eq:hopital-semplice}
  \lim_{x\to x_{0}}\frac{f(x)}{g(x)}=\frac{f'(x_{0})}{g'(x_{0})} .
\end{equation}
Se inoltre $f'$ e $g'$ esistono in un intorno di $x_{0}$ e sono continue in $x_{0}$, il
secondo membro è uguale a $\lim_{x\to x_{0}}\frac{f'(x)}{g'(x)}$.
\end{proposizione}

\begin{schema}
\Ipotesi{$f,g$ derivabili in $x_{0}$; $f(x_{0})=g(x_{0})=0$; $g'(x_{0})\neq0$.}
\Tesi{$\displaystyle\frac{f(x)}{g(x)}\to\frac{f'(x_{0})}{g'(x_{0})}$ per $x\to x_{0}$.}
\end{schema}

\begin{strategia}
Poiché $f(x_{0})=g(x_{0})=0$, possiamo scrivere $f(x)=f(x)-f(x_{0})$ e
$g(x)=g(x)-g(x_{0})$; dividendo numeratore e denominatore per $x-x_{0}$, il rapporto
$\frac{f}{g}$ diventa il quoziente di due rapporti incrementali. Ciascuno ha limite
finito, e quello del denominatore è diverso da zero: si applica il teorema sul limite
del quoziente.
\end{strategia}

\begin{dimostrazione}
Per $x\neq x_{0}$ poniamo
\[
  F(x)=\frac{f(x)-f(x_{0})}{x-x_{0}}=\frac{f(x)}{x-x_{0}}, \qquad
  G(x)=\frac{g(x)-g(x_{0})}{x-x_{0}}=\frac{g(x)}{x-x_{0}} .
\]
Per definizione di derivata, $F(x)\to f'(x_{0})$ e $G(x)\to g'(x_{0})\neq0$. Per la
permanenza del segno (\cref{teo:confronto-funzioni}\,(i), applicato a $G$ oppure a $-G$),
esiste $\delta>0$ tale che $G(x)\neq0$, e quindi $g(x)\neq0$, per
$0<\abs{x-x_{0}}<\delta$. In tali punti
\[
  \frac{f(x)}{g(x)}=\frac{f(x)/(x-x_{0})}{g(x)/(x-x_{0})}=\frac{F(x)}{G(x)}
  \ \longrightarrow\ \frac{f'(x_{0})}{g'(x_{0})}
\]
per il teorema sulle operazioni con i limiti (\cref{teo:operazioni-funzioni}). Se infine
$f'$, $g'$ sono continue in $x_{0}$, risulta $\frac{f'(x)}{g'(x)}\to\frac{f'(x_{0})}{g'(x_{0})}$,
di nuovo per il limite del quoziente.
\end{dimostrazione}

\subsection{Esempi: le forme \texorpdfstring{$\frac00$}{0/0} e \texorpdfstring{$\frac\infty\infty$}{inf/inf}}

Calcoliamo il limite
\begin{equation}\label{eq:hopital-es1}
  \lim_{x\to0}\frac{e^{x}-1}{\sen2x} .
\end{equation}
Si tratta di una forma indeterminata $\frac00$, che verifica le ipotesi della
\cref{prop:hopital-semplice}: numeratore e denominatore sono nulli e derivabili in $0$, e
$g'(0)=2\cos0=2\neq0$. Quindi
\[
  \lim_{x\to0}\frac{e^{x}-1}{\sen2x}=\frac{e^{0}}{2\cos0}=\frac12 ,
\]
come suggeriva la \cref{fig:hopital-tangenti}. Un'altra situazione dello stesso tipo è
\[
  \lim_{x\to1}\frac{x^{2}-1}{\log x}=\frac{2x}{1/x}\Big|_{x=1}=2 .
\]

Consideriamo ora due limiti notevoli, analoghi ai limiti di successioni della
gerarchia degli infiniti (\cref{teo:gerarchia}), che si calcolano facilmente con il
teorema di L'Hôpital nella forma $\frac\infty\infty$. Sia $b$ un numero reale
positivo. Si ha
\begin{equation}\label{eq:hopital-log-potenza}
  \lim_{x\to+\infty}\frac{\log x}{x^{b}}=\lim_{x\to+\infty}\frac{1/x}{b\,x^{b-1}}
  =\lim_{x\to+\infty}\frac1{b\,x^{b}}=0 :
\end{equation}
il logaritmo tende all'infinito più lentamente di qualunque potenza positiva. Per il
confronto tra potenza ed esponenziale, sia $n$ il più piccolo intero $\ge b$ (cioè
$n=b$ se $b$ è intero, $n=\lfloor b\rfloor+1$ altrimenti). Derivando $n$ volte il
numeratore si
ottiene
\begin{equation}\label{eq:hopital-potenza-exp}
  \lim_{x\to+\infty}\frac{x^{b}}{e^{x}}=\lim_{x\to+\infty}\frac{b\,x^{b-1}}{e^{x}}=\dots
  =\lim_{x\to+\infty}\frac{b(b-1)\cdots(b-n+1)\,x^{b-n}}{e^{x}}=0 ,
\end{equation}
perché $b-n\le0$ e quindi $x^{b-n}$ è limitata per $x\ge1$. Il ragionamento va letto
\emph{da destra verso sinistra}: l'ultimo limite esiste ed è nullo; le forme intermedie
$\frac{x^{b-k}}{e^{x}}$ con $k<n$ sono del tipo $\frac\infty\infty$ (perché $b-k>0$), e il
teorema di L'Hôpital, applicato $n$ volte risalendo la catena, assicura che anche i
limiti precedenti esistono e valgono $0$. Le \eqref{eq:hopital-log-potenza} e
\eqref{eq:hopital-potenza-exp} estendono alle funzioni la gerarchia
$\log n\ll n^{b}\ll a^{n}$ delle successioni.

\subsection{Le altre forme indeterminate}

Il teorema di L'Hôpital è utile anche per le forme indeterminate $0\cdot\infty$ e
$\infty-\infty$, che si riconducono a un quoziente con semplici manipolazioni. Per il
prodotto si scrive $fg=\frac{f}{1/g}$, oppure $fg=\frac{g}{1/f}$, ottenendo una forma
$\frac00$ oppure $\frac\infty\infty$. Ad esempio, se $b>0$,
\begin{equation}\label{eq:hopital-xb-log}
  \lim_{x\to0^{+}}x^{b}\log x=\lim_{x\to0^{+}}\frac{\log x}{x^{-b}}
  =\lim_{x\to0^{+}}\frac{1/x}{-b\,x^{-b-1}}=\lim_{x\to0^{+}}\frac{x^{b}}{-b}=0 .
\end{equation}
Per una differenza si cerca di scriverla come un unico quoziente, ad esempio
riducendo allo stesso denominatore:
\begin{align*}
  \lim_{x\to0}\Bigl(\frac1{\sen x}-\frac1x\Bigr)&=\lim_{x\to0}\frac{x-\sen x}{x\sen x}
  =\lim_{x\to0}\frac{1-\cos x}{\sen x+x\cos x}\\
  &=\lim_{x\to0}\frac{\sen x}{2\cos x-x\sen x}=\frac02=0,
\end{align*}
dove il teorema è stato applicato due volte (le prime due forme sono $\frac00$, l'ultimo
rapporto è continuo in $0$).

Le forme indeterminate esponenziali $0^{0}$, $1^{\infty}$, $\infty^{0}$ si trattano
scrivendo $f(x)^{g(x)}=e^{g(x)\log f(x)}$ (per $f>0$): l'esponente $g(x)\log f(x)$ è una
forma $0\cdot\infty$, e per la continuità della funzione esponenziale
(\cref{teo:continuita-esponenziale,cor:limite-composta-continua}) basta calcolarne il
limite. Ad esempio, usando la \eqref{eq:hopital-xb-log} con $b=1$,
\begin{equation}\label{eq:hopital-xx}
  \lim_{x\to0^{+}}x^{x}=\lim_{x\to0^{+}}e^{x\log x}=e^{\,\lim_{x\to0^{+}}x\log x}=e^{0}=1 ;
\end{equation}
per la forma $1^{\infty}$,
\begin{equation}\label{eq:hopital-1-infinito}
  \lim_{x\to0}(1+\sen x)^{1/x}=e^{\,\lim\limits_{x\to0}\frac{\log(1+\sen x)}{x}}
  =e^{\,\lim\limits_{x\to0}\frac{\cos x}{1+\sen x}}=e ;
\end{equation}
e per la forma $\infty^{0}$, per la \eqref{eq:hopital-log-potenza} con $b=1$,
\[
  \lim_{x\to+\infty}x^{1/x}=e^{\,\lim\limits_{x\to+\infty}\frac{\log x}{x}}=e^{0}=1 .
\]

Il programma seguente verifica tutti questi limiti con la funzione \texttt{limit} di
\texttt{sympy}. Si osservi che \texttt{sympy} non usa il teorema di L'Hôpital, ma
algoritmi basati sugli sviluppi asintotici, di cui la formula di Taylor del
\cref{sec:taylor} è il primo esempio.

\begin{python}[Verifica dei limiti con sympy]
import sympy as sp

x = sp.symbols("x", real=True)
b = sp.Rational(1, 2)                    # un esponente positivo
esempi = [((sp.exp(x) - 1) / sp.sin(2*x), 0, "+-"),
          ((x**2 - 1) / sp.log(x), 1, "+-"),
          (sp.log(x) / x**b, sp.oo, "-"),
          (x**5 / sp.exp(x), sp.oo, "-"),
          (x**b * sp.log(x), 0, "+"),
          (1 / sp.sin(x) - 1 / x, 0, "+-"),
          (x**x, 0, "+"),
          ((1 + sp.sin(x))**(1/x), 0, "+-"),
          (x**(1/x), sp.oo, "-")]
for f, x0, verso in esempi:
    print(f"lim {str(f):>24} per x -> {x0}:",
          sp.limit(f, x, x0, verso))
\end{python}
\begin{risultato}
lim    (exp(x) - 1)/sin(2*x) per x -> 0: 1/2
lim        (x**2 - 1)/log(x) per x -> 1: 2
lim           log(x)/sqrt(x) per x -> oo: 0
lim             x**5*exp(-x) per x -> oo: 0
lim           sqrt(x)*log(x) per x -> 0: 0
lim           1/sin(x) - 1/x per x -> 0: 0
lim                     x**x per x -> 0: 1
lim      (sin(x) + 1)**(1/x) per x -> 0: E
lim                 x**(1/x) per x -> oo: 1
\end{risultato}

\subsection{Tre avvertenze}

\begin{attenzione}[Prima di derivare, controllare che la forma sia indeterminata]
Il teorema si applica soltanto alle forme $\frac00$ e ai rapporti il cui denominatore
tende all'infinito (in particolare alle forme $\frac\infty\infty$). Applicato a
sproposito dà risultati sbagliati: $\lim_{x\to0}\frac{1+x}{2+x}=\frac12$ per
sostituzione, mentre il rapporto delle derivate vale costantemente $1$. Inoltre si
derivano numeratore e denominatore \emph{separatamente}: la derivata del quoziente
$\bigl(\frac fg\bigr)'$ non ha nulla a che fare con il teorema.
\end{attenzione}

\begin{attenzione}[Il teorema fornisce una condizione sufficiente]
Se il limite di $\frac{f'}{g'}$ non esiste, il teorema non dice nulla: il limite di
$\frac fg$ può esistere ugualmente. Per $f(x)=x+\sen x$ e $g(x)=x$, per $x\to+\infty$,
\[
  \frac{f(x)}{g(x)}=1+\frac{\sen x}{x}\ \to\ 1, \qquad\text{mentre}\qquad
  \frac{f'(x)}{g'(x)}=1+\cos x
\]
non ha limite, perché oscilla tra $0$ e $2$ (\cref{cor:non-esistenza}). Il programma
seguente lo mostra numericamente nei punti $x=\frac{k\pi}{2}$. Un esempio dello stesso
tipo, con $x\to0$, è $\frac{x^{2}\sen(1/x)}{\sen x}$, che tende a $0$, mentre il rapporto
delle derivate contiene il termine $\frac{\cos(1/x)}{\cos x}$, privo di limite (si
ricordino le funzioni $x^{n}\sen\frac1x$ del \cref{sec:retta-tangente}).
\end{attenzione}

\begin{python}[Quando il rapporto delle derivate non ha limite]
import math

# f(x) = x + sen x, g(x) = x per x -> +oo: f/g -> 1, f'/g' oscilla
print("     x = k pi/2        f/g     f'/g' = 1 + cos x")
for k in (2000, 2001, 2002, 2003, 200000, 200001, 200002):
    x = k * math.pi / 2
    print(f"{x:15.4f}   {(x + math.sin(x)) / x:.6f}   "
          f"{1 + math.cos(x):.6f}")
\end{python}
\begin{risultato}
     x = k pi/2        f/g     f'/g' = 1 + cos x
      3141.5927   1.000000   2.000000
      3143.1634   1.000318   1.000000
      3144.7342   1.000000   0.000000
      3146.3050   0.999682   1.000000
    314159.2654   1.000000   2.000000
    314160.8362   1.000003   1.000000
    314162.4070   1.000000   0.000000
\end{risultato}

\begin{attenzione}[Il teorema di L'Hôpital non dimostra i limiti notevoli]
Si è tentati di «dimostrare» il limite notevole $\frac{\sen x}{x}\to1$ con il teorema di
L'Hôpital: $\lim\frac{\sen x}{x}=\lim\frac{\cos x}{1}=1$. Il calcolo è corretto, ma non è
una dimostrazione: la formula $D\sen x=\cos x$ è stata ricavata nel
\cref{sec:derivate-elementari} proprio dal limite $\frac{\sen x}x\to1$, e usarla per
dimostrare quel limite sarebbe un circolo vizioso. Lo stesso vale per
$\frac{e^{x}-1}{x}\to1$ e $\frac{\log(1+x)}{x}\to1$, da cui sono state ricavate le
derivate dell'esponenziale e del logaritmo. Il teorema di L'Hôpital è uno strumento per
calcolare limiti \emph{nuovi} a partire dalle derivate \emph{già note}.
\end{attenzione}

\begin{notastorica}[Il marchese e il suo maestro]
Guillaume François Antoine de L'Hôpital (1661--1704), marchese di Sainte-Mesme,
pubblicò nel 1696 l'\emph{Analyse des infiniment petits pour l'intelligence des lignes
courbes}, il primo manuale di calcolo differenziale della storia, che ebbe un enorme
successo e diffuse in Europa i metodi di Leibniz. La regola che porta il suo nome vi
compare nella sezione IX. Già allora si sapeva che il marchese aveva preso lezioni dal
giovane Johann Bernoulli, ma solo nel Novecento, con la pubblicazione della loro
corrispondenza, si scoprì che nel 1694 L'Hôpital aveva offerto a Bernoulli un vitalizio
in cambio della comunicazione esclusiva delle sue scoperte matematiche, con il diritto di
usarle a suo piacimento. Il teorema «di L'Hôpital» è dunque, con ogni probabilità, di
Bernoulli. Va detto, a onore del marchese, che nella prefazione dell'\emph{Analyse} egli
riconosceva il suo debito verso Leibniz e verso i fratelli Bernoulli. La grafia
originale del nome era «L'Hospital»; nella grafia francese moderna la \emph{s} muta è
sostituita dall'accento circonflesso.
\end{notastorica}

%% ----------------------------------------------------------------------------
%%  §6.6  Studio del grafico di una funzione
%% ----------------------------------------------------------------------------
\section{Studio del grafico di una funzione}
\label{sec:studio-funzione}

I risultati ottenuti in questo capitolo ci permettono di studiare l'andamento di una
funzione $f(x)$ e di disegnarne il grafico in modo qualitativamente corretto. Si può
procedere secondo lo schema seguente, che non va inteso come una ricetta rigida: in
ciascun caso conviene adattarlo, saltando i passi superflui e approfondendo quelli
rilevanti.

\begin{enumerate}[label=\textbf{\Alph*.}, leftmargin=2em]
  \item Si determina il \emph{dominio} (o \emph{insieme di definizione}) della funzione.
  \item Si esamina se la funzione gode di qualche \emph{simmetria}: se è \emph{pari},
  $f(-x)=f(x)$, il grafico è simmetrico rispetto all'asse $y$; se è \emph{dispari},
  $f(-x)=-f(x)$, è simmetrico rispetto all'origine; se è \emph{periodica} di periodo $T$,
  $f(x+T)=f(x)$, basta studiarla in un intervallo di lunghezza $T$. Quando è semplice
  farlo, si calcolano le intersezioni con gli assi e si studia il segno della funzione.
  \item Si calcolano i limiti agli estremi del dominio e si determinano gli eventuali
  \emph{asintoti orizzontali e verticali}.
  \item Si determinano, studiando il segno della derivata prima, gli \emph{intervalli in
  cui la funzione è crescente o decrescente} e i \emph{punti di massimo e di minimo
  relativo}, e si calcolano i valori di $f$ in tali punti.
  \item Si determinano, studiando il segno della derivata seconda, gli \emph{intervalli
  in cui la funzione è convessa o concava} e gli eventuali \emph{punti di flesso}, e si
  calcolano i valori di $f$ in tali punti.
  \item Si determinano gli eventuali \emph{asintoti obliqui}.
\end{enumerate}

\subsection{Gli asintoti}

Precisiamo le nozioni di asintoto che compaiono ai punti C ed F.

\begin{definizione}[Asintoti]\label{def:asintoti}\index{asintoto}
\begin{enumerate}[label=\textup{(\roman*)}]
  \item La retta $y=\ell$ è un \emph{asintoto orizzontale} per $x\to+\infty$ se
  $\lim_{x\to+\infty}f(x)=\ell\in\R$; analogamente per $x\to-\infty$.
  \item La retta $x=x_{0}$ è un \emph{asintoto verticale} se
  $\lim_{x\to x_{0}}f(x)=\pm\infty$; si parla di asintoto verticale destro o sinistro se
  il limite infinito è soltanto quello per $x\to x_{0}^{+}$ o per $x\to x_{0}^{-}$.
  \item La retta $y=mx+q$ è un \emph{asintoto obliquo} per $x\to+\infty$ se
  \begin{equation}\label{eq:asintoto-obliquo}
    \lim_{x\to+\infty}\bigl[f(x)-(mx+q)\bigr]=0 ;
  \end{equation}
  analogamente per $x\to-\infty$.
\end{enumerate}
\end{definizione}

La \eqref{eq:asintoto-obliquo} significa che, per $x\to+\infty$, la distanza verticale
tra il grafico della funzione e la retta $y=mx+q$ tende a zero: il grafico «si appoggia»
alla retta. L'asintoto orizzontale è il caso particolare $m=0$. La definizione non dice
come trovare $m$ e $q$; lo dice la proposizione seguente.

\begin{proposizione}[Ricerca dell'asintoto obliquo]\label{prop:asintoto-obliquo}
La retta $y=mx+q$ è un asintoto obliquo per $x\to+\infty$ se e solo se esistono finiti
i due limiti
\begin{equation}\label{eq:m-q}
  m=\lim_{x\to+\infty}\frac{f(x)}{x}, \qquad q=\lim_{x\to+\infty}\bigl[f(x)-mx\bigr].
\end{equation}
Analogamente per $x\to-\infty$.
\end{proposizione}

\begin{schema}
\Ipotesi{$f$ definita in una semiretta $(c,+\infty)$.}
\Tesi{$y=mx+q$ asintoto obliquo per $x\to+\infty$ \ $\sse$ \ valgono le \eqref{eq:m-q}
con $m,q\in\R$.}
\end{schema}

\begin{strategia}
Supponiamo che l'asintoto esista e cerchiamo di ricavare $m$ e $q$. L'osservazione chiave
è che, se la quantità $f(x)-mx-q$ tende a zero, \emph{a maggior ragione} tende a zero la
stessa quantità divisa per $x$, che tende all'infinito: in questo modo $q$ scompare e
resta un'equazione per $m$. Noto $m$, il valore di $q$ si legge direttamente nella
definizione. Il viceversa è immediato.
\end{strategia}

\begin{dimostrazione}
Supponiamo che valga la \eqref{eq:asintoto-obliquo}. Allora, per il teorema sul limite
del prodotto (una funzione infinitesima per una funzione infinitesima $\frac1x$),
\[
  0=\lim_{x\to+\infty}\frac{f(x)-(mx+q)}{x}=\lim_{x\to+\infty}\Bigl[\frac{f(x)}{x}-m-\frac qx\Bigr],
\]
e poiché $\frac qx\to0$, risulta $\lim_{x\to+\infty}\frac{f(x)}x=m$. Inoltre, sempre per la
\eqref{eq:asintoto-obliquo}, $f(x)-mx=\bigl[f(x)-mx-q\bigr]+q\to q$. Viceversa, se
valgono le \eqref{eq:m-q} con $m$, $q$ finiti, la seconda di esse dice che
$f(x)-mx-q\to0$, cioè la \eqref{eq:asintoto-obliquo}.
\end{dimostrazione}

Si noti che, se esiste l'asintoto orizzontale $y=\ell$ per $x\to+\infty$, è inutile
cercare un asintoto obliquo dalla stessa parte: infatti $\frac{f(x)}x\to0$, e quindi
$m=0$ e $q=\ell$. Si noti anche che può esistere finito $m$ senza che esista l'asintoto:
per $f(x)=x+\log x$ si ha $\frac{f(x)}x\to1$, ma $f(x)-x=\log x\to+\infty$. In tal caso
il grafico ha la «direzione» della retta $y=x$ senza avvicinarsi ad alcuna retta.

\subsection{Un esempio completo}

Secondo lo schema proposto, studiamo la funzione
\begin{equation}\label{eq:studio-xe1x}
  f(x)=x\,e^{1/x} .
\end{equation}

\textbf{A.} La funzione è definita per ogni $x\neq0$: il dominio è
$(-\infty,0)\cup(0,+\infty)$.

\textbf{B.} Non ci sono simmetrie evidenti: $f(-x)=-x\,e^{-1/x}$ non coincide né con
$f(x)$ né con $-f(x)$. Poiché $e^{1/x}>0$, la funzione non si annulla mai: il grafico non
interseca gli assi, e $f(x)$ ha il segno di $x$.

\textbf{C.} Calcoliamo i limiti agli estremi del dominio. Per $x\to\pm\infty$ si ha
$\frac1x\to0$ ed $e^{1/x}\to e^{0}=1$, quindi
\[
  \lim_{x\to\pm\infty}x\,e^{1/x}=\pm\infty ,
\]
e non ci sono asintoti orizzontali. Per $x\to0^{-}$ il limite non è una forma
indeterminata: $\frac1x\to-\infty$, quindi $e^{1/x}\to0$ e $x\,e^{1/x}\to0\cdot0=0$; per
$x\to0^{-}$ non c'è asintoto verticale. Per $x\to0^{+}$ si ha invece la forma
$0\cdot\infty$; con il teorema di L'Hôpital nella forma $\frac\infty\infty$,
\[
  \lim_{x\to0^{+}}x\,e^{1/x}=\lim_{x\to0^{+}}\frac{e^{1/x}}{1/x}
  =\lim_{x\to0^{+}}\frac{e^{1/x}\cdot(-1/x^{2})}{-1/x^{2}}=\lim_{x\to0^{+}}e^{1/x}=+\infty .
\]
(Allo stesso risultato si arriva ponendo $t=\frac1x\to+\infty$: il limite diventa
$\frac{e^{t}}{t}\to+\infty$, per la \eqref{eq:hopital-potenza-exp}.) La retta $x=0$ è un
asintoto verticale destro.

\textbf{D.} La derivata prima vale
\begin{equation}\label{eq:studio-derivata-prima}
  f'(x)=e^{1/x}+x\,e^{1/x}\Bigl(-\frac1{x^{2}}\Bigr)=e^{1/x}\Bigl(1-\frac1x\Bigr)
  =e^{1/x}\,\frac{x-1}{x} .
\end{equation}
Il fattore $e^{1/x}$ è positivo; quindi $f'(x)>0$ se $x<0$ oppure $x>1$, e $f'(x)<0$ se
$0<x<1$. Per il criterio di monotonia (in ciascun intervallo del dominio), la funzione
è crescente in $(-\infty,0)$ e in $[1,+\infty)$, ed è decrescente in $(0,1]$. Il punto
$x=1$ è di minimo relativo, e $f(1)=e$. Osserviamo anche che
$f'(x)=e^{1/x}-\frac{e^{1/x}}{x}\to0$ per $x\to0^{-}$ (posto $t=\frac1x\to-\infty$, si ha
$t\,e^{t}\to0$): il grafico arriva verso l'origine con tangente orizzontale.

\textbf{E.} La derivata seconda vale
\begin{equation}\label{eq:studio-derivata-seconda}
  f''(x)=e^{1/x}\Bigl(-\frac1{x^{2}}\Bigr)\Bigl(1-\frac1x\Bigr)+e^{1/x}\,\frac1{x^{2}}
  =\frac{e^{1/x}}{x^{2}}\Bigl(-1+\frac1x+1\Bigr)=\frac{e^{1/x}}{x^{3}} .
\end{equation}
È positiva per $x>0$ e negativa per $x<0$: la funzione è convessa in $(0,+\infty)$ e
concava in $(-\infty,0)$. Poiché $f$ non è definita in $0$, non ci sono punti di
flesso. Il criterio della derivata seconda conferma il minimo: $f''(1)=e>0$.

\textbf{F.} Poiché $f(x)\to\pm\infty$ per $x\to\pm\infty$, cerchiamo gli asintoti
obliqui con la \cref{prop:asintoto-obliquo}:
\begin{equation}\label{eq:studio-m}
  m=\lim_{x\to\pm\infty}\frac{f(x)}{x}=\lim_{x\to\pm\infty}e^{1/x}=e^{0}=1 ,
\end{equation}
e, ponendo $t=\frac1x\to0$ e usando il limite notevole \eqref{eq:ln-exp},
\begin{equation}\label{eq:studio-q}
  q=\lim_{x\to\pm\infty}\bigl(x\,e^{1/x}-x\bigr)=\lim_{x\to\pm\infty}\frac{e^{1/x}-1}{1/x}
  =\lim_{t\to0}\frac{e^{t}-1}{t}=1
\end{equation}
(lo stesso risultato si ottiene con il teorema di L'Hôpital). La retta $y=x+1$ è un
asintoto obliquo sia per $x\to+\infty$ sia per $x\to-\infty$.

Possiamo anche stabilire da che parte dell'asintoto si trova il grafico. Con la stessa
sostituzione $t=\frac1x$,
\[
  f(x)-(x+1)=\frac{e^{t}-1-t}{t},
\]
e il numeratore è positivo per $t\neq0$ per la \eqref{eq:exp-log-disuguaglianze}.
Quindi il grafico sta sopra l'asintoto per $x>0$ e sotto per $x<0$. Con tutti questi
elementi si esegue il disegno della \cref{fig:studio-xe1x}.

\begin{figure}[htb]
\centering
\begin{tikzpicture}
\begin{axis}[width=0.9\textwidth, height=7.5cm, xmin=-4.3, xmax=4.3, ymin=-3.7, ymax=6.4,
  xtick={1}, ytick={1,2.71828}, yticklabels={$1$,$e$}, xlabel={$x$}, ylabel={$y$},
  clip=false]
  \addplot[tratteggio, domain=-4.2:4.2, samples=2] {x+1};
  \addplot[curva, domain=-4.2:-0.02, samples=200] {x*exp(1/x)};
  \addplot[curva, domain=0.3532:4.2, samples=200] {x*exp(1/x)};
  \addplot[tratteggio] coordinates {(1,0) (1,2.71828) (0,2.71828)};
  \addplot[only marks, mark=*, mark size=1.6pt, rame] coordinates {(1,2.71828)};
  \addplot[only marks, mark=o, mark size=1.6pt, blu, fill=white] coordinates {(0,0)};
  \node[font=\scriptsize, blu, anchor=west] at (axis cs:0.45,5.6) {$y=x\,e^{1/x}$};
  \node[font=\scriptsize, grigio, anchor=north west] at (axis cs:2.6,3.4) {$y=x+1$};
\end{axis}
\end{tikzpicture}
\caption{Il grafico di $f(x)=x\,e^{1/x}$: asintoto verticale destro $x=0$, asintoto
obliquo $y=x+1$ per $x\to\pm\infty$ (sopra il grafico per $x<0$, sotto per $x>0$),
minimo relativo in $(1,e)$, convessità per $x>0$ e concavità per $x<0$. Per
$x\to0^{-}$ il grafico tende all'origine, che non gli appartiene, con tangente
orizzontale.}
\label{fig:studio-xe1x}
\end{figure}

Il programma seguente ripete con \texttt{sympy} i calcoli dei punti C, D, E, F.
Strumenti di questo tipo sono un ottimo controllo, ma non sostituiscono lo studio: il
programma fornisce i risultati, lo studio ne fornisce le ragioni.

\needspace{12\baselineskip}
\begin{python}[Lo studio di $x\,e^{1/x}$ con sympy]
import sympy as sp

x = sp.symbols("x", real=True)
f = x * sp.exp(1 / x)
print("limiti:", [sp.limit(f, x, p, d) for p, d in
                  ((-sp.oo, "+"), (0, "-"), (0, "+"), (sp.oo, "-"))])
d1 = sp.factor(sp.diff(f, x))
d2 = sp.factor(sp.diff(f, x, 2))
print("f'  =", d1, "   f' = 0 in", sp.solve(d1, x))
print("f'' =", d2)
print("f(1) =", f.subs(x, 1), "   f''(1) =", d2.subs(x, 1))
for p in (sp.oo, -sp.oo):                  # asintoti obliqui
    m = sp.limit(f / x, x, p)
    q = sp.limit(f - m * x, x, p)
    print(f"x -> {p}:  m = {m},  q = {q}")
\end{python}
\begin{risultato}
limiti: [-oo, 0, oo, oo]
f'  = (x - 1)*exp(1/x)/x    f' = 0 in [1]
f'' = exp(1/x)/x**3
f(1) = E    f''(1) = E
x -> oo:  m = 1,  q = 1
x -> -oo:  m = 1,  q = 1
\end{risultato}

\subsection{Un esempio aeronautico: resistenza e potenza necessaria}

Consideriamo un velivolo di peso $W$ e superficie alare $S$, in volo rettilineo
livellato alla velocità $V$ nell'aria di densità $\rho$. L'equilibrio verticale impone
che la portanza eguagli il peso, $\frac12\rho V^{2}SC_{L}=W$, da cui
$C_{L}=\frac{2W}{\rho V^{2}S}$; con la polare parabolica $C_{D}=C_{D0}+kC_{L}^{2}$, la
resistenza $D=\frac12\rho V^{2}SC_{D}$ diventa una funzione della sola velocità:
\begin{equation}\label{eq:resistenza-V}
  D(V)=aV^{2}+\frac{b}{V^{2}}, \qquad a=\frac12\rho SC_{D0}, \qquad b=\frac{2kW^{2}}{\rho S} .
\end{equation}
Il primo termine è la \emph{resistenza parassita}, che cresce con il quadrato della
velocità; il secondo è la \emph{resistenza indotta}, legata alla portanza, che
diminuisce al crescere della velocità. Studiamo $D$ nell'intervallo $(0,+\infty)$,
ricordando che il modello ha senso fisico soltanto al di sopra della velocità di stallo.

\textbf{C.} Si ha $D(V)\to+\infty$ sia per $V\to0^{+}$ (asintoto verticale $V=0$) sia
per $V\to+\infty$. Non ci sono asintoti obliqui, perché $\frac{D(V)}V\to+\infty$. Si noti
però che $D(V)-aV^{2}=\frac b{V^{2}}\to0$ per $V\to+\infty$: il grafico si avvicina
alla \emph{parabola} $y=aV^{2}$, che si comporta come un «asintoto curvilineo»; allo
stesso modo, per $V\to0^{+}$, il grafico si avvicina all'iperbole $y=\frac b{V^{2}}$.

\textbf{D.} La derivata è
\[
  D'(V)=2aV-\frac{2b}{V^{3}}=\frac{2\,(aV^{4}-b)}{V^{3}},
\]
negativa per $V<V_{\min D}$ e positiva per $V>V_{\min D}$, dove
\begin{equation}\label{eq:V-min-D}
  V_{\min D}=\Bigl(\frac ba\Bigr)^{1/4}=\sqrt{\frac{2W}{\rho S}}\,\Bigl(\frac{k}{C_{D0}}\Bigr)^{1/4} .
\end{equation}
Per il criterio di monotonia, $V_{\min D}$ è il punto di \emph{minimo assoluto} della
resistenza in $(0,+\infty)$. In esso $aV^{4}=b$, cioè $aV^{2}=\frac b{V^{2}}$: \emph{alla
velocità di minima resistenza la resistenza parassita e quella indotta sono uguali}. Il
valore minimo è $D_{\min}=2aV_{\min D}^{2}=2\sqrt{ab}=2W\sqrt{kC_{D0}}=\frac{W}{E_{\max}}$,
e il coefficiente di portanza corrispondente è proprio $C_{L}^{*}=\sqrt{C_{D0}/k}$:
volare alla velocità di minima resistenza significa volare all'assetto di massima
efficienza.

\textbf{E.} La derivata seconda $D''(V)=2a+\frac{6b}{V^{4}}$ è positiva: la curva è
convessa.

Più interessante è la \emph{potenza necessaria} al volo livellato, $P(V)=D(V)\,V$:
\[
  P(V)=aV^{3}+\frac bV, \qquad P'(V)=3aV^{2}-\frac b{V^{2}}=\frac{3aV^{4}-b}{V^{2}},
\]
con $P''(V)=6aV+\frac{2b}{V^{3}}>0$.
La potenza è convessa e ha il minimo assoluto in
\begin{equation}\label{eq:V-min-P}
  V_{\min P}=\Bigl(\frac{b}{3a}\Bigr)^{1/4}=\frac{V_{\min D}}{3^{1/4}}\approx0{,}760\,V_{\min D} .
\end{equation}
Il rapporto $\frac{V_{\min D}}{V_{\min P}}=3^{1/4}\approx1{,}316$ non dipende dal
velivolo: è una conseguenza della sola forma della polare parabolica. C'è anche
un'interpretazione geometrica: poiché $\frac{P(V)}V=D(V)$, la resistenza è la pendenza
della retta che congiunge l'origine con il punto $(V,P(V))$; la velocità di minima
resistenza è quindi il punto di contatto della \emph{tangente per l'origine} alla curva
della potenza (\cref{fig:resistenza-potenza}), un'altra versione della costruzione della
\cref{fig:polare-tangente}. Per un aliante, la velocità di minima resistenza è quella
di massima efficienza, cioè di massima distanza percorribile in planata, mentre quella
di minima potenza è la velocità di minima discesa, che si usa per rimanere in volo il più
a lungo possibile.

\begin{python}[Velocità di minima resistenza e di minima potenza]
import math

# velivolo leggero, volo livellato al livello del mare
W, S, rho = 10000.0, 16.0, 1.225        # peso (N), superficie (m^2)
CD0, AR, e = 0.025, 8.0, 0.8
k = 1 / (math.pi * e * AR)
a = 0.5 * rho * S * CD0                 # D(V) = a V^2 + b / V^2
b = 2 * k * W**2 / (rho * S)
D = lambda V: a * V**2 + b / V**2
P = lambda V: D(V) * V                  # potenza necessaria

V_D = (b / a) ** 0.25                   # minimo di D: D'(V) = 0
V_P = (b / (3 * a)) ** 0.25             # minimo di P: P'(V) = 0
print(f"k = {k:.5f}   a = {a:.4f}   b = {b:.4e}")
print(f"V_minD = {V_D:.2f} m/s   D_min = {D(V_D):.1f} N"
      f"   W/D_min = {W / D(V_D):.2f}")
print(f"V_minP = {V_P:.2f} m/s   P_min = {P(V_P) / 1000:.2f} kW")
print(f"V_minD / V_minP = {V_D / V_P:.4f}   3^(1/4) = {3**0.25:.4f}")
CLmax = 1.5                             # C_L massimo (stallo)
print(f"V_stallo = {math.sqrt(2 * W / (rho * S * CLmax)):.2f} m/s")
\end{python}
\begin{risultato}
k = 0.04974   a = 0.2450   b = 5.0751e+05
V_minD = 37.94 m/s   D_min = 705.2 N   W/D_min = 14.18
V_minP = 28.83 m/s   P_min = 23.47 kW
V_minD / V_minP = 1.3161   3^(1/4) = 1.3161
V_stallo = 26.08 m/s
\end{risultato}

Il rapporto $\frac W{D_{\min}}\approx14{,}18$ coincide con l'efficienza massima trovata
nel \cref{cap:funzioni-continue} per la stessa polare, e le due velocità ottime sono
entrambe superiori alla velocità di stallo, calcolata con $C_{L,\max}=1{,}5$: il
modello è applicabile.

\begin{figure}[htb]
\centering
\includegraphics[width=\textwidth]{resistenza_potenza}
\caption{Resistenza (a sinistra) e potenza necessaria (a destra) in volo livellato, in
funzione della velocità, per il velivolo del programma precedente. La resistenza è
minima dove le componenti parassita e indotta si eguagliano; la tangente per l'origine
alla curva della potenza la tocca nella velocità di minima resistenza
(\texttt{python/fig\_resistenza\_potenza.py}).}
\label{fig:resistenza-potenza}
\end{figure}

%% ----------------------------------------------------------------------------
%%  §6.7  La formula di Taylor
%% ----------------------------------------------------------------------------
\section{La formula di Taylor}
\label{sec:taylor}

Nel \cref{sec:retta-tangente} abbiamo visto che una funzione derivabile in $x_{0}$ si
può approssimare, vicino a $x_{0}$, con un polinomio di primo grado:
\begin{equation}\label{eq:taylor-ordine-1}
  f(x)=f(x_{0})+f'(x_{0})(x-x_{0})+R_{1}(x), \qquad
  \lim_{x\to x_{0}}\frac{R_{1}(x)}{x-x_{0}}=0 ,
\end{equation}
cioè $R_{1}(x)=o(x-x_{0})$ (\cref{prop:approssimazione-lineare}): il \emph{resto} $R_{1}$
tende a zero più rapidamente di $x-x_{0}$. Sapendo che una funzione è derivabile fino
a un ordine $n>1$, ci si può domandare se sia possibile ottenere un'approssimazione
migliore, cioè decomporre $f(x)$ nella somma di un polinomio di grado $n$ e di un resto
$R_{n}(x)$ che tende a zero più rapidamente di $(x-x_{0})^{n}$. La formula di Taylor
risponde affermativamente.

Useremo il simbolo di sommatoria introdotto nel \cref{sec:binomio}, e il fattoriale
$k!=1\cdot2\cdots k$ \eqref{eq:fattoriale}, con la convenzione $0!=1$, e la nozione di
funzione \emph{derivabile $n$ volte in un punto} precisata nella
\cref{not:derivabile-n-volte}.

\begin{definizione}[Polinomio di Taylor]\label{def:polinomio-taylor}
\index{polinomio di Taylor}
Sia $f$ derivabile $n$ volte in $x_{0}$. Si chiama \emph{polinomio di Taylor di ordine
$n$} di $f$ di punto iniziale $x_{0}$ il polinomio
\begin{equation}\label{eq:polinomio-taylor}
\begin{split}
  T_{n}(x)&=\sum_{k=0}^{n}\frac{f^{(k)}(x_{0})}{k!}(x-x_{0})^{k}\\
  &=f(x_{0})+f'(x_{0})(x-x_{0})+\frac{f''(x_{0})}{2}(x-x_{0})^{2}\\
  &\quad+\dots+\frac{f^{(n)}(x_{0})}{n!}(x-x_{0})^{n}
\end{split}
\end{equation}
(per $k=0$ si intende $f^{(0)}(x_{0})=f(x_{0})$). Se $x_{0}=0$ si parla anche di
\emph{polinomio di Mac Laurin}.
\end{definizione}

Perché proprio questi coefficienti? Il lemma seguente mostra che $T_{n}$ è il polinomio
di grado $\le n$ che ha \emph{le stesse derivate di $f$ in $x_{0}$, fino all'ordine
$n$}: generalizza la retta tangente, che ha in comune con $f$ il valore e la derivata
prima.

\begin{lemma}\label{lem:derivate-taylor}
Per $k=0,1,\dots,n$ risulta $T_{n}^{(k)}(x_{0})=f^{(k)}(x_{0})$. Inoltre
$T_{n}^{(n-1)}(x)=f^{(n-1)}(x_{0})+f^{(n)}(x_{0})(x-x_{0})$ per ogni $x$.
\end{lemma}

\begin{schema}
\Ipotesi{$f$ derivabile $n$ volte in $x_{0}$; $T_{n}$ definito dalla
\eqref{eq:polinomio-taylor}.}
\Tesi{$T_{n}^{(k)}(x_{0})=f^{(k)}(x_{0})$ per $0\le k\le n$; \
$T_{n}^{(n-1)}(x)=f^{(n-1)}(x_{0})+f^{(n)}(x_{0})(x-x_{0})$.}
\end{schema}

\needspace{6\baselineskip}
\begin{strategia}
$T_{n}$ è una combinazione delle potenze $(x-x_{0})^{j}$: basta sapere quanto vale la
derivata $k$-esima di ciascuna potenza nel punto $x_{0}$, e osservare che soltanto la
potenza di esponente $j=k$ dà un contributo non nullo.
\end{strategia}

\begin{dimostrazione}
Derivando $k$ volte la potenza $(x-x_{0})^{j}$ si ottiene
$j(j-1)\cdots(j-k+1)(x-x_{0})^{j-k}$ se $k\le j$, e $0$ se $k>j$. Nel punto $x=x_{0}$
questa espressione è nulla, tranne nel caso $j=k$, in cui vale $k!$. Quindi, nella
derivata $k$-esima di $T_{n}$ calcolata in $x_{0}$, sopravvive soltanto il termine con
$j=k$, che vale $\frac{f^{(k)}(x_{0})}{k!}\cdot k!=f^{(k)}(x_{0})$. Per la seconda
affermazione, derivando $n-1$ volte $T_{n}$ si annullano i termini con $j<n-1$, e
restano $\frac{f^{(n-1)}(x_{0})}{(n-1)!}(n-1)!$ e
$\frac{f^{(n)}(x_{0})}{n!}\,n!\,(x-x_{0})$.
\end{dimostrazione}

\begin{teorema}[Formula di Taylor con il resto di Peano]\label{teo:taylor}
\index{formula di Taylor}\index{teorema!di Taylor}
Sia $f$ una funzione derivabile $n$ volte nel punto $x_{0}$ ($n\ge1$). Posto
$R_{n}(x)=f(x)-T_{n}(x)$, risulta
\begin{equation}\label{eq:taylor}
  f(x)=\sum_{k=0}^{n}\frac{f^{(k)}(x_{0})}{k!}(x-x_{0})^{k}+R_{n}(x), \qquad
  \lim_{x\to x_{0}}\frac{R_{n}(x)}{(x-x_{0})^{n}}=0 ,
\end{equation}
cioè $R_{n}(x)=o\bigl((x-x_{0})^{n}\bigr)$ per $x\to x_{0}$.
\end{teorema}

\begin{schema}
\Ipotesi{$f$ derivabile $n$ volte in $x_{0}$.}
\Tesi{$f(x)-T_{n}(x)=o\bigl((x-x_{0})^{n}\bigr)$ per $x\to x_{0}$.}
\end{schema}

\begin{strategia}
Dobbiamo calcolare il limite di $\frac{R_{n}(x)}{(x-x_{0})^{n}}$, che è una forma
indeterminata $\frac00$. L'idea è di applicare ripetutamente il teorema di L'Hôpital:
per il \cref{lem:derivate-taylor} le derivate di $R_{n}=f-T_{n}$ fino all'ordine $n$ si
annullano in $x_{0}$, e quindi a ogni passo si ritrova una forma $\frac00$. Dopo $n-1$
applicazioni il numeratore è $f^{(n-1)}(x)-f^{(n-1)}(x_{0})-f^{(n)}(x_{0})(x-x_{0})$ e il
denominatore è un multiplo di $x-x_{0}$: il limite si calcola direttamente con la
definizione di $f^{(n)}(x_{0})$. Come per la \eqref{eq:hopital-potenza-exp}, la catena
di uguaglianze va poi letta \emph{all'indietro}: è l'esistenza dell'ultimo limite a
garantire, passo dopo passo, quella dei precedenti.
\end{strategia}

\begin{dimostrazione}
Per $n=1$ la tesi è la \eqref{eq:taylor-ordine-1}, cioè la
\cref{prop:approssimazione-lineare}. Sia $n\ge2$, e sia $I$ un intorno di $x_{0}$ in cui
esiste $f^{(n-1)}$. Poniamo, per $k=0,1,\dots,n-1$,
\begin{align*}
  N_{k}(x)&=R_{n}^{(k)}(x)=f^{(k)}(x)-T_{n}^{(k)}(x),\\
  D_{k}(x)&=\frac{d^{k}}{dx^{k}}(x-x_{0})^{n}=\frac{n!}{(n-k)!}(x-x_{0})^{n-k} .
\end{align*}
Osserviamo che:
\begin{itemize}
  \item per $k\le n-2$, $N_{k}$ e $D_{k}$ sono derivabili in $I$, con $N_{k}'=N_{k+1}$ e
  $D_{k}'=D_{k+1}$; inoltre $D_{k+1}(x)\neq0$ per $x\neq x_{0}$;
  \item per $k\le n-2$, $N_{k}(x)\to N_{k}(x_{0})=0$ (per la continuità e il
  \cref{lem:derivate-taylor}) e $D_{k}(x)\to0$ per $x\to x_{0}$: il rapporto
  $\frac{N_{k}}{D_{k}}$ è una forma indeterminata $\frac00$.
\end{itemize}
Calcoliamo l'ultimo rapporto. Per il \cref{lem:derivate-taylor},
$N_{n-1}(x)=f^{(n-1)}(x)-f^{(n-1)}(x_{0})-f^{(n)}(x_{0})(x-x_{0})$, e $D_{n-1}(x)=n!\,(x-x_{0})$;
quindi
\begin{equation}\label{eq:taylor-ultimo}
\begin{split}
  \frac{N_{n-1}(x)}{D_{n-1}(x)}&=\frac1{n!}\Bigl[\frac{f^{(n-1)}(x)-f^{(n-1)}(x_{0})}{x-x_{0}}
  -f^{(n)}(x_{0})\Bigr]\\
  &\longrightarrow\ \frac1{n!}\bigl[f^{(n)}(x_{0})-f^{(n)}(x_{0})\bigr]=0,
\end{split}
\end{equation}
per la definizione di derivata di $f^{(n-1)}$ nel punto $x_{0}$. Ora risaliamo: poiché
$\frac{N_{n-1}}{D_{n-1}}=\frac{N_{n-2}'}{D_{n-2}'}\to0$ e $\frac{N_{n-2}}{D_{n-2}}$ è una
forma $\frac00$ che verifica le ipotesi del teorema di L'Hôpital (\cref{teo:hopital}),
anche $\frac{N_{n-2}}{D_{n-2}}\to0$. Ripetendo il ragionamento per
$k=n-3,\dots,1,0$ si ottiene infine
\[
  \lim_{x\to x_{0}}\frac{N_{0}(x)}{D_{0}(x)}=\lim_{x\to x_{0}}\frac{R_{n}(x)}{(x-x_{0})^{n}}=0 .
  \qedhere
\]
\end{dimostrazione}

\begin{osservazione}[Sulle ipotesi e sull'uso del teorema di L'Hôpital]
\label{oss:taylor-ipotesi}
Una dimostrazione più breve, che si trova in molti testi, suppone che $f^{(n)}$ esista in
un intorno di $x_{0}$ e sia continua in $x_{0}$, e applica il teorema di L'Hôpital $n$
volte; la dimostrazione precedente, fermandosi a $n-1$ applicazioni e concludendo con la
definizione di derivata, richiede soltanto l'esistenza di $f^{(n)}(x_{0})$. Si noti che
per $n\ge2$ le derivate dei denominatori $D_{k}$ si annullano in $x_{0}$, e quindi non
basta la forma semplice del teorema di L'Hôpital (\cref{prop:hopital-semplice}): serve
il caso generale, dimostrato nel \cref{sec:cauchy-hopital} con il teorema di Cauchy, che
a sua volta poggia soltanto sul teorema di Rolle. Non c'è dunque alcun circolo vizioso.
\end{osservazione}

Per $n=2$ la formula di Taylor si scrive esplicitamente
\begin{equation}\label{eq:taylor-ordine-2}
  f(x)=f(x_{0})+f'(x_{0})(x-x_{0})+\frac{f''(x_{0})}{2}(x-x_{0})^{2}+o\bigl((x-x_{0})^{2}\bigr):
\end{equation}
la parabola $y=T_{2}(x)$, detta \emph{parabola osculatrice}, approssima il grafico
meglio della retta tangente, perché ne riproduce anche la curvatura. Il polinomio di
Taylor è inoltre l'\emph{unico} polinomio di grado $\le n$ che approssima $f$ con un
errore $o\bigl((x-x_{0})^{n}\bigr)$.

\begin{proposizione}[Unicità del polinomio di Taylor]\label{prop:unicita-taylor}
Se $P$ è un polinomio di grado $\le n$ tale che $f(x)-P(x)=o\bigl((x-x_{0})^{n}\bigr)$
per $x\to x_{0}$, e $f$ è derivabile $n$ volte in $x_{0}$, allora $P=T_{n}$.
\end{proposizione}

\begin{schema}
\Ipotesi{$f$ derivabile $n$ volte in $x_{0}$; $P$ polinomio di grado $\le n$ con
$f-P=o\bigl((x-x_{0})^{n}\bigr)$.}
\Tesi{$P=T_{n}$.}
\end{schema}

\begin{strategia}
Per la formula di Taylor anche $f-T_{n}$ è $o\bigl((x-x_{0})^{n}\bigr)$; quindi il polinomio
$Q=T_{n}-P$ è un infinitesimo di ordine superiore a $n$. Ma un polinomio non nullo di
grado $\le n$ è, vicino a $x_{0}$, un infinitesimo di ordine pari all'esponente $j$ del suo
primo coefficiente non nullo, che è $\le n$: si arriva a una contraddizione dividendo
per $(x-x_{0})^{j}$.
\end{strategia}

\begin{dimostrazione}
Per differenza, il polinomio $Q=T_{n}-P$ verifica $Q(x)=o\bigl((x-x_{0})^{n}\bigr)$.
Scriviamolo nella forma $Q(x)=\sum_{k=0}^{n}c_{k}(x-x_{0})^{k}$ e supponiamo per assurdo
che i coefficienti non siano tutti nulli; sia $j$ il più piccolo indice con
$c_{j}\neq0$. Allora
\[
  \frac{Q(x)}{(x-x_{0})^{j}}=c_{j}+c_{j+1}(x-x_{0})+\dots+c_{n}(x-x_{0})^{n-j}\ \to\ c_{j}\neq0,
\]
mentre d'altra parte
$\frac{Q(x)}{(x-x_{0})^{j}}=\frac{Q(x)}{(x-x_{0})^{n}}\,(x-x_{0})^{n-j}$ è il prodotto della
funzione infinitesima $\frac{Q(x)}{(x-x_{0})^{n}}$ per la funzione $(x-x_{0})^{n-j}$, limitata
vicino a $x_{0}$, e quindi tende a $0$. Contraddizione: quindi $Q=0$.
\end{dimostrazione}

L'unicità è molto utile in pratica: \emph{qualunque} procedimento produca un polinomio
che approssima $f$ con l'errore giusto fornisce il polinomio di Taylor, senza bisogno di
calcolare le derivate. Ad esempio, dallo sviluppo $e^{t}=1+t+\frac{t^{2}}2+\frac{t^{3}}6+o(t^{3})$
che vedremo tra poco, sostituendo $t=-x^{2}$ si ottiene
$e^{-x^{2}}=1-x^{2}+\frac{x^{4}}2-\frac{x^{6}}6+o(x^{6})=1-x^{2}+\frac{x^{4}}2+o(x^{5})$
(per il teorema sul limite delle funzioni composte, $o(t^{3})$ con $t=-x^{2}$ è $o(x^{6})$):
$1-x^{2}+\frac{x^{4}}2$ è dunque il polinomio di Taylor di ordine $4$, e anche di ordine
$5$, della curva di Gauss.

\subsection{Gli sviluppi delle funzioni elementari}

Sviluppiamo con la formula di Taylor, con $x_{0}=0$, alcune funzioni elementari.

\emph{Esponenziale.} Se $f(x)=e^{x}$, risulta $f^{(k)}(x)=e^{x}$ per ogni $k$, e quindi
$f^{(k)}(0)=1$:
\begin{equation}\label{eq:taylor-exp}
  e^{x}=1+x+\frac{x^{2}}{2}+\frac{x^{3}}{3!}+\dots+\frac{x^{n}}{n!}+o(x^{n})
  =\sum_{k=0}^{n}\frac{x^{k}}{k!}+o(x^{n}) .
\end{equation}

\emph{Logaritmo.} Se $f(x)=\log(1+x)$, per induzione si verifica che
$f^{(k)}(x)=\frac{(-1)^{k+1}(k-1)!}{(1+x)^{k}}$ per $k\ge1$: infatti per $k=1$ si ha
$f'(x)=\frac1{1+x}$, e derivando la formula per $k$ si ottiene quella per $k+1$. Quindi
$\frac{f^{(k)}(0)}{k!}=\frac{(-1)^{k+1}}{k}$ e
\begin{equation}\label{eq:taylor-log}
  \log(1+x)=x-\frac{x^{2}}2+\frac{x^{3}}3-\dots+(-1)^{n+1}\frac{x^{n}}{n}+o(x^{n}) .
\end{equation}

\emph{Seno e coseno.} Per le formule di addizione \eqref{eq:addizione-sen},
$\cos x=\sen\bigl(x+\frac\pi2\bigr)$ e $-\sen x=\cos\bigl(x+\frac\pi2\bigr)$; quindi, per
induzione, $D^{k}\sen x=\sen\bigl(x+k\frac\pi2\bigr)$ e
$D^{k}\cos x=\cos\bigl(x+k\frac\pi2\bigr)$. In $x=0$ le derivate del seno valgono
ciclicamente $0,1,0,-1$, quelle del coseno $1,0,-1,0$; nei polinomi compaiono soltanto
le potenze dispari per il seno, soltanto quelle pari per il coseno:
\begin{align}
  \sen x&=x-\frac{x^{3}}{3!}+\frac{x^{5}}{5!}-\dots+(-1)^{n}\frac{x^{2n+1}}{(2n+1)!}
  +o(x^{2n+2}), \label{eq:taylor-sen}\\
  \cos x&=1-\frac{x^{2}}{2}+\frac{x^{4}}{4!}-\dots+(-1)^{n}\frac{x^{2n}}{(2n)!}
  +o(x^{2n+1}). \label{eq:taylor-cos}
\end{align}
Si noti che il resto del seno è $o(x^{2n+2})$, e non soltanto $o(x^{2n+1})$: il
polinomio di ordine $2n+1$ coincide con quello di ordine $2n+2$, perché il coefficiente
di $x^{2n+2}$ è nullo. Lo stesso vale per il coseno.

\emph{Potenza.} Se $f(x)=(1+x)^{b}$, con $b\in\R$, si ha
$f^{(k)}(x)=b(b-1)\cdots(b-k+1)(1+x)^{b-k}$, e quindi
\begin{equation}\label{eq:taylor-potenza}
  (1+x)^{b}=1+bx+\frac{b(b-1)}{2}x^{2}+\dots+\frac{b(b-1)\cdots(b-n+1)}{n!}x^{n}+o(x^{n}).
\end{equation}
Se $b$ è un intero positivo e $n\ge b$, il resto è nullo e si ritrova la formula del
binomio di Newton (\cref{sec:binomio}). Ad esempio,
$\sqrt{1+x}=1+\frac x2-\frac{x^{2}}8+o(x^{2})$ e
$\frac1{1+x}=1-x+x^{2}-\dots+(-1)^{n}x^{n}+o(x^{n})$.

\begin{python}[Gli sviluppi di Mac Laurin con sympy]
import sympy as sp

x, b = sp.symbols("x b", real=True)
for f in (sp.exp(x), sp.log(1 + x), sp.sin(x), sp.cos(x)):
    print(f"{str(f):>10} =", sp.series(f, x, 0, 6))
print("  (1+x)**b =", sp.series((1 + x)**b, x, 0, 3))
print("  e^(-x^2) =", sp.series(sp.exp(-x**2), x, 0, 6))
\end{python}
\begin{risultato}
    exp(x) = 1 + x + x**2/2 + x**3/6 + x**4/24 + x**5/120 + O(x**6)
log(x + 1) = x - x**2/2 + x**3/3 - x**4/4 + x**5/5 + O(x**6)
    sin(x) = x - x**3/6 + x**5/120 + O(x**6)
    cos(x) = 1 - x**2/2 + x**4/24 + O(x**6)
  (1+x)**b = 1 + b*x + b*x**2*(b - 1)/2 + O(x**3)
  e^(-x^2) = 1 - x**2 + x**4/2 + O(x**6)
\end{risultato}

Il simbolo $O(x^{6})$ usato da \texttt{sympy} indica un resto che, diviso per $x^{6}$,
si mantiene limitato: è un'informazione più precisa di $o(x^{5})$. Per verificare
graficamente quanto ottenuto, la \cref{fig:taylor-seno} riporta i grafici dei polinomi
di Mac Laurin del seno di grado $1,3,5,\dots,23$, cioè
\[
  T_{1}(x)=x, \qquad T_{3}(x)=x-\frac{x^{3}}6, \qquad T_{5}(x)=x-\frac{x^{3}}6+\frac{x^{5}}{120},
  \qquad\dots
\]
Si vede chiaramente che i polinomi di Taylor hanno un grafico tanto più simile a quello
del seno, e su un intervallo tanto più ampio, quanto più il grado è alto. La formula di
Taylor, però, garantisce soltanto il comportamento \emph{vicino a $x_{0}$}: che per il
seno l'approssimazione migliori su ogni intervallo limitato al crescere del grado è un
fatto più profondo, che dimostreremo nel capitolo sulle serie (\cref{teo:sviluppabilita}).

\begin{figure}[htb]
\centering
\includegraphics[width=\textwidth]{taylor_seno}
\caption{Il seno (blu) e i suoi polinomi di Mac Laurin (rame) di grado $1,3,5,\dots,23$;
ciascuno è etichettato con il suo grado nel punto in cui esce dalla finestra
(\texttt{python/fig\_taylor\_seno.py}).}
\label{fig:taylor-seno}
\end{figure}

\subsection{Due applicazioni: limiti e piccoli angoli}

\emph{Calcolo di limiti.} La formula di Taylor è spesso più rapida del teorema di
L'Hôpital, perché mostra in un colpo solo l'ordine di infinitesimo del numeratore e del
denominatore. Ad esempio, per la \eqref{eq:taylor-sen} con $n=1$,
\[
  \lim_{x\to0}\frac{x-\sen x}{x^{3}}=\lim_{x\to0}\frac{x-\bigl(x-\frac{x^{3}}6+o(x^{4})\bigr)}{x^{3}}
  =\lim_{x\to0}\Bigl(\frac16+\frac{o(x^{4})}{x^{3}}\Bigr)=\frac16 ,
\]
mentre con il teorema di L'Hôpital occorrerebbe derivare tre volte. Allo stesso modo
$\frac{1-\cos x}{x^{2}}\to\frac12$, in accordo con la \eqref{eq:lim-1-cos}, e
$\frac{e^{x}-1-x}{x^{2}}\to\frac12$.

\emph{Piccoli angoli.} In meccanica del volo si usano di continuo le approssimazioni
$\sen\varphi\approx\varphi$ e $\cos\varphi\approx1-\frac{\varphi^{2}}2$ per angoli
«piccoli» (l'incidenza, l'angolo di rampa, l'angolo di sbandata). La formula di Taylor
dice quanto piccoli: poiché $\varphi-\sen\varphi=\frac{\varphi^{3}}6+o(\varphi^{4})$,
l'errore relativo che si commette sostituendo $\sen\varphi$ con $\varphi$ è
\[
  \frac{\varphi-\sen\varphi}{\sen\varphi}\approx\frac{\varphi^{2}}{6},
\]
che resta sotto l'$1\%$ finché $\varphi^{2}<0{,}06$, cioè $\varphi\lesssim0{,}245$~rad
$\approx14^{\circ}$. Analogamente, in una virata corretta con angolo di sbandata
$\varphi$ il fattore di carico (che qui indichiamo con $n_{z}$, per non confonderlo con
l'ordine dello sviluppo) è $n_{z}=\frac1{\cos\varphi}$; poiché questa funzione è pari,
ha $f(0)=1$, $f'(0)=0$, $f''(0)=1$ e $f'''(0)=0$, e quindi
$n_{z}=1+\frac{\varphi^{2}}2+o(\varphi^{3})$. Il programma confronta i valori esatti con le
approssimazioni.

\begin{python}[Approssimazioni per piccoli angoli]
import math

print(" phi(gradi)  sen phi    phi    errore   1/cos phi  1+phi^2/2")
for g in (5, 10, 15, 20, 30, 45, 60):
    phi = math.radians(g)
    s = math.sin(phi)
    print(f"{g:8d}    {s:.5f}  {phi:.5f}  {(phi - s) / s:7.2%}"
          f"    {1 / math.cos(phi):.4f}    {1 + phi**2 / 2:.4f}")
\end{python}
\begin{risultato}
 phi(gradi)  sen phi    phi    errore   1/cos phi  1+phi^2/2
       5    0.08716  0.08727    0.13%    1.0038    1.0038
      10    0.17365  0.17453    0.51%    1.0154    1.0152
      15    0.25882  0.26180    1.15%    1.0353    1.0343
      20    0.34202  0.34907    2.06%    1.0642    1.0609
      30    0.50000  0.52360    4.72%    1.1547    1.1371
      45    0.70711  0.78540   11.07%    1.4142    1.3084
      60    0.86603  1.04720   20.92%    2.0000    1.5483
\end{risultato}

A $15^{\circ}$ l'errore relativo è dell'$1{,}15\%$, molto vicino alla stima
$\frac{\varphi^{2}}6\approx1{,}14\%$; a $60^{\circ}$ è ormai del $21\%$. Per il fattore di
carico l'approssimazione quadratica è eccellente fino a $20^{\circ}$ circa, ma a
$60^{\circ}$ dà $1{,}55$ invece di $2$: le virate strette vanno calcolate con la formula
esatta.

\subsection{Il criterio delle derivate successive}

Per mezzo della formula di Taylor si può generalizzare il criterio della derivata
seconda (\cref{teo:test-derivata-seconda}) al caso in cui anche $f''(x_{0})$ sia nulla.

\begin{teorema}[Criterio per i punti di massimo e di minimo]\label{teo:criterio-ordine-n}
\index{criterio!delle derivate successive}
Sia $f$ derivabile $n$ volte nel punto $x_{0}$, con $n\ge2$, e supponiamo che
\begin{equation}\label{eq:criterio-ordine-n}
  f'(x_{0})=f''(x_{0})=\dots=f^{(n-1)}(x_{0})=0, \qquad f^{(n)}(x_{0})\neq0 .
\end{equation}
\begin{enumerate}[label=\textup{(\roman*)}]
  \item Se $n$ è pari e $f^{(n)}(x_{0})>0$, $x_{0}$ è un punto di minimo relativo stretto.
  \item Se $n$ è pari e $f^{(n)}(x_{0})<0$, $x_{0}$ è un punto di massimo relativo stretto.
  \item Se $n$ è dispari, $x_{0}$ non è né di massimo né di minimo relativo.
\end{enumerate}
\end{teorema}

\begin{schema}
\Ipotesi{$f$ derivabile $n$ volte in $x_{0}$; le derivate di ordine $1,\dots,n-1$ sono
nulle in $x_{0}$; $f^{(n)}(x_{0})\neq0$.}
\Tesi{$n$ pari: minimo se $f^{(n)}(x_{0})>0$, massimo se $f^{(n)}(x_{0})<0$; \ $n$ dispari:
né massimo né minimo.}
\end{schema}

\begin{strategia}
Per l'annullarsi delle derivate, la formula di Taylor si riduce a
$f(x)-f(x_{0})=\frac{f^{(n)}(x_{0})}{n!}(x-x_{0})^{n}+o\bigl((x-x_{0})^{n}\bigr)$: vicino a
$x_{0}$ l'incremento $f(x)-f(x_{0})$ ha lo stesso segno del termine principale
$\frac{f^{(n)}(x_{0})}{n!}(x-x_{0})^{n}$. Si tratta allora di discutere il segno di
$(x-x_{0})^{n}$, che è sempre positivo se $n$ è pari, e cambia segno attraversando
$x_{0}$ se $n$ è dispari. Lo strumento per rendere rigoroso «ha lo stesso segno» è la
permanenza del segno.
\end{strategia}

\begin{dimostrazione}
Consideriamo il caso $f^{(n)}(x_{0})>0$ (il caso opposto si ottiene considerando $-f$).
Per le \eqref{eq:criterio-ordine-n} il polinomio di Taylor si riduce a
$T_{n}(x)=f(x_{0})+\frac{f^{(n)}(x_{0})}{n!}(x-x_{0})^{n}$, e per il \cref{teo:taylor}
\begin{equation}\label{eq:criterio-ordine-n-limite}
  \lim_{x\to x_{0}}\frac{f(x)-f(x_{0})}{(x-x_{0})^{n}}
  =\lim_{x\to x_{0}}\Bigl[\frac{f^{(n)}(x_{0})}{n!}+\frac{R_{n}(x)}{(x-x_{0})^{n}}\Bigr]
  =\frac{f^{(n)}(x_{0})}{n!}>0 .
\end{equation}
Per la permanenza del segno (\cref{teo:confronto-funzioni}\,(i)) esiste $\delta>0$ tale
che
\begin{equation}\label{eq:criterio-ordine-n-segno}
  \frac{f(x)-f(x_{0})}{(x-x_{0})^{n}}>0 \qquad\perogni x:\ 0<\abs{x-x_{0}}<\delta .
\end{equation}
Se $n$ è pari, il denominatore $(x-x_{0})^{n}$ è positivo per ogni $x\neq x_{0}$; quindi
$f(x)>f(x_{0})$ per $0<\abs{x-x_{0}}<\delta$, e $x_{0}$ è di minimo relativo stretto. Se
$n$ è dispari, il denominatore è positivo per $x>x_{0}$ e negativo per $x<x_{0}$; quindi
$f(x)>f(x_{0})$ per $x_{0}<x<x_{0}+\delta$ e $f(x)<f(x_{0})$ per $x_{0}-\delta<x<x_{0}$:
in ogni intorno di $x_{0}$ la funzione assume valori sia maggiori sia minori di
$f(x_{0})$, e quindi $x_{0}$ non è né di massimo né di minimo relativo.
\end{dimostrazione}

Il criterio si può riassumere così: \emph{si cerca la prima derivata non nulla in
$x_{0}$; se il suo ordine è pari, $x_{0}$ è di estremo (minimo se è positiva, massimo se
è negativa); se è dispari, non lo è}. Per $f(x)=x^{4}$ risulta
$f'(0)=f''(0)=f'''(0)=0$ e $f^{(4)}(0)=4!=24>0$: l'origine è di minimo, come si verifica
anche direttamente ($x^{4}\ge0=f(0)$). Per $g(x)=x^{3}$ si ha $g'(0)=g''(0)=0$ e
$g'''(0)=3!=6\neq0$, con ordine dispari: l'origine non è di estremo; in effetti è un
punto di flesso (\cref{fig:x3-iperbole}). Quest'ultima proprietà vale in generale.

\begin{proposizione}[Flessi e derivate successive]\label{prop:flesso-ordine-n}
Sia $f$ derivabile $n$ volte in $x_{0}$, con $n\ge3$ dispari, e supponiamo che
$f''(x_{0})=\dots=f^{(n-1)}(x_{0})=0$ e $f^{(n)}(x_{0})\neq0$. Allora $x_{0}$ è un punto di
flesso.
\end{proposizione}

\begin{schema}
\Ipotesi{$f$ derivabile $n$ volte in $x_{0}$, $n\ge3$ dispari; derivate di ordine
$2,\dots,n-1$ nulle in $x_{0}$; $f^{(n)}(x_{0})\neq0$.}
\Tesi{$x_{0}$ è un punto di flesso.}
\end{schema}

\begin{strategia}
Si ripete il ragionamento del teorema precedente per la funzione $f''$ invece che per
$f$: la formula di Taylor mostra che $f''$ ha, vicino a $x_{0}$, il segno di
$f^{(n)}(x_{0})(x-x_{0})^{n-2}$, che cambia attraversando $x_{0}$ perché $n-2$ è dispari.
Il criterio di convessità traduce il cambiamento di segno di $f''$ in un cambiamento di
concavità.
\end{strategia}

\begin{dimostrazione}
La funzione $f''$ è derivabile $n-2\ge1$ volte in $x_{0}$, con derivate nulle fino
all'ordine $n-3$ e derivata di ordine $n-2$ uguale a $f^{(n)}(x_{0})\neq0$. Ragionando
come nella dimostrazione precedente, con la formula di Taylor applicata a $f''$, si
ottiene $\frac{f''(x)}{(x-x_{0})^{n-2}}\to\frac{f^{(n)}(x_{0})}{(n-2)!}\neq0$, e quindi, per
la permanenza del segno, $f''(x)$ ha in un intorno bucato $0<\abs{x-x_{0}}<\delta$ il
segno di $f^{(n)}(x_{0})(x-x_{0})^{n-2}$. Poiché $n-2$ è dispari, $f''$ ha segni opposti
in $(x_{0}-\delta,x_{0})$ e in $(x_{0},x_{0}+\delta)$. Per il criterio di convessità
(\cref{teo:criterio-convessita}), applicato agli intervalli $[x_{0}-\delta',x_{0}]$ e
$[x_{0},x_{0}+\delta']$ con $0<\delta'<\delta$, la funzione è convessa da una parte e
concava dall'altra: $x_{0}$ è di flesso.
\end{dimostrazione}

Si osservi che la proposizione non richiede $f'(x_{0})=0$. Per $\sen x$ nel punto
$x_{0}=\pi$ si ha $f''(\pi)=-\sen\pi=0$ e $f'''(\pi)=-\cos\pi=1\neq0$: con $n=3$ si
ritrova il flesso a tangente obliqua dell'\cref{es:grafico-seno}.

\begin{notastorica}[Taylor, Mac Laurin, Peano]
Brook Taylor (1685--1731), seguace di Newton e segretario della Royal Society, pubblicò
lo sviluppo che porta il suo nome nel trattato \emph{Methodus incrementorum directa et
inversa} (1715), ricavandolo come limite di una formula alle differenze finite; il
risultato era però già noto a James Gregory, che lo aveva comunicato per lettera nel
1671, e in parte a Newton e a Leibniz. Colin Mac Laurin ne fece largo uso, nel caso
$x_{0}=0$, nel \emph{Treatise of Fluxions} (1742). Il significato dello sviluppo come
approssimazione, con una valutazione esplicita del resto, si deve a Lagrange (1797), che
scrisse il resto in una forma simile alla \eqref{eq:accrescimenti-finiti}. La forma del
resto che abbiamo dimostrato, $R_{n}(x)=o\bigl((x-x_{0})^{n}\bigr)$, che richiede la sola
esistenza delle derivate nel punto, è legata al nome di Giuseppe Peano, che la presentò
nel 1889 nella nota \emph{Une nouvelle forme du reste dans la formule de Taylor},
pubblicata sulla rivista \emph{Mathesis}. Peano si era già occupato a fondo dei
fondamenti del calcolo differenziale nel trattato \emph{Calcolo differenziale e principii
di calcolo integrale} (1884), pubblicato come rielaborazione delle lezioni di Angelo
Genocchi.
\end{notastorica}

%% ----------------------------------------------------------------------------
%%  §6.8  Appendice: il teorema di Cauchy e il teorema di L'Hôpital
%% ----------------------------------------------------------------------------
\section{Appendice: il teorema di Cauchy e il teorema di L'Hôpital nel caso generale}
\label{sec:cauchy-hopital}

Allo scopo di dimostrare il teorema di L'Hôpital nel caso generale, è utile una
generalizzazione del teorema di Lagrange che coinvolge \emph{due} funzioni.

\needspace{9\baselineskip}
\begin{teorema}[di Cauchy]\label{teo:cauchy}\index{teorema!di Cauchy}
Siano $f$ e $g$ due funzioni continue in $[a,b]$ e derivabili in $(a,b)$. Se
$g'(x)\neq0$ per ogni $x\in(a,b)$, allora $g(b)\neq g(a)$ ed esiste un punto
$x_{0}\in(a,b)$ tale che
\begin{equation}\label{eq:cauchy-valor-medio}
  \frac{f'(x_{0})}{g'(x_{0})}=\frac{f(b)-f(a)}{g(b)-g(a)} .
\end{equation}
\end{teorema}

\begin{schema}
\Ipotesi{$f,g$ continue in $[a,b]$, derivabili in $(a,b)$; \ $g'(x)\neq0$ in $(a,b)$.}
\Tesi{$g(b)\neq g(a)$ ed esiste $x_{0}\in(a,b)$ con
$\dfrac{f'(x_{0})}{g'(x_{0})}=\dfrac{f(b)-f(a)}{g(b)-g(a)}$.}
\end{schema}

\begin{strategia}
Si procede come per il teorema di Lagrange, costruendo una funzione ausiliaria alla
quale applicare il teorema di Rolle: al posto della corda $f(a)+m(x-a)$ si sottrae a $f$
l'espressione $f(a)+m\bigl(g(x)-g(a)\bigr)$, con la costante $m$ scelta in modo che la
differenza si annulli in entrambi gli estremi. Prima, però, bisogna assicurarsi che il
denominatore $g(b)-g(a)$ non sia nullo: anche questo è una conseguenza del teorema di
Rolle, applicato a $g$.
\end{strategia}

\begin{dimostrazione}
Se fosse $g(a)=g(b)$, per il teorema di Rolle (\cref{teo:rolle}) applicato a $g$
esisterebbe un punto di $(a,b)$ in cui $g'$ si annulla, contro l'ipotesi. Dunque
$g(b)-g(a)\neq0$, e la funzione
\begin{equation}\label{eq:cauchy-h}
  h(x)=f(x)-\Bigl[f(a)+\frac{f(b)-f(a)}{g(b)-g(a)}\bigl(g(x)-g(a)\bigr)\Bigr],
  \qquad x\in[a,b],
\end{equation}
è ben definita, continua in $[a,b]$ e derivabile in $(a,b)$, con
\begin{equation}\label{eq:cauchy-h-derivata}
  h'(x)=f'(x)-\frac{f(b)-f(a)}{g(b)-g(a)}\,g'(x) .
\end{equation}
Inoltre $h(a)=0$ e $h(b)=f(b)-\bigl[f(a)+f(b)-f(a)\bigr]=0$. Per il teorema di Rolle
esiste $x_{0}\in(a,b)$ con $h'(x_{0})=0$, che, dividendo la
\eqref{eq:cauchy-h-derivata} per $g'(x_{0})\neq0$, equivale alla \eqref{eq:cauchy-valor-medio}.
\end{dimostrazione}

\begin{attenzione}[Perché non basta applicare due volte il teorema di Lagrange]
Si potrebbe pensare di ottenere la \eqref{eq:cauchy-valor-medio} dividendo membro a membro le
uguaglianze $f(b)-f(a)=f'(x_{1})(b-a)$ e $g(b)-g(a)=g'(x_{2})(b-a)$ fornite dal teorema di
Lagrange. Si otterrebbe però $\frac{f(b)-f(a)}{g(b)-g(a)}=\frac{f'(x_{1})}{g'(x_{2})}$, con
due punti $x_{1}$, $x_{2}$ \emph{in generale diversi}; il teorema di Cauchy afferma
invece che si può usare lo \emph{stesso} punto a numeratore e a denominatore, ed è
proprio questo che servirà nella dimostrazione del teorema di L'Hôpital. Un'interpretazione
geometrica aiuta: se si pensa al punto $\bigl(g(t),f(t)\bigr)$ che descrive una curva nel
piano al variare di $t\in[a,b]$, il secondo membro della \eqref{eq:cauchy-valor-medio} è la pendenza
della corda che congiunge gli estremi della curva, e il primo è la pendenza della
tangente alla curva nel punto di parametro $x_{0}$: il teorema di Cauchy è il teorema di
Lagrange per le curve in forma parametrica.
\end{attenzione}

Se nel teorema di Cauchy si sceglie $g(x)=x$, si ritrova il teorema di Lagrange; se nel
teorema di Lagrange si suppone $f(a)=f(b)$, si ritrova il teorema di Rolle, che a sua
volta è alla base delle dimostrazioni dei teoremi di Lagrange e di Cauchy. I tre teoremi
sono quindi da considerarsi fra loro equivalenti: ciascuno è un caso particolare del
successivo, e ciascuno si deduce dal primo.

\begin{osservazione}[Il teorema di Cauchy senza quozienti]\label{oss:cauchy-simmetrica}
Se $f$ e $g$ sono continue in $[a,b]$ e derivabili in $(a,b)$, \emph{senza alcuna
ipotesi su $g'$}, esiste comunque un punto $x_{0}\in(a,b)$ tale che
\begin{equation}\label{eq:cauchy-simmetrica}
  \bigl[f(b)-f(a)\bigr]\,g'(x_{0})=\bigl[g(b)-g(a)\bigr]\,f'(x_{0}) .
\end{equation}
Basta applicare il teorema di Rolle alla funzione
$k(x)=f(x)\bigl[g(b)-g(a)\bigr]-g(x)\bigl[f(b)-f(a)\bigr]$, come fa Giusti
\cite{giusti}: un calcolo diretto dà $k(a)=k(b)=f(a)g(b)-f(b)g(a)$, e la condizione
$k'(x_{0})=0$ è la \eqref{eq:cauchy-simmetrica}. In questa forma simmetrica non si divide
per nulla, e con $g(x)=x$ si ottiene subito il teorema di Lagrange. L'ipotesi $g'\neq0$
serve soltanto per passare alla forma con i quozienti \eqref{eq:cauchy-valor-medio}:
garantisce che $g(b)\neq g(a)$ e permette di dividere per $g'(x_{0})$.
\end{osservazione}

\subsection{Un lemma sul segno della derivata}

Nella dimostrazione del teorema di L'Hôpital per le forme $\frac\infty\infty$ ci servirà
sapere che una derivata che non si annulla mai in un intervallo ha segno costante. Per
una derivata \emph{continua} è una conseguenza del teorema degli zeri; ma la derivata di
una funzione derivabile può non essere continua (\cref{oss:f2-non-C1}). Il fatto vale
ugualmente, per un'argomentazione elegante dovuta a Darboux.

\begin{lemma}[di Darboux]\label{lem:darboux}\index{teorema!di Darboux}
Sia $g$ derivabile in un intervallo $J$, con $g'(x)\neq0$ per ogni $x\in J$. Allora
$g'$ ha segno costante in $J$: è positiva in tutto $J$, oppure negativa in tutto $J$.
\end{lemma}

\begin{schema}
\Ipotesi{$g$ derivabile nell'intervallo $J$; \ $g'(x)\neq0$ per ogni $x\in J$.}
\Tesi{$g'>0$ in $J$, oppure $g'<0$ in $J$.}
\end{schema}

\begin{strategia}
Ragioniamo per assurdo: se $g'(x_{1})<0<g'(x_{2})$, la funzione scende vicino a $x_{1}$ e
sale vicino a $x_{2}$. Il suo minimo nell'intervallo di estremi $x_{1}$, $x_{2}$, che
esiste per il teorema di Weierstrass, non può allora cadere negli estremi, e il teorema
di Fermat lo costringe a essere un punto in cui la derivata si annulla. Lo strumento per
escludere gli estremi è la \cref{prop:fermat-estremi}.
\end{strategia}

\begin{dimostrazione}
Supponiamo per assurdo che esistano $x_{1},x_{2}\in J$ con $g'(x_{1})<0<g'(x_{2})$, e
supponiamo $x_{1}<x_{2}$ (l'altro caso è analogo). La funzione $g$ è continua in
$[x_{1},x_{2}]$ e quindi, per il teorema di Weierstrass, ha un punto di minimo assoluto
$\xi\in[x_{1},x_{2}]$. Se fosse $\xi=x_{1}$, l'estremo sinistro sarebbe di minimo, e per
la \cref{prop:fermat-estremi} si avrebbe $g'(x_{1})\ge0$: falso. Se fosse $\xi=x_{2}$,
l'estremo destro sarebbe di minimo, e si avrebbe $g'(x_{2})\le0$: falso. Dunque
$\xi\in(x_{1},x_{2})$, e per il teorema di Fermat $g'(\xi)=0$, contro l'ipotesi.
\end{dimostrazione}

Più in generale, lo stesso ragionamento prova che la derivata di una funzione
derivabile in un intervallo assume tutti i valori compresi tra due suoi valori, anche
se non è continua (si veda l'\cref{es:darboux}).

Ci servirà anche una piccola estensione del \cref{teo:limite-monotone-funzioni}, che era
enunciato per le funzioni monotòne in un intervallo chiuso e limitato e forniva limiti
finiti.

\begin{osservazione}[Limiti delle funzioni monotòne in un intervallo aperto]
\label{oss:monotone-aperto}
Sia $h$ decrescente in un intervallo $(c,d)$ e sia $\lambda=\inf\{h(x): c<x<d\}$, con
$\lambda=-\infty$ se $h$ non è limitata inferiormente. Allora
$\lim_{x\to d^{-}}h(x)=\lambda$. Infatti, se $\lambda\in\R$, per ogni $\varepsilon>0$ esiste
$\bar x\in(c,d)$ con $h(\bar x)<\lambda+\varepsilon$, e per $\bar x<x<d$ si ha
$\lambda\le h(x)\le h(\bar x)<\lambda+\varepsilon$; se $\lambda=-\infty$, per ogni $M$ esiste
$\bar x$ con $h(\bar x)<M$, e $h(x)\le h(\bar x)<M$ per $\bar x<x<d$. Analogamente, una
funzione crescente in $(c,d)$ tende, per $x\to d^{-}$, al suo estremo superiore, finito
oppure $+\infty$.
\end{osservazione}

\subsection{La dimostrazione del teorema di L'Hôpital}

Siamo ora in grado di dimostrare il teorema di L'Hôpital (\cref{teo:hopital}) nelle
ipotesi generali in cui lo abbiamo enunciato.

\begin{strategia}
Nel caso $\frac00$ estendiamo $f$ e $g$ per continuità ponendo $f(x_{0})=g(x_{0})=0$.
Allora $\frac{f(x)}{g(x)}=\frac{f(x)-f(x_{0})}{g(x)-g(x_{0})}$, e il teorema di Cauchy
trasforma questo rapporto in $\frac{f'(\xi)}{g'(\xi)}$, con $\xi$ compreso tra $x_{0}$ e
$x$: quando $x\to x_{0}$ anche $\xi\to x_{0}$, e il rapporto delle derivate tende a
$\ell$. Il linguaggio delle successioni rende il ragionamento rigoroso. Nel caso
$\frac\infty\infty$ questo trucco non funziona, perché $f$ e $g$ non si possono
estendere a $x_{0}$; useremo il teorema di Lagrange in forma di monotonia: dalla
disuguaglianza $\frac{f'}{g'}<\ell+\varepsilon$ ricaveremo che una opportuna combinazione
$f-c\,g$ è decrescente, e quindi tende a $-\infty$ o a un limite finito. Infine il caso
$x\to\pm\infty$ si riconduce ai precedenti con il cambio di variabile $t=\frac1x$.
\end{strategia}

\begin{dimostrazione}
\emph{Passo 1: forma $\frac00$ per $x\to x_{0}\in\R$.} Siano $f$, $g$ derivabili in
$I\setminus\{x_{0}\}$, con $g'\neq0$, $f(x)\to0$, $g(x)\to0$ e
$\frac{f'(x)}{g'(x)}\to\ell\in\overline\R$. Ponendo $f(x_{0})=g(x_{0})=0$ (con abuso di
notazione usiamo lo stesso simbolo per le funzioni estese), $f$ e $g$ diventano
continue in tutto $I$: in $x_{0}$ perché il limite coincide con il valore, negli altri
punti perché derivabili. Osserviamo che $g(x)\neq0$ per ogni $x\in I\setminus\{x_{0}\}$:
se fosse $g(x_{1})=0$ con $x_{1}\neq x_{0}$, per il teorema di Rolle applicato a $g$
nell'intervallo di estremi $x_{0}$, $x_{1}$ esisterebbe un punto in cui $g'$ si annulla.

Sia ora $x_{n}$ una qualunque successione di punti di $I\setminus\{x_{0}\}$ convergente
a $x_{0}$. Per il teorema di Cauchy, applicato nell'intervallo di estremi $x_{0}$ e
$x_{n}$ (nel quale $f$, $g$ sono continue, e derivabili nei punti interni con $g'\neq0$),
per ogni $n$ esiste un punto $x_{n}'$ strettamente compreso tra $x_{0}$ e $x_{n}$ tale che
\begin{equation}\label{eq:hopital-cauchy}
  \frac{f(x_{n})}{g(x_{n})}=\frac{f(x_{n})-f(x_{0})}{g(x_{n})-g(x_{0})}=\frac{f'(x_{n}')}{g'(x_{n}')} .
\end{equation}
Poiché $0<\abs{x_{n}'-x_{0}}<\abs{x_{n}-x_{0}}$, la successione $x_{n}'$ converge a $x_{0}$
(teorema dei carabinieri) ed è formata da punti diversi da $x_{0}$. Per la definizione di
limite di funzione mediante le successioni (\cref{def:limite-funzione}, con le varianti
per i limiti infiniti, \cref{oss:ponte-infinito}),
$\frac{f'(x_{n}')}{g'(x_{n}')}\to\ell$, e quindi, per la \eqref{eq:hopital-cauchy},
$\frac{f(x_{n})}{g(x_{n})}\to\ell$. Poiché la successione $x_{n}\to x_{0}$ è arbitraria,
ancora per la \cref{def:limite-funzione} si conclude che
\[
  \lim_{x\to x_{0}}\frac{f(x)}{g(x)}=\ell=\lim_{x\to x_{0}}\frac{f'(x)}{g'(x)} .
\]

\emph{Passo 2: limiti unilaterali.} Per il limite destro si suppone che $f$, $g$ siano
derivabili in un intervallo $(x_{0},x_{0}+r)$ e si considerano soltanto successioni
$x_{n}\to x_{0}$ con $x_{n}>x_{0}$; il ragionamento del passo 1 si ripete parola per
parola. Lo stesso vale per il limite sinistro.

\emph{Passo 3: forma $\frac\infty\infty$.} Consideriamo il limite sinistro per
$x\to x_{0}^{-}$ (il limite destro è analogo, e combinando i due si ottiene il limite
completo, \cref{prop:dx-sx}). Siano $f$, $g$ derivabili in $(x_{0}-r,x_{0})$, con
$g'\neq0$, $\abs{g(x)}\to+\infty$ e $\frac{f'(x)}{g'(x)}\to\ell$. Per il
\cref{lem:darboux} $g'$ ha segno costante in $(x_{0}-r,x_{0})$, e sostituendo
eventualmente $f$, $g$ con $-f$, $-g$ (il che non cambia i rapporti $\frac fg$ e
$\frac{f'}{g'}$, né l'ipotesi $\abs g\to+\infty$) possiamo supporre $g'>0$. Allora $g$ è
strettamente crescente (\cref{cor:derivata-positiva}); fissato $\bar x\in(x_{0}-r,x_{0})$,
si ha $g(x)\ge g(\bar x)$ per $\bar x<x<x_{0}$, cioè $g$ è limitata inferiormente vicino a
$x_{0}$, e quindi l'ipotesi $\abs{g(x)}\to+\infty$ impone $g(x)\to+\infty$. In
particolare $g(x)>0$ vicino a $x_{0}$.

Proviamo anzitutto la seguente affermazione: \emph{se $c<c'$ sono due numeri reali e
$\frac{f'(x)}{g'(x)}<c$ in un intervallo $(x_{0}-\delta_{1},x_{0})$, allora
$\frac{f(x)}{g(x)}<c'$ in un intervallo $(x_{0}-\delta,x_{0})$.} Infatti, poiché $g'>0$,
le funzioni
\[
  h_{1}(x)=f(x)-c\,g(x), \qquad h_{2}(x)=f(x)-c'\,g(x)
\]
hanno derivate $h_{1}'=g'\bigl(\frac{f'}{g'}-c\bigr)<0$ e
$h_{2}'=h_{1}'-(c'-c)g'<0$ in $(x_{0}-\delta_{1},x_{0})$; sono quindi strettamente
decrescenti, e per l'\cref{oss:monotone-aperto} ammettono limite per
$x\to x_{0}^{-}$, uguale al loro estremo inferiore: un numero reale oppure $-\infty$. I due
limiti non possono essere entrambi finiti, perché la differenza
$h_{1}(x)-h_{2}(x)=(c'-c)g(x)$ tende a $+\infty$. Dunque almeno una delle due funzioni,
diciamo $h_{i}$, tende a $-\infty$, e per la permanenza del segno esiste $\delta_{2}>0$
tale che $h_{i}(x)<0$ in $(x_{0}-\delta_{2},x_{0})$. Sia $\delta\le\delta_{2}$ tale che
$g>0$ in $(x_{0}-\delta,x_{0})$. Dividendo $h_{i}(x)<0$ per $g(x)>0$ si ottiene
$\frac{f(x)}{g(x)}<c$ se $i=1$, oppure $\frac{f(x)}{g(x)}<c'$ se $i=2$; in entrambi i casi
$\frac{f(x)}{g(x)}<c'$ in $(x_{0}-\delta,x_{0})$, come volevamo. In modo del tutto
simmetrico (con funzioni crescenti che tendono a $+\infty$) si prova che \emph{se
$c''<c$ e $\frac{f'}{g'}>c$ vicino a $x_{0}$, allora $\frac fg>c''$ vicino a $x_{0}$.}

Ora concludiamo. Se $\ell\in\R$, fissato $\varepsilon>0$ si ha
$\ell-\varepsilon<\frac{f'}{g'}<\ell+\varepsilon$ vicino a $x_{0}$; applicando le due
affermazioni con $c=\ell+\varepsilon$, $c'=\ell+2\varepsilon$ e con $c=\ell-\varepsilon$,
$c''=\ell-2\varepsilon$, otteniamo $\ell-2\varepsilon<\frac{f(x)}{g(x)}<\ell+2\varepsilon$ in un
intorno sinistro di $x_{0}$. Per l'arbitrarietà di $\varepsilon$ (\cref{teo:eps-delta}),
$\frac fg\to\ell$. Se $\ell=+\infty$, per ogni $M$ si ha $\frac{f'}{g'}>M+1$ vicino a
$x_{0}$, e la seconda affermazione dà $\frac fg>M$ vicino a $x_{0}$: quindi
$\frac fg\to+\infty$. Il caso $\ell=-\infty$ è analogo.

\emph{Passo 4: limiti per $x\to\pm\infty$.} Siano $f$, $g$ derivabili in una semiretta
$[a,+\infty)$, con $a>0$ e $g'(x)\neq0$. Definiamo nell'intervallo $(0,\frac1a]$ le
funzioni
\[
  \tilde f(t)=f\Bigl(\frac1t\Bigr), \qquad \tilde g(t)=g\Bigl(\frac1t\Bigr)
\]
(se $t\in(0,\frac1a]$, allora $x=\frac1t\in[a,+\infty)$). Per il teorema sul limite delle
funzioni composte (\cref{teo:limite-composta}), $\tilde f$ e $\tilde g$ hanno per
$t\to0^{+}$ gli stessi limiti che $f$ e $g$ hanno per $x\to+\infty$: sono quindi
infinitesime, oppure $\abs{\tilde g}\to+\infty$. Per la regola di derivazione delle
funzioni composte (\cref{teo:derivata-composta}),
\[
  \frac{\tilde f'(t)}{\tilde g'(t)}=\frac{f'(1/t)\cdot(-1/t^{2})}{g'(1/t)\cdot(-1/t^{2})}
  =\frac{f'(1/t)}{g'(1/t)}\ \longrightarrow\ \ell \qquad (t\to0^{+}),
\]
dove il limite segue dall'ipotesi $\frac{f'(x)}{g'(x)}\to\ell$ per $x\to+\infty$ e, di
nuovo, dal teorema sul limite delle funzioni composte; inoltre $\tilde g'(t)\neq0$. Per i passi 2 e 3, applicati al limite destro per $t\to0^{+}$,
\[
  \lim_{x\to+\infty}\frac{f(x)}{g(x)}=\lim_{t\to0^{+}}\frac{\tilde f(t)}{\tilde g(t)}
  =\lim_{t\to0^{+}}\frac{\tilde f'(t)}{\tilde g'(t)}=\ell ,
\]
dove nella prima uguaglianza si è usato ancora il teorema sul limite delle funzioni
composte. Il caso $x\to-\infty$ si tratta allo stesso modo, con $t\to0^{-}$.
\end{dimostrazione}

\begin{strategia}[Un'altra via per il passo 3]
Il passo 3 si può dimostrare con lo stesso strumento del passo 1, il teorema di Cauchy,
senza il lemma di Darboux. Poiché $f$ e $g$ non sono definite in $x_{0}$, lo si applica
in $[y,x]$, con $y$ fisso e tanto vicino a $x_{0}$ che, alla sua destra, $\frac{f'}{g'}$
disti da $\ell$ meno di $\varepsilon$. Il quoziente $\frac{f(x)-f(y)}{g(x)-g(y)}$ è allora vicino a $\ell$, e
differisce da $\frac{f(x)}{g(x)}$ per termini come $\frac{f(y)}{g(x)}$ e
$\frac{g(y)}{g(x)}$, che tendono a zero perché $y$ è fisso e $\abs{g(x)}\to+\infty$. È,
con qualche ritocco, la via seguita da Giusti \cite{giusti}.
\end{strategia}

\begin{dimostrazione}[Seconda dimostrazione del passo 3]
Siano, come nel passo 3, $f$ e $g$ derivabili in $(x_{0}-r,x_{0})$, con $g'\neq0$,
$\abs{g(x)}\to+\infty$ e $\frac{f'(x)}{g'(x)}\to\ell$ per $x\to x_{0}^{-}$. Se
$x_{0}-r<y<x<x_{0}$, le funzioni $f$ e $g$ sono continue in $[y,x]$ e derivabili in
$(y,x)$, con $g'\neq0$; per il teorema di Cauchy (\cref{teo:cauchy}) è $g(x)\neq g(y)$,
ed esiste $\xi\in(y,x)$ tale che
\begin{equation}\label{eq:hopital-Q}
  Q(x)=\frac{f(x)-f(y)}{g(x)-g(y)}=\frac{f'(\xi)}{g'(\xi)} .
\end{equation}
Se inoltre $g(x)\neq0$, vale l'identità
\begin{equation}\label{eq:hopital-identita}
  \frac{f(x)}{g(x)}=\frac{f(y)}{g(x)}+\Bigl(1-\frac{g(y)}{g(x)}\Bigr)Q(x),
\end{equation}
che si verifica osservando che
\[
  \Bigl(1-\frac{g(y)}{g(x)}\Bigr)Q(x)=\frac{g(x)-g(y)}{g(x)}\cdot
  \frac{f(x)-f(y)}{g(x)-g(y)}=\frac{f(x)-f(y)}{g(x)} .
\]

\emph{Caso $\ell\in\R$.} Fissato $\varepsilon>0$, esiste $\delta_{1}\in(0,r)$ tale che
$\bigl|\frac{f'(\xi)}{g'(\xi)}-\ell\bigr|<\varepsilon$ per $x_{0}-\delta_{1}<\xi<x_{0}$.
Scegliamo $y=x_{0}-\delta_{1}$ e teniamolo fisso. Per $y<x<x_{0}$ il punto $\xi$ della
\eqref{eq:hopital-Q} sta in $(x_{0}-\delta_{1},x_{0})$, e quindi
$\abs{Q(x)-\ell}<\varepsilon$; in particolare $\abs{Q(x)}<\abs\ell+\varepsilon$. Dalla
\eqref{eq:hopital-identita}, nei punti $x\in(y,x_{0})$ in cui $g(x)\neq0$,
\[
  \frac{f(x)}{g(x)}-\ell=\bigl(Q(x)-\ell\bigr)+\frac{f(y)-g(y)\,Q(x)}{g(x)},
\]
e quindi, per la disuguaglianza triangolare,
\[
  \Bigl|\frac{f(x)}{g(x)}-\ell\Bigr|<\varepsilon+\frac{C}{\abs{g(x)}}, \qquad
  C=\abs{f(y)}+\abs{g(y)}\bigl(\abs\ell+\varepsilon\bigr).
\]
La costante $C\ge0$ non dipende da $x$, perché $y$ è fisso. Poiché
$\abs{g(x)}\to+\infty$, esiste $\delta\in(0,\delta_{1})$ tale che
$\abs{g(x)}>\frac{C}{\varepsilon}$, e in particolare $g(x)\neq0$, per
$x_{0}-\delta<x<x_{0}$. Per tali $x$ risulta $\frac{C}{\abs{g(x)}}<\varepsilon$, e quindi
$\bigl|\frac{f(x)}{g(x)}-\ell\bigr|<2\varepsilon$. Per l'arbitrarietà di $\varepsilon$
(\cref{teo:eps-delta}), $\frac fg\to\ell$.

\emph{Caso $\ell=+\infty$.} Fissato $M>0$, esiste $\delta_{1}\in(0,r)$ tale che
$\frac{f'(\xi)}{g'(\xi)}>2(M+1)$ per $x_{0}-\delta_{1}<\xi<x_{0}$; con $y=x_{0}-\delta_{1}$
si ha allora $Q(x)>2(M+1)>0$ per $y<x<x_{0}$. Poiché $\abs{g(x)}\to+\infty$, esiste
$\delta\in(0,\delta_{1})$ tale che $\abs{g(x)}>\max\bigl\{\abs{f(y)},2\abs{g(y)}\bigr\}$
per $x_{0}-\delta<x<x_{0}$; per tali $x$ è $g(x)\neq0$,
$\bigl|\frac{f(y)}{g(x)}\bigr|<1$ e $1-\frac{g(y)}{g(x)}>\frac12$. Dalla
\eqref{eq:hopital-identita}, essendo $Q(x)>0$,
\[
  \frac{f(x)}{g(x)}>-1+\frac12\,Q(x)>-1+(M+1)=M .
\]
Dunque $\frac fg\to+\infty$. Il caso $\ell=-\infty$ si riconduce a questo sostituendo $f$
con $-f$.
\end{dimostrazione}

\begin{osservazione}[Le due vie a confronto]\label{oss:hopital-due-vie}
Con la seconda dimostrazione tutto il teorema di L'Hôpital poggia sul solo teorema di
Cauchy: non occorre sapere che $g'$ ha segno costante né che $g$ tende a $+\infty$
anziché a $-\infty$. Il suo ingrediente è un ragionamento «in due tempi»: prima si fissa
un punto ausiliario $y$, scelto in funzione di $\varepsilon$, poi si fa tendere $x$ a
$x_{0}$ con $y$ fermo, così che i termini che contengono $y$ diventino trascurabili. La
prima dimostrazione traduce invece le disuguaglianze tra le derivate in monotonia delle
funzioni $f-c\,g$, ed evita l'identità \eqref{eq:hopital-identita}, che a prima vista
sembra piovuta dal cielo.
\end{osservazione}

\begin{notastorica}[Cauchy e Darboux]
Il teorema che porta il nome di Augustin-Louis Cauchy compare nelle lezioni di calcolo
differenziale che egli tenne all'École Polytechnique e pubblicò negli anni Venti
dell'Ottocento, come strumento per studiare proprio le forme indeterminate: è con esso
che Cauchy diede una dimostrazione rigorosa del teorema di L'Hôpital, oltre un secolo
dopo la sua pubblicazione. Il lemma sul segno
della derivata è un caso particolare di un risultato di Gaston Darboux, che nel 1875,
nella \emph{Mémoire sur les fonctions discontinues}, dimostrò che ogni derivata gode
della proprietà dei valori intermedi, anche quando è discontinua; nella stessa memoria
Darboux costruì funzioni continue prive di derivata e funzioni derivabili con derivata
discontinua, contribuendo a sgombrare il campo dalle intuizioni geometriche ingannevoli
che avevano accompagnato la nascita del calcolo.
\end{notastorica}

%% ----------------------------------------------------------------------------
%%  §6.9  Esercizi di riepilogo
%% ----------------------------------------------------------------------------
\section{Esercizi di riepilogo}
\label{sec:esercizi-cap6}

Come nei capitoli precedenti, ogni soluzione è accompagnata da una nota sulla
\emph{scelta logica} che la guida.

\needspace{8\baselineskip}
\subsection*{I teoremi di Rolle, di Lagrange e di Darboux}

\begin{esercizio}\label{es:cubica-unica}
Dimostrare che l'equazione $x^{3}+x-1=0$ \eqref{eq:cubica} ha una e una sola soluzione
reale, e che essa appartiene all'intervallo $(0,1)$.
\begin{soluzione}
Sia $f(x)=x^{3}+x-1$. \emph{Esistenza:} $f$ è continua, $f(0)=-1<0$ e $f(1)=1>0$; per il
teorema degli zeri (\cref{teo:zeri}) esiste una soluzione in $(0,1)$. \emph{Unicità:} se
ci fossero due soluzioni $x_{1}<x_{2}$, sarebbe $f(x_{1})=f(x_{2})=0$ e, per il teorema di
Rolle, $f'(\xi)=3\xi^{2}+1=0$ per qualche $\xi$, il che è impossibile perché
$3\xi^{2}+1\ge1$. In alternativa: $f'>0$ su $\R$, quindi $f$ è strettamente crescente
(\cref{cor:derivata-positiva}) e iniettiva. \emph{Scelta logica:} esistenza e unicità
richiedono strumenti diversi; il teorema degli zeri non dice nulla sul numero delle
soluzioni, e il teorema di Rolle (o la monotonia) lo controlla: tra due zeri di $f$ c'è
uno zero di $f'$.
\end{soluzione}
\end{esercizio}

\begin{esercizio}\label{es:log-lagrange}
Dimostrare con il teorema di Lagrange che, per ogni $x>0$,
$\dfrac{x}{1+x}<\log(1+x)<x$.
\begin{soluzione}
Applichiamo il teorema di Lagrange a $\log t$ nell'intervallo $[1,1+x]$: esiste
$\xi\in(1,1+x)$ tale che $\log(1+x)-\log1=\frac1\xi\,x$. Poiché $1<\xi<1+x$, risulta
$\frac1{1+x}<\frac1\xi<1$, e moltiplicando per $x>0$ si ottiene la tesi. La
disuguaglianza di destra era già nell'\cref{es:disuguaglianze-exp-log}.
\emph{Scelta logica:} il teorema di Lagrange trasforma un'uguaglianza esatta ma con un
punto sconosciuto in una doppia disuguaglianza, sostituendo al punto sconosciuto gli
estremi dell'intervallo in cui si trova; è il meccanismo dell'\cref{es:stima-radice}.
\end{soluzione}
\end{esercizio}

\begin{esercizio}\label{es:darboux}
(\emph{Teorema di Darboux}.) Sia $g$ derivabile in $[a,b]$ e sia $\lambda$ un numero
strettamente compreso tra $g'(a)$ e $g'(b)$. Dimostrare che esiste $\xi\in(a,b)$ con
$g'(\xi)=\lambda$. Dedurne che la funzione che vale $0$ per $x<0$ e $1$ per $x\ge0$ non è
la derivata di alcuna funzione in $[-1,1]$.
\begin{soluzione}
La funzione $\varphi(x)=g(x)-\lambda x$ è derivabile in $[a,b]$, con
$\varphi'(x)=g'(x)-\lambda$; per la scelta di $\lambda$, $\varphi'(a)$ e $\varphi'(b)$ hanno
segni opposti. Se $\varphi'$ non si annullasse in alcun punto di $[a,b]$, per il
\cref{lem:darboux} avrebbe segno costante: contraddizione. Dunque $\varphi'(\xi)=0$ per
qualche $\xi$, che non può essere né $a$ né $b$ (dove $\varphi'\neq0$); cioè
$g'(\xi)=\lambda$ con $\xi\in(a,b)$. Una funzione con un salto non assume i valori
intermedi, e quindi non può essere una derivata. \emph{Scelta logica:} si sposta il
problema «$g'$ assume il valore $\lambda$» nel problema «$\varphi'$ si annulla»,
sottraendo a $g$ la retta di pendenza $\lambda$; è lo stesso trucco della dimostrazione
del teorema di Lagrange.
\end{soluzione}
\end{esercizio}

\begin{esercizio}\label{es:corde-convesse}
Sia $f$ convessa in $[a,b]$ nel senso della \cref{def:convessa}. Dimostrare che ogni
corda sta sopra il grafico, cioè che
$f\bigl((1-t)x+ty\bigr)\le(1-t)f(x)+tf(y)$ per ogni $x,y\in[a,b]$ e ogni $t\in[0,1]$.
\begin{soluzione}
Sia $z=(1-t)x+ty\in[a,b]$. Scriviamo la \eqref{eq:convessa} con punto di tangenza $z$,
prima nel punto $x$ e poi nel punto $y$:
$f(x)\ge f(z)+f'(z)(x-z)$, $f(y)\ge f(z)+f'(z)(y-z)$. Moltiplichiamo la prima per
$1-t\ge0$, la seconda per $t\ge0$ e sommiamo:
\[
  (1-t)f(x)+tf(y)\ge f(z)+f'(z)\bigl[(1-t)x+ty-z\bigr]=f(z) .
\]
\emph{Scelta logica:} si sceglie come punto di tangenza proprio il punto $z$ in cui si
vuole stimare $f$; la combinazione con pesi $1-t$ e $t$ è quella che fa sparire il termine
con la derivata, perché $(1-t)x+ty-z=0$.
\end{soluzione}
\end{esercizio}

\needspace{8\baselineskip}
\subsection*{Disuguaglianze e problemi di ottimo}

\begin{esercizio}\label{es:bernoulli-convessita}
Dimostrare che $(1+x)^{b}\ge1+bx$ per ogni $x>-1$ e ogni $b\ge1$, e che la disuguaglianza
si inverte se $0\le b\le1$. Confrontare con la disuguaglianza di Bernoulli
(\cref{teo:bernoulli}).
\begin{soluzione}
La funzione $f(x)=(1+x)^{b}$ è derivabile due volte in $(-1,+\infty)$, con
$f''(x)=b(b-1)(1+x)^{b-2}$, che è $\ge0$ se $b\ge1$ e $\le0$ se $0\le b\le1$. Nel primo
caso $f$ è convessa, e la \eqref{eq:convessa} con $x_{0}=0$ dà
$(1+x)^{b}\ge f(0)+f'(0)x=1+bx$; nel secondo caso $f$ è concava e la disuguaglianza si
inverte. La disuguaglianza di Bernoulli del \cref{cap:numeri}, dimostrata per induzione
per gli esponenti naturali, si estende così a ogni esponente reale $b\ge1$.
\emph{Scelta logica:} $y=1+bx$ è la retta tangente al grafico di $(1+x)^{b}$ nel punto
$x=0$, e la disuguaglianza dice da che parte della tangente sta il grafico.
\end{soluzione}
\end{esercizio}

\begin{esercizio}[La trave di massima resistenza]\label{es:trave}
Da un tronco cilindrico di diametro $d$ si vuole ricavare una trave a sezione
rettangolare, di base $b$ e altezza $h$, inflessa nel piano dell'altezza. La resistenza a
flessione è proporzionale al \emph{modulo di resistenza} $W_{f}=\frac{bh^{2}}6$.
Determinare le dimensioni che rendono massimo $W_{f}$.
\begin{soluzione}
Il rettangolo è inscritto nella circonferenza, quindi $b^{2}+h^{2}=d^{2}$ e
\[
  W_{f}(b)=\frac{b(d^{2}-b^{2})}{6}, \qquad b\in[0,d] .
\]
La funzione è continua nell'intervallo chiuso, nulla negli estremi e positiva
all'interno; per il metodo del \cref{sec:fermat} il massimo assoluto è in un punto
stazionario interno. Da $W_{f}'(b)=\frac{d^{2}-3b^{2}}6=0$ si ricava $b^{*}=\frac d{\sqrt3}$,
e quindi $h^{*}=\sqrt{d^{2}-b^{*2}}=d\sqrt{\frac23}$: il rapporto ottimo è
$\frac{h^{*}}{b^{*}}=\sqrt2$. Inoltre $W_{f}''(b)=-b<0$ conferma il massimo. Se invece
si vuole la trave più \emph{rigida}, si massimizza il momento d'inerzia del
second'ordine $\frac{bh^{3}}{12}$, e si trova $\frac hb=\sqrt3$. Nelle costruzioni
aeronautiche lo stesso tipo di calcolo guida il dimensionamento delle solette dei
longheroni. \emph{Scelta logica:} il vincolo geometrico permette di eliminare una
variabile e di ridurre il problema a una funzione di una sola variabile su un intervallo
chiuso; il teorema di Weierstrass garantisce l'esistenza del massimo, e il teorema di
Fermat lo localizza.
\end{soluzione}
\end{esercizio}

\begin{python}[La trave di massima resistenza con sympy]
import sympy as sp

b, d = sp.symbols("b d", positive=True)
h2 = d**2 - b**2                  # Pitagora: b^2 + h^2 = d^2
Wf = b * h2 / 6                   # modulo di resistenza b h^2 / 6
bs = sp.solve(sp.diff(Wf, b), b)[0]
print("b* =", bs, "  h* =", sp.sqrt(h2.subs(b, bs)),
      "  h*/b* =", sp.simplify(sp.sqrt(h2.subs(b, bs)) / bs))
print("W'' (b*) =", sp.simplify(sp.diff(Wf, b, 2).subs(b, bs)))
print("W_max =", sp.simplify(Wf.subs(b, bs)))
\end{python}
\begin{risultato}
b* = sqrt(3)*d/3   h* = sqrt(6)*d/3   h*/b* = sqrt(2)
W'' (b*) = -sqrt(3)*d/3
W_max = sqrt(3)*d**3/27
\end{risultato}

\begin{esercizio}[Il principio di Fermat e la legge di Snell]\label{es:snell}
Un raggio di luce va dal punto $A=(0,h_{1})$, nel mezzo $y>0$ dove la velocità della
luce è $v_{1}$, al punto $B=(L,-h_{2})$, nel mezzo $y<0$ dove la velocità è $v_{2}$,
attraversando l'asse $x$ nel punto $(x,0)$. Dimostrare che il tempo di percorrenza è
minimo quando
\begin{equation}\label{eq:snell}
  \frac{\sen\theta_{1}}{v_{1}}=\frac{\sen\theta_{2}}{v_{2}},
\end{equation}
dove $\theta_{1}$, $\theta_{2}$ sono gli angoli che i due tratti del percorso formano con
la normale alla superficie di separazione.
\begin{soluzione}
Il tempo di percorrenza è
\begin{align*}
  T(x)&=\frac{\sqrt{h_{1}^{2}+x^{2}}}{v_{1}}+\frac{\sqrt{h_{2}^{2}+(L-x)^{2}}}{v_{2}},\\
  T'(x)&=\frac{x}{v_{1}\sqrt{h_{1}^{2}+x^{2}}}-\frac{L-x}{v_{2}\sqrt{h_{2}^{2}+(L-x)^{2}}} .
\end{align*}
Le frazioni $\frac{x}{\sqrt{h_{1}^{2}+x^{2}}}$ e $\frac{L-x}{\sqrt{h_{2}^{2}+(L-x)^{2}}}$ sono
proprio $\sen\theta_{1}$ e $\sen\theta_{2}$; quindi $T'(x)=0$ equivale alla
\eqref{eq:snell}. Resta da provare che il punto stazionario esiste, è unico ed è il
minimo. Si ha $T'(0)<0<T'(L)$, e un calcolo diretto dà
\[
  T''(x)=\frac{h_{1}^{2}}{v_{1}(h_{1}^{2}+x^{2})^{3/2}}+\frac{h_{2}^{2}}{v_{2}(h_{2}^{2}+(L-x)^{2})^{3/2}}>0 ,
\]
quindi $T'$ è strettamente crescente: per il teorema degli zeri si annulla in un unico
punto $x^{*}\in(0,L)$, nel quale $T'$ passa da negativa a positiva. Per il criterio di
monotonia, $T$ è strettamente decrescente in $[0,x^{*}]$ e strettamente crescente in
$[x^{*},L]$, e quindi $x^{*}$ è il punto di minimo assoluto di $T$ in $[0,L]$ (e anche su $\R$,
perché $T$ è strettamente convessa). Il programma seguente risolve numericamente
l'equazione $T'(x)=0$ per $h_{1}=h_{2}=1$, $L=2$, $v_{1}=1$, $v_{2}=0{,}75$ e verifica la
\eqref{eq:snell}. \emph{Scelta logica:} la convessità risolve in un colpo solo
esistenza, unicità e natura del punto stazionario.
\end{soluzione}
\end{esercizio}

\begin{python}[Il percorso di tempo minimo]
import math
from scipy.optimize import brentq

# raggio da A=(0, h1) a B=(L, -h2); velocità v1 sopra, v2 sotto
h1, h2, L, v1, v2 = 1.0, 1.0, 2.0, 1.0, 0.75
T = lambda x: math.hypot(h1, x) / v1 + math.hypot(h2, L - x) / v2
dT = lambda x: (x / math.hypot(h1, x) / v1
                - (L - x) / math.hypot(h2, L - x) / v2)
xs = brentq(dT, 0, L)                    # T'(x) = 0, unico zero
s1, s2 = xs / math.hypot(h1, xs), (L - xs) / math.hypot(h2, L - xs)
print(f"x* = {xs:.6f}   T(x*) = {T(xs):.6f}")
print(f"sen t1 / v1 = {s1 / v1:.6f}   sen t2 / v2 = {s2 / v2:.6f}")
print(f"T(0) = {T(0):.6f}   T(L) = {T(L):.6f}")
\end{python}
\begin{risultato}
x* = 1.270505   T(x*) = 3.267251
sen t1 / v1 = 0.785794   sen t2 / v2 = 0.785794
T(0) = 3.981424   T(L) = 3.569401
\end{risultato}

\needspace{8\baselineskip}
\subsection*{Limiti con il teorema di L'Hôpital e con la formula di Taylor}

\begin{esercizio}
Calcolare $\displaystyle\lim_{x\to0}\frac{e^{x}-e^{-x}-2x}{x-\sen x}$ in due modi: con il
teorema di L'Hôpital e con la formula di Taylor.
\begin{soluzione}
È una forma $\frac00$. Derivando numeratore e denominatore si ottengono
successivamente i rapporti
\[
  \frac{e^{x}+e^{-x}-2}{1-\cos x}, \qquad \frac{e^{x}-e^{-x}}{\sen x}, \qquad
  \frac{e^{x}+e^{-x}}{\cos x};
\]
i primi due sono ancora forme $\frac00$, l'ultimo tende a $2$. Applicando tre volte il
teorema di L'Hôpital, e risalendo la catena, il limite cercato vale $2$. Con la formula
di Taylor \eqref{eq:taylor-exp} di ordine $3$, applicata a $e^{x}$ e a $e^{-x}$,
\[
\begin{aligned}
  e^{x}-e^{-x}-2x&=\Bigl(1+x+\frac{x^{2}}2+\frac{x^{3}}6\Bigr)
  -\Bigl(1-x+\frac{x^{2}}2-\frac{x^{3}}6\Bigr)-2x+o(x^{3})\\
  &=\frac{x^{3}}3+o(x^{3})
\end{aligned}
\]
(i termini di grado pari si cancellano), e $x-\sen x=\frac{x^{3}}6+o(x^{4})$ per la
\eqref{eq:taylor-sen}; quindi il rapporto tende a $\frac{1/3}{1/6}=2$.
\emph{Scelta logica:} la formula di Taylor rivela subito che numeratore e denominatore
sono infinitesimi del terzo ordine, il che spiega perché il teorema di L'Hôpital
richiede tre applicazioni.
\end{soluzione}
\end{esercizio}

\needspace{8\baselineskip}
\begin{esercizio}
Calcolare $\displaystyle\lim_{x\to+\infty}x\Bigl(\frac\pi2-\operatorname{arctg}x\Bigr)$.
\begin{soluzione}
È una forma $\infty\cdot0$. Scrivendola come quoziente,
\[
  \lim_{x\to+\infty}\frac{\frac\pi2-\operatorname{arctg}x}{1/x}
  =\lim_{x\to+\infty}\frac{-\frac1{1+x^{2}}}{-\frac1{x^{2}}}
  =\lim_{x\to+\infty}\frac{x^{2}}{1+x^{2}}=1 .
\]
Il risultato dice che $\operatorname{arctg}x=\frac\pi2-\frac1x+o\bigl(\frac1x\bigr)$ per
$x\to+\infty$: il grafico si avvicina all'asintoto $y=\frac\pi2$ come l'iperbole
$y=\frac\pi2-\frac1x$. \emph{Scelta logica:} conviene mettere a denominatore il fattore
$\frac1x$, la cui derivata è semplice, e lasciare a numeratore la funzione trascendente,
la cui derivata è razionale.
\end{soluzione}
\end{esercizio}

\begin{esercizio}
Stabilire la natura del punto stazionario $x=0$ della funzione
$f(x)=e^{x}+e^{-x}-x^{2}-2$.
\begin{soluzione}
Si ha $f'(x)=e^{x}-e^{-x}-2x$, $f''(x)=e^{x}+e^{-x}-2$, $f'''(x)=e^{x}-e^{-x}$,
$f^{(4)}(x)=e^{x}+e^{-x}$; quindi $f'(0)=f''(0)=f'''(0)=0$ e $f^{(4)}(0)=2>0$. Per il
\cref{teo:criterio-ordine-n} con $n=4$ (pari), l'origine è di minimo relativo stretto.
In accordo, dalla \eqref{eq:taylor-exp} si ricava
$e^{x}+e^{-x}=2+x^{2}+\frac{x^{4}}{12}+o(x^{4})$, cioè $f(x)=\frac{x^{4}}{12}+o(x^{4})$.
\emph{Scelta logica:} il criterio della derivata seconda non decide, perché
$f''(0)=0$; si prosegue fino alla prima derivata non nulla, la cui parità decide.
\end{soluzione}
\end{esercizio}

\begin{esercizio}
Calcolare un valore approssimato di $\sqrt{1{,}1}$ con il polinomio di Taylor di ordine
$2$ di $\sqrt{1+x}$, e stimare l'errore.
\begin{soluzione}
Dalla \eqref{eq:taylor-potenza} con $b=\frac12$:
$\sqrt{1+x}=1+\frac x2-\frac{x^{2}}8+o(x^{2})$. Per $x=0{,}1$ si ottiene
$1+0{,}05-0{,}00125=1{,}04875$, mentre $\sqrt{1{,}1}=1{,}048809\dots$: l'errore è circa
$5{,}9\cdot10^{-5}$. Il termine successivo dello sviluppo è
$\frac{\frac12(-\frac12)(-\frac32)}{3!}x^{3}=\frac{x^{3}}{16}=6{,}25\cdot10^{-5}$, che dà
l'ordine di grandezza corretto dell'errore. \emph{Scelta logica:} la formula di Taylor
con il resto di Peano dice che l'errore è piccolo rispetto a $x^{2}$, ma non quanto; una
stima effettiva si ottiene guardando il primo termine trascurato (e, in modo rigoroso,
con la forma di Lagrange del resto, \cref{teo:resto-lagrange}).
\end{soluzione}
\end{esercizio}

\needspace{8\baselineskip}
\subsection*{Studi di funzione}

\begin{esercizio}\label{es:studio-razionale}
Studiare la funzione $f(x)=\dfrac{x^{2}}{x-1}$ e disegnarne il grafico.
\begin{soluzione}
\textbf{A.} Dominio: $x\neq1$. \textbf{B.} Nessuna simmetria; $f(x)=0$ solo per $x=0$;
$f>0$ per $x>1$, $f\le0$ per $x<1$. \textbf{C.} $f(x)\to\pm\infty$ per $x\to\pm\infty$ e
per $x\to1^{\pm}$: la retta $x=1$ è asintoto verticale. \textbf{D.} Si ha
\[
  f'(x)=\frac{2x(x-1)-x^{2}}{(x-1)^{2}}=\frac{x(x-2)}{(x-1)^{2}} ,
\]
positiva per $x<0$ e per $x>2$, negativa in $(0,1)$ e in $(1,2)$: $x=0$ è di massimo
relativo, con $f(0)=0$, e $x=2$ è di minimo relativo, con $f(2)=4$. \textbf{E.}
Scrivendo $f(x)=x+1+\frac1{x-1}$ (divisione tra polinomi) si ricava facilmente
$f''(x)=\frac{2}{(x-1)^{3}}$: $f$ è concava per $x<1$ e convessa per $x>1$; non ci sono
flessi. \textbf{F.} Dalla stessa scrittura, $f(x)-(x+1)=\frac1{x-1}\to0$ per
$x\to\pm\infty$: la retta $y=x+1$ è asintoto obliquo da entrambe le parti (con la
\cref{prop:asintoto-obliquo}: $m=\lim f(x)/x=\lim x/(x-1)=1$ e
$q=\lim\bigl(f(x)-x\bigr)=\lim x/(x-1)=1$). Il grafico è nella
\cref{fig:studio-razionale}. Si noti che il massimo relativo $f(0)=0$ è minore del minimo
relativo $f(2)=4$: «massimo» e «minimo» relativi sono nozioni locali.
\emph{Scelta logica:} la divisione tra polinomi semplifica tutti i passi successivi:
rende evidenti l'asintoto obliquo, la posizione del grafico rispetto a esso (sopra per
$x>1$, sotto per $x<1$) e il calcolo di $f''$.
\end{soluzione}
\end{esercizio}

\begin{figure}[htb]
\centering
\begin{tikzpicture}
\begin{axis}[width=0.8\textwidth, height=7cm, xmin=-3.7, xmax=5.7, ymin=-5, ymax=8,
  xtick={2}, ytick={4}, xlabel={$x$}, ylabel={$y$}, clip=false]
  \addplot[tratteggio, domain=-3.5:5.5, samples=2] {x+1};
  \addplot[tratteggio] coordinates {(1,-5) (1,8)};
  \addplot[curva, domain=-3.5:0.84, samples=150] {x^2/(x-1)};
  \addplot[curva, domain=1.2:5.5, samples=150] {x^2/(x-1)};
  \addplot[tratteggio] coordinates {(2,0) (2,4) (0,4)};
  \addplot[only marks, mark=*, mark size=1.6pt, rame] coordinates {(0,0) (2,4)};
  \node[font=\scriptsize, grigio, anchor=west] at (axis cs:4.6,4.9) {$y=x+1$};
  \node[font=\scriptsize, blu, anchor=west] at (axis cs:1.25,7.4) {$y=\frac{x^{2}}{x-1}$};
\end{axis}
\end{tikzpicture}
\caption{Il grafico di $f(x)=\frac{x^{2}}{x-1}$ (\cref{es:studio-razionale}): asintoto
verticale $x=1$, asintoto obliquo $y=x+1$, massimo relativo in $(0,0)$ e minimo relativo
in $(2,4)$.}
\label{fig:studio-razionale}
\end{figure}

\begin{esercizio}\label{es:e-pi}
Studiare la funzione $f(x)=\dfrac{\log x}{x}$ e dedurne quale dei due numeri $e^{\pi}$ e
$\pi^{e}$ sia il maggiore.
\begin{soluzione}
Il dominio è $(0,+\infty)$. Si ha $f(x)\to-\infty$ per $x\to0^{+}$ (asintoto verticale
$x=0$) e $f(x)\to0$ per $x\to+\infty$ per la \eqref{eq:hopital-log-potenza} (asintoto
orizzontale $y=0$). La derivata $f'(x)=\frac{1-\log x}{x^{2}}$ è positiva per $x<e$ e
negativa per $x>e$: $f$ è strettamente crescente in $(0,e]$ e strettamente decrescente in
$[e,+\infty)$, con massimo assoluto $f(e)=\frac1e$. La derivata seconda
$f''(x)=\frac{2\log x-3}{x^{3}}$ cambia segno in $x=e^{3/2}$, che è un punto di flesso.
Poiché $\pi>e$ e $f$ è strettamente decrescente in $[e,+\infty)$, risulta
$\frac{\log\pi}{\pi}<\frac{\log e}{e}=\frac1e$, cioè $e\log\pi<\pi$, e quindi, per la
monotonia dell'esponenziale, $\pi^{e}=e^{e\log\pi}<e^{\pi}$. Numericamente,
$\pi^{e}\approx22{,}46$ ed $e^{\pi}\approx23{,}14$. \emph{Scelta logica:} il confronto tra
$a^{b}$ e $b^{a}$ si trasforma, prendendo i logaritmi e dividendo per $ab$, nel confronto
tra $\frac{\log a}a$ e $\frac{\log b}b$, cioè tra due valori della stessa funzione: lo
studio della monotonia risolve il problema per tutte le coppie di numeri distinti non
minori di $e$.
\end{soluzione}
\end{esercizio}

%% ----------------------------------------------------------------------------
%%  §6.10  Punti fondamentali del capitolo
%% ----------------------------------------------------------------------------
\section{Punti fondamentali del capitolo}
\label{sec:sintesi-cap06}

La derivata è un'informazione locale; il teorema di Lagrange la trasforma in
un'informazione globale. Dal segno di $f'$ e di $f''$ in un intervallo si deducono la
monotonia e la convessità di $f$, e quindi i massimi, i minimi e la forma del grafico.
Il teorema di L'Hôpital e la formula di Taylor completano il quadro con due strumenti
per il calcolo dei limiti e l'approssimazione delle funzioni.

\begin{puntichiave}
\item \textbf{Teorema di Fermat.} Se $x_{0}$ è un punto \emph{interno} di massimo o di
minimo relativo e $f$ è derivabile in $x_{0}$, allora $f'(x_{0})=0$ (\cref{teo:fermat}).
La condizione è necessaria ma non sufficiente ($x^{3}$ in $0$); negli estremi
dell'intervallo vale solo una disuguaglianza sul segno di $f'$
(\cref{prop:fermat-estremi}).
\item \textbf{Estremi assoluti in $[a,b]$.} Se $f$ è continua in $[a,b]$, il massimo e il
minimo assoluti, che esistono per il \cref{teo:weierstrass}, vanno cercati tra gli
estremi $a$, $b$, i punti interni in cui $f$ non è derivabile e i punti stazionari:
basta confrontare i valori di $f$ in questi punti.
\item \textbf{Rolle e Lagrange.} Se $f$ è continua in $[a,b]$ e derivabile in $(a,b)$,
esiste $x_{0}\in(a,b)$ con $f(b)-f(a)=f'(x_{0})(b-a)$ (\cref{teo:lagrange}); il teorema di
Rolle è il caso $f(a)=f(b)$, e ciascuna ipotesi è indispensabile
(\cref{oss:rolle-ipotesi}). Se $\abs{f'}\le M$ in un intervallo, allora
$\abs{f(x)-f(y)}\le M\abs{x-y}$ (\cref{cor:lipschitz}).
\item \textbf{Criterio di monotonia.} Se $f$ è continua in un intervallo e derivabile nei
punti interni, $f$ è crescente se e solo se $f'\ge0$ (\cref{teo:criterio-monotonia}), e
$f'>0$ implica la stretta monotonia (\cref{cor:derivata-positiva}), ma non viceversa
($x^{3}$). Errore tipico: dimenticare che il dominio deve essere un intervallo, come
mostra $1/x$.
\item \textbf{Funzioni con derivata nulla.} In un intervallo, $f'=0$ implica $f$ costante
(\cref{cor:funzioni-costanti}), e due funzioni con la stessa derivata differiscono per
una costante (\cref{cor:stessa-derivata}); ne segue, ad esempio, l'unicità della legge
esponenziale della pressione nella stratosfera (\cref{es:atmosfera-isoterma}).
\item \textbf{Test per gli estremi.} Se $f$ è continua in un intorno di $x_{0}$, derivabile
per $x\neq x_{0}$, con $f'\ge0$ a sinistra e $f'\le0$ a destra di $x_{0}$, il punto è di
massimo relativo, anche se $f$ non è derivabile in $x_{0}$
(\cref{cor:test-derivata-prima}). Se $f'(x_{0})=0$ e $f''(x_{0})>0$ (oppure
$<0$), il punto è di minimo (oppure di massimo) relativo stretto
(\cref{teo:test-derivata-seconda,oss:test-derivata-seconda}).
\item \textbf{Convessità e flessi.} Per $f$ derivabile in $[a,b]$ e derivabile due volte in
$(a,b)$, le condizioni «$f$ convessa», «$f'$ crescente» e «$f''\ge0$ in $(a,b)$» sono
equivalenti (\cref{teo:criterio-convessita}). In un flesso in cui $f$ è derivabile due
volte risulta $f''(x_{0})=0$ (\cref{cor:flesso-necessaria}), ma non viceversa ($x^{4}$ in
$0$).
\item \textbf{Teorema di L'Hôpital.} Se $f$ e $g$ sono derivabili vicino a $x_{0}$ (escluso
al più $x_{0}$), con $g'\neq0$, se $f,g\to0$ oppure $\abs g\to+\infty$, e se esiste
$\lim f'/g'=\ell\in\overline\R$, allora anche $f/g\to\ell$
(\cref{teo:hopital}); la dimostrazione generale poggia sul teorema di Cauchy
(\cref{teo:cauchy}).
\item \textbf{Errori tipici con L'Hôpital.} Il teorema si applica solo alle forme
indeterminate; se $f'/g'$ non ha limite, non si può concludere nulla su
$f/g$; non dimostra i limiti notevoli, da cui sono state ricavate le derivate delle
funzioni elementari.
\item \textbf{Asintoti e studio di funzione.} La retta $y=mx+q$ è asintoto obliquo per
$x\to+\infty$ se e solo se $m=\lim f(x)/x$ e $q=\lim\bigl(f(x)-mx\bigr)$ esistono
finiti (\cref{prop:asintoto-obliquo}); lo studio di funzione segue lo schema dominio,
simmetrie, limiti, $f'$, $f''$, asintoti obliqui (\cref{sec:studio-funzione}).
\item \textbf{Formula di Taylor con il resto di Peano.} Se $f$ è derivabile $n$ volte in
$x_{0}$, allora $f(x)=T_{n}(x)+o\bigl((x-x_{0})^{n}\bigr)$ per $x\to x_{0}$
(\cref{teo:taylor}), e $T_{n}$ è l'unico polinomio di grado $\le n$ con questa proprietà
(\cref{prop:unicita-taylor}); gli sviluppi di $e^{x}$, $\log(1+x)$, $\sen x$, $\cos x$,
$(1+x)^{b}$ vanno ricordati.
\item \textbf{Criterio delle derivate successive.} Sia $n\ge2$; se le derivate di $f$ di
ordine $1,\dots,n-1$ sono nulle in $x_{0}$ e $f^{(n)}(x_{0})\neq0$, il punto $x_{0}$ è di
estremo relativo stretto se $n$ è pari (minimo se $f^{(n)}(x_{0})>0$, massimo se
$f^{(n)}(x_{0})<0$), mentre non è di estremo se $n$ è dispari
(\cref{teo:criterio-ordine-n}).
\item \textbf{Applicazioni aeronautiche.} L'efficienza massima si trova con la tangente
alla polare condotta dall'origine (\cref{fig:polare-tangente}); alla velocità di minima
resistenza la resistenza parassita e quella indotta si eguagliano, e
$V_{\min D}/V_{\min P}=3^{1/4}$ \eqref{eq:V-min-P}; l'errore relativo di
$\sen\varphi\approx\varphi$ è circa $\varphi^{2}/6$.
\end{puntichiave}

\paragraph{Il filo del discorso.} Il capitolo lascia aperte due questioni. La formula di
Taylor con il resto di Peano descrive soltanto il comportamento vicino a $x_{0}$, senza
una stima quantitativa dell'errore: la daranno i resti integrale e di Lagrange del
\cref{cap:taylor}. Inoltre tutti i risultati riguardano funzioni di una sola variabile.
Il prossimo capitolo li estende alle funzioni di più variabili: la derivata diventa la
derivata parziale, la condizione di Fermat $f'(x_{0})=0$ diventa l'annullarsi del
gradiente, e il criterio della derivata seconda diventa quello dell'hessiano
(\cref{sec:estremi-due}).

